\documentclass[11pt]{article}

\usepackage[utf8]{inputenc}
% microtype 의 폰트 확장은 scalable 폰트를 요구한다. 기본 CM 이 비트맵(Type3)으로
% 잡히는 배포판에서는 lmodern 없이는 "auto expansion is only possible with
% scalable fonts" 로 컴파일이 죽는다.
\usepackage{lmodern}
\usepackage[T1]{fontenc}
\usepackage[margin=1in]{geometry}
\usepackage{amsmath,amssymb}
\usepackage{graphicx}
\usepackage{longtable,booktabs}
\usepackage{microtype}
% svgnames 를 줘야 teal 이 정의된다. 없으면 \cite 하나마다
% "Undefined color `teal'" 이 뜬다 (본편 산출물에서 732회).
\usepackage[svgnames]{xcolor}
% hyperref goes last; it makes \cite clickable through to the bibliography.
\usepackage[colorlinks=true,linkcolor=blue,citecolor=teal,urlcolor=magenta]{hyperref}

% pandoc emits \tightlist but expects the preamble to define it.
\providecommand{\tightlist}{%
  \setlength{\itemsep}{0pt}\setlength{\parskip}{0pt}}

% If the compile times out on Overleaf, comment out \tableofcontents below first.
\title{Evaluation of Large Language Models: A Comprehensive Survey}
\author{Generated by AutoSurvey}
\date{}

\begin{document}
\maketitle
\tableofcontents
\newpage

\section{Introduction to Large Language Models (LLMs)}

\subsection{Emergence and Importance of Large Language Models (LLMs)}

The emergence and rapid advancement of large language models (LLMs) have been nothing short of revolutionary, transforming the landscape of natural language processing and ushering in a new era of artificial intelligence. These sophisticated models, trained on vast troves of textual data, have demonstrated a remarkable ability to understand, generate, and reason about language in ways that were once unimaginable \cite{ref1,ref2}.

The journey of LLMs began with the groundbreaking work on transformer architectures and self-attention mechanisms, which laid the foundation for models like GPT and BERT \cite{ref3}. These early models showcased the potential of large-scale language modeling, paving the way for increasingly sophisticated iterations, such as GPT-3 and its successors \cite{ref2}. The sheer scale of these models, with parameter sizes ranging from billions to trillions, has unlocked unprecedented capabilities, enabling them to excel at a wide array of natural language tasks, from text generation and translation to question answering and content creation \cite{ref1}.

The significance of LLMs lies not only in their impressive performance but also in their potential to revolutionize various industries and domains. These models have found applications in fields as diverse as customer service, healthcare, education, and scientific research \cite{ref4,ref5,ref6}. Their ability to understand and generate human-like text has led to the development of intelligent chatbots, virtual assistants, and content-generation tools that can enhance productivity, automate repetitive tasks, and aid in decision-making \cite{ref7,ref8}.

Moreover, the versatility of LLMs has extended beyond traditional natural language processing, with their capabilities being leveraged in multimodal tasks, such as image captioning and visual question answering \cite{ref9}. This integration of language understanding with other modalities has the potential to unlock new frontiers in artificial intelligence, paving the way for more holistic and contextual interactions between humans and machines.

The rapid progress of LLMs has also raised important questions and concerns regarding their safety, reliability, and alignment with human values \cite{ref10,ref11}. As these models become increasingly influential in various domains, there is a growing need to ensure that they are deployed responsibly and in a manner that minimizes potential risks, such as the generation of biased or harmful content, privacy breaches, and the undermining of human agency \cite{ref7,ref12}.

To address these challenges and unlock the full potential of LLMs, the research community has devoted significant efforts to developing comprehensive evaluation frameworks and benchmarks \cite{ref11,ref13}. These initiatives aim to assess the models' capabilities across a wide range of tasks, domains, and societal impacts, providing valuable insights into their strengths, weaknesses, and areas for improvement \cite{ref14,ref15}.

As the field of LLMs continues to evolve, the importance of rigorous evaluation and responsible development cannot be overstated. By understanding the capabilities and limitations of these models, researchers, policymakers, and practitioners can work towards harnessing their transformative potential while addressing the challenges that come with their increasing prominence in our lives \cite{ref16,ref17}.

\subsection{Scaling and Capabilities of LLMs}

The rapid progress in large language models (LLMs) over the past few years can be largely attributed to the scaling of model size, training data, and computational resources. As highlighted in \cite{ref18}, the performance of LLMs has been shown to improve remarkably with the increase in the number of model parameters, following the predictions of scaling laws \cite{ref19,ref20}.

Large-scale language models have grown from the early successes of GPT-3 \cite{ref21}, which had 175 billion parameters, to even more powerful models such as GPT-4 with 1.5 trillion parameters \cite{ref22}. This exponential growth in model size has been enabled by advancements in hardware, software, and training techniques, allowing researchers to push the boundaries of what is computationally feasible.

Alongside the scaling of model size, the availability of vast and diverse training data has been crucial for the remarkable capabilities exhibited by modern LLMs. As described in \cite{ref23}, the training data for LLMs has grown from the initial corpus of web pages, books, and articles to a comprehensive collection of datasets spanning a wide range of domains, including natural language, programming languages, and specialized knowledge. This expansive training data has enabled LLMs to acquire a deep understanding of language, knowledge, and reasoning, allowing them to excel in various natural language processing tasks.

The scaling of computational resources has been another key factor in the development of powerful LLMs. The training of these models requires immense computational power, with the training of GPT-3 alone estimated to have consumed over 1.2 million GPU-hours \cite{ref21}. The advent of specialized hardware, such as GPUs and TPUs, along with advancements in distributed training techniques, have made it possible to train these large-scale models efficiently \cite{ref24}.

The synergistic effect of scaling model size, training data, and computational resources has resulted in LLMs with remarkable capabilities in language understanding and generation. As described in \cite{ref1}, these models have demonstrated strong performance across a wide range of natural language processing tasks, including text generation, question answering, sentiment analysis, and task-specific applications.

One of the most notable capabilities of LLMs is their ability to perform in-context learning, where the models can adapt to new tasks or domains with minimal fine-tuning or task-specific training \cite{ref21}. This capability is particularly valuable in scenarios where labeled data is scarce, as LLMs can leverage their broad knowledge and language understanding to generalize to new tasks.

Furthermore, LLMs have shown impressive performance in tasks that require reasoning, abstraction, and causal understanding \cite{ref25}. This suggests that the scaling of these models has not only improved their language modeling capabilities but has also enhanced their ability to capture and reason about complex relationships and patterns in language.

However, the scaling of LLMs also introduces new challenges and limitations. As discussed in \cite{ref26}, reducing the size of LLMs can lead to a significant deterioration in their ability to recall factual information from their pre-training data, while their in-context learning capabilities may be more resilient to model scaling. This highlights the need for a nuanced understanding of the trade-offs involved in scaling LLMs and the potential for developing specialized models that excel in specific domains or tasks.

Additionally, the scaling of LLMs has raised concerns about their transparency, reliability, and potential for misuse \cite{ref27}. As these models become more powerful and influential, it is crucial to develop robust methods for evaluating their capabilities, understanding their limitations, and ensuring their alignment with human values and preferences.

\subsection{Applications and Societal Impact of LLMs}

The widespread adoption and application of large language models (LLMs) across various industries and domains have been truly transformative. LLMs have demonstrated unprecedented capabilities in natural language processing tasks, enabling significant advancements in fields ranging from content generation and information retrieval to medical diagnosis and legal assistance \cite{ref28,ref29}.

One of the most prominent applications of LLMs has been in the field of content generation. These models have proven adept at producing human-like text, allowing for the automated creation of articles, stories, and even creative works \cite{ref22}. This has led to the rise of AI-generated content, which has both benefits and drawbacks. On the one hand, LLMs can help alleviate content creation bottlenecks, enabling faster and more efficient production of written materials. This can be particularly useful in industries such as journalism, marketing, and education, where content generation is a crucial task. Additionally, LLMs can assist in the personalization of content, tailoring it to the specific needs and preferences of individual users \cite{ref30}.

However, the use of LLMs in content generation also raises concerns about the authenticity and trustworthiness of the produced content. There is a risk of LLMs generating ``fake news'' or misinformation, which could have serious societal implications \cite{ref31}. Moreover, the potential for LLMs to be used in the creation of deepfakes, where synthetic media is generated to misrepresent individuals or events, poses a significant threat to the integrity of digital communication \cite{ref32}.

The widespread application of LLMs extends beyond content generation, with the models being increasingly integrated into decision-support systems across various industries. In the healthcare domain, for example, LLMs have shown promise in tasks such as clinical decision support, medical image analysis, and drug discovery \cite{ref33,ref34}. While these applications can potentially improve healthcare outcomes and enhance patient care, the high-stakes nature of the medical field necessitates rigorous evaluation and safeguards to ensure the reliability and safety of LLM-powered systems \cite{ref35}.

In the legal domain, LLMs have been utilized for tasks such as legal research, document review, and contract analysis \cite{ref36,ref37}. These applications can potentially improve the efficiency and accessibility of legal services, particularly for individuals with limited resources. However, the use of LLMs in legal settings also raises concerns about the potential for biased or erroneous outputs, which could have significant consequences for individuals and society \cite{ref38}.

Beyond these specific applications, LLMs have also demonstrated the potential to transform various other sectors, including finance \cite{ref39}, education \cite{ref40}, and transportation \cite{ref41}. These models have the capacity to automate a wide range of tasks, enhance productivity, and provide personalized assistance to users.

However, the societal implications of the widespread adoption of LLMs are complex and multifaceted. There are concerns about the potential displacement of human labor, the amplification of biases and inequalities, and the erosion of privacy and data rights \cite{ref42}. Additionally, the rapid advancement of LLM technology has outpaced the development of appropriate regulatory frameworks, raising questions about the ethical and responsible deployment of these models \cite{ref7}.

To address these challenges, there is a growing need for interdisciplinary collaboration between researchers, policymakers, and industry stakeholders to establish guidelines, standards, and oversight mechanisms for the development and deployment of LLMs. This includes addressing issues related to data privacy, algorithmic bias, and the potential misuse of these models \cite{ref43}.

Overall, the applications and societal impact of LLMs are far-reaching and multifaceted. While these models offer immense potential for transformation across various industries and domains, their widespread adoption also raises critical concerns that must be carefully navigated to ensure the responsible and beneficial development of these technologies.

\subsection{Challenges and Motivations for Evaluating LLMs}

The rapid progress and widespread adoption of large language models (LLMs) have brought about a host of challenges and concerns that necessitate rigorous evaluation. As these models become increasingly capable and ubiquitous, issues related to safety, reliability, fairness, and transparency have emerged as pressing priorities \cite{ref32,ref35}.

One of the key challenges in evaluating LLMs is their potential to generate unsafe or harmful content. LLMs can be susceptible to biases and errors, leading to the production of misinformation, hate speech, or even content that could be used for malicious purposes \cite{ref44,ref45}. This raises concerns about the responsible deployment of these models, particularly in high-stakes domains such as healthcare, finance, and education \cite{ref46,ref47}.

Moreover, the reliability of LLM outputs is another major challenge. These models can exhibit inconsistencies, hallucinations, and unpredictable behaviors, which undermine their trustworthiness \cite{ref48,ref34}. This is especially problematic in applications where accuracy and reliability are of paramount importance, such as medical diagnosis, legal advice, or scientific research \cite{ref49,ref50}.

Fairness is another crucial concern with LLMs, as these models can perpetuate and amplify societal biases related to race, gender, and other demographic factors \cite{ref51,ref52}. Biased outputs from LLMs can have far-reaching consequences, exacerbating inequalities and undermining the principles of equity and inclusivity \cite{ref53,ref54}. The need to ensure that LLMs are designed and deployed in a fair and equitable manner is, therefore, a pressing challenge.

Lastly, the issue of transparency and interpretability in LLMs is a significant concern. The complex and opaque nature of these models makes it challenging to understand their inner workings, decision-making processes, and the reasons behind their outputs \cite{ref55,ref44}. This lack of transparency can hinder our ability to hold these models accountable, as well as to diagnose and mitigate potential issues \cite{ref56,ref57}.

Given the multifaceted challenges posed by LLMs, the need for rigorous evaluation has become increasingly apparent. Comprehensive and systematic assessment of these models is essential to ensure their safety, reliability, fairness, and transparency, ultimately enabling their responsible and beneficial deployment \cite{ref57,ref58}. Such evaluation efforts should address the ``what, where, and how'' of LLM assessment, considering their capabilities, societal impacts, and underlying mechanisms \cite{ref59,ref60}.

By addressing these challenges through robust evaluation, the AI community can work towards unlocking the full potential of LLMs while mitigating the risks and ensuring their alignment with human values and societal well-being. This endeavor is crucial for guiding the responsible development and deployment of these transformative technologies, ultimately shaping a future where the benefits of LLMs are realized in a safe, equitable, and transparent manner.

\section{Benchmark Datasets and Evaluation Frameworks}

\subsection{Evaluation Dimensions}

The rapid advancements of large language models (LLMs) have necessitated the development of comprehensive evaluation frameworks to assess their capabilities across diverse tasks and domains. These evaluation frameworks can be broadly categorized into three key dimensions: task-level, societal-level, and multi-faceted assessments \cite{ref61,ref53}.

Task-level Evaluation
At the task-level, the evaluation of LLMs has traditionally focused on their performance on specific natural language processing (NLP) benchmarks, such as text generation, question answering, and language understanding \cite{ref61,ref62}. These benchmarks provide a standardized way to measure the models' capabilities in handling well-defined tasks, allowing for direct comparisons between different LLMs and tracking progress over time.

However, as LLMs continue to demonstrate impressive results on these task-specific benchmarks, the research community has recognized the need to explore more complex and open-ended tasks that better reflect the real-world challenges faced by these models \cite{ref63,ref64}. These tasks may involve reasoning, commonsense understanding, and the ability to handle ambiguity and contextual information, which are crucial for the successful deployment of LLMs in diverse applications.

Societal-level Evaluation
Beyond task-level assessments, the evaluation of LLMs has also shifted towards understanding their societal impact and alignment with human values \cite{ref65,ref66}. This societal-level evaluation focuses on critical aspects such as fairness, safety, transparency, and ethical considerations.

Researchers have highlighted the need to assess the potential biases, discriminatory behaviors, and harmful outputs that LLMs may exhibit, particularly when deployed in high-stakes domains like healthcare, education, and public policy \cite{ref67,ref68}. Evaluating the alignment of LLMs with human values, such as being helpful, honest, and harmless, has become a crucial aspect of this societal-level assessment \cite{ref69}.

Multi-faceted Evaluation
To address the limitations of task-specific and societal-level evaluations, researchers have proposed more comprehensive, multi-faceted evaluation frameworks that consider a broader range of capabilities and potential impacts \cite{ref11,ref70}. These frameworks aim to provide a holistic assessment of LLMs, encompassing their knowledge, reasoning abilities, safety, reliability, transparency, and alignment with human preferences.

One key aspect of multi-faceted evaluation is the integration of various assessment techniques, including human evaluation, automated metrics, and the use of LLMs as evaluators themselves \cite{ref71,ref72}. By leveraging the complementary strengths of these different evaluation methods, researchers can obtain a more reliable and nuanced understanding of the LLMs' capabilities and limitations \cite{ref56,ref73}.

Furthermore, multi-faceted evaluation frameworks often aim to assess LLMs across a wide range of domains, from general language understanding to specialized tasks in fields like medicine, education, and law \cite{ref74,ref75}. This diversity of evaluation tasks helps to uncover the models' strengths, weaknesses, and potential biases in different contexts, providing valuable insights for their responsible development and deployment.

\subsection{Comprehensive Evaluation Frameworks}

The rapid development of large language models (LLMs) has necessitated the creation of comprehensive evaluation frameworks that can holistically assess their performance across diverse tasks and domains. These frameworks aim to address the ``what, where, and how'' of LLM evaluation, providing a structured approach to understanding the capabilities and limitations of these advanced language models.

One prominent example of a comprehensive evaluation framework is FreeEval \cite{ref76}. FreeEval is designed to enable trustworthy and efficient automatic evaluations of LLMs by addressing several key challenges. Firstly, FreeEval's unified abstractions simplify the integration and improve the transparency of diverse evaluation methodologies, including dynamic evaluation that demands sophisticated LLM interactions. This modular approach allows researchers and practitioners to seamlessly incorporate state-of-the-art evaluation techniques into their workflows, fostering a more transparent and reproducible evaluation ecosystem.

Secondly, FreeEval integrates meta-evaluation techniques, such as human evaluation and data contamination detection, which enhance the reliability and fairness of the evaluation outcomes. By introducing these additional layers of assessment, FreeEval ensures that the evaluation results are not biased or skewed by potential data issues or flaws in the evaluation process itself.

Moreover, FreeEval is designed with a high-performance infrastructure, including distributed computation and caching strategies, enabling extensive evaluations across multi-node, multi-GPU clusters for both open-source and proprietary LLMs. This efficient implementation allows researchers to rapidly assess the performance of various LLMs, accelerating the development and deployment of these powerful models.

Another comprehensive evaluation framework, CheckEval \cite{ref77}, addresses the challenges of ambiguity and inconsistency in current evaluation methods by introducing a checklist-based approach. CheckEval divides evaluation criteria into detailed sub-aspects and constructs a checklist of Boolean questions for each, simplifying the evaluation process and enhancing its robustness and reliability. This framework not only improves the interpretability of the evaluation but also demonstrates strong correlations with human judgments and highly consistent inter-annotator agreement, setting a new standard for LLM-based evaluation.

The HD-Eval \cite{ref78} framework takes a different approach, focusing on aligning LLM-based evaluators with human preferences through Hierarchical Criteria Decomposition. HD-Eval iteratively decomposes evaluation tasks into finer-grained criteria, aggregating them according to estimated human preferences, pruning insignificant criteria, and further decomposing significant criteria. By integrating these steps within an iterative alignment training process, HD-Eval obtains a hierarchical decomposition of criteria that comprehensively captures aspects of natural language at multiple levels of granularity. This approach has been shown to outperform state-of-the-art evaluators and provide deeper insights into the explanation of evaluation results and the task itself.

In addition to these comprehensive frameworks, the research community has also introduced several other evaluation platforms and suites that address the diverse aspects of LLM assessment. UltraEval \cite{ref79} is a user-friendly evaluation framework characterized by its lightweight, comprehensiveness, modularity, and efficiency. UltraEval enables the free combination of different models, tasks, prompts, and metrics within a unified evaluation workflow, supporting diverse models through a unified HTTP service and providing sufficient inference acceleration.

The PRE \cite{ref80} framework proposes a novel peer-review-based approach to automatically evaluate LLMs. Inspired by the peer review systems widely used in academic publication processes, PRE constructs a small qualification exam to select ``reviewers'' from a set of powerful LLMs, who then rate or compare the ``submissions'' written by different candidate LLMs. The final ranking of the evaluated LLMs is generated based on the results provided by all reviewers, addressing the issues of bias and low generalizability inherent in existing evaluation methods.

Furthermore, the research community has also introduced specialized evaluation frameworks targeting specific domains or aspects of LLM performance. For instance, the EquityMedQA dataset and accompanying assessment framework \cite{ref81} focuses on evaluating the potential biases and equity-related harms in LLM-generated content within the medical domain. Similarly, the ExpertMedQA benchmark \cite{ref82} rigorously evaluates LLM performance on open-ended, expert-level clinical questions, demanding an in-depth understanding and critical appraisal of up-to-date clinical literature.

These comprehensive evaluation frameworks, along with the continued development of specialized benchmarks and platforms, play a crucial role in advancing the responsible and transparent development of LLMs. By providing structured approaches to assess the capabilities, alignment, and safety of these models, the research community can better guide the evolution of LLMs towards maximizing societal benefit while minimizing potential risks.

\subsection{Benchmark Datasets and Evaluation Platforms}

The rapid advancements in large language models (LLMs) have highlighted the pressing need for comprehensive and standardized evaluation frameworks to assess their capabilities across diverse domains \cite{ref61}. To address this need, there has been a growing emphasis on the development of open-source materials and ongoing initiatives that aim to facilitate and advance the evaluation of LLMs.

One prominent example is the GPT-Fathom evaluation suite, which was introduced to provide a systematic and reproducible assessment of leading LLMs, including OpenAI's legacy models \cite{ref83}. The GPT-Fathom project is built upon the OpenAI Evals framework and offers a curated set of 20+ benchmarks across seven capability categories, enabling a comprehensive evaluation of LLMs under aligned settings. By conducting a retrospective study on OpenAI's earlier models, the GPT-Fathom project provides valuable insights into the evolutionary path from GPT-3 to GPT-4, shedding light on aspects such as the impact of adding code data on reasoning capabilities, the benefits of supervised fine-tuning and reinforcement learning from human feedback, and the trade-offs associated with model alignment.

Another significant open-source initiative is the PiCO (Peer Review in LLMs based on the Consistency Optimization) framework, which explores the concept of using peer-review mechanisms to evaluate LLMs in an unsupervised manner \cite{ref84}. In this setting, both open-source and closed-source LLMs are positioned in the same environment, capable of answering unlabeled questions and evaluating each other. The PiCO framework assigns each LLM a learnable capability parameter and aims to maximize the consistency of the LLMs' capabilities and scores, under the assumption that higher-level LLMs can evaluate others more accurately than lower-level ones.

The medical domain has also seen the development of several open-source initiatives to facilitate the evaluation and advancement of LLMs in healthcare applications. The MedAlpaca dataset and the Hippocrates open-source framework are two prominent examples \cite{ref85,ref86}. MedAlpaca provides a collection of over 160,000 entries specifically curated to fine-tune LLMs for effective medical applications, while Hippocrates offers an open-source ecosystem, including training datasets, codebases, and evaluation protocols, to enable collaborative research and the development of medical LLMs.

Recognizing the importance of evaluating LLMs' clinical capabilities, researchers have also introduced frameworks like the one proposed in ``Towards Automatic Evaluation for LLMs' Clinical Capabilities'' \cite{ref87}. This framework aims to automatically assess the performance of LLMs in delivering clinical services, such as disease diagnosis and treatment, by leveraging a multi-agent simulation environment and a retrieval-augmented evaluation approach.

The open-source community has also made significant contributions to the evaluation of LLMs in specialized domains, such as the legal and educational fields. For instance, the GPT Models in Construction Industry study explores the opportunities and limitations of using GPT models in the construction industry \cite{ref88}, while the Apprentices to Research Assistants work examines the potential of leveraging LLMs as research tools across various domains \cite{ref6}.

In addition to domain-specific initiatives, there are also ongoing efforts to develop comprehensive evaluation frameworks that cover a wide range of LLM capabilities, alignment, safety, and applicability. The ``Comprehensive Evaluation Frameworks for LLMs'' work aims to introduce open-source platforms that enable the assessment of LLMs' performance across diverse tasks and applications \cite{ref61}.

Furthermore, the research community has recognized the importance of evaluating LLMs' robustness, reliability, and alignment with human values and preferences. The EquityMedQA dataset and associated methodologies, for instance, focus on surfacing biases and equity-related harms in LLM-generated content \cite{ref81}. Similarly, the PiCO framework's peer-review mechanism and the L-Eval benchmark for long-context language models \cite{ref89} contribute to the evaluation of LLMs' reliability and alignment.

Overall, the open-source materials and ongoing initiatives discussed in this subsection highlight the research community's concerted efforts to develop comprehensive and standardized evaluation frameworks for LLMs. These efforts aim to facilitate the responsible development and deployment of LLMs, ensuring their capabilities are thoroughly assessed across diverse tasks, domains, and societal impact considerations. By fostering transparency, collaboration, and the adoption of rigorous evaluation practices, these initiatives pave the way for the continued advancement and responsible use of large language models.

\section{Evaluation Methodologies and Metrics}

\subsection{Human Evaluation}

Human evaluation has long been considered the gold standard for assessing the quality of open-ended natural language generation (NLG) tasks, where models are required to produce creative and diverse outputs \cite{ref90,ref91}. This is primarily because the nuanced and subjective nature of open-ended generation makes it challenging to capture all the relevant quality criteria using automated metrics alone \cite{ref92}.

The role of human evaluation in open-ended NLG tasks is twofold. First, it helps assess the overall quality of the generated outputs, capturing aspects such as fluency, coherence, relevance, and creativity that are difficult to quantify algorithmically \cite{ref93}. Second, it provides valuable insights into the human experience of interacting with the generated text, shedding light on aspects like the naturalness, trustworthiness, and alignment with human values and preferences \cite{ref94}.

One of the key strengths of human evaluation is its ability to capture the holistic and subjective nature of open-ended generation \cite{ref95}. By leveraging human judgment, researchers can gain a deeper understanding of how the generated outputs are perceived and interpreted by end-users, going beyond the limitations of purely algorithmic assessments \cite{ref96}.

However, human evaluation also faces several limitations and challenges. First, it is inherently subjective and prone to biases, as different human evaluators may have varying perceptions and preferences \cite{ref97}. This can lead to inconsistent and unreliable results, making it difficult to compare the performance of different NLG models across studies \cite{ref98}.

Second, human evaluation is often resource-intensive and time-consuming, requiring the recruitment and coordination of a large number of human assessors \cite{ref91}. This can limit the scalability and feasibility of human evaluation, especially in the face of rapidly evolving NLG models and the need for frequent evaluations \cite{ref99}.

To address these challenges, researchers have proposed various best practices and guidelines for conducting more reliable and reproducible human evaluations. These include standardizing the evaluation protocols, improving the recruitment and training of human assessors, and exploring innovative evaluation frameworks that leverage the complementary strengths of humans and machines \cite{ref100,ref66}.

For example, the GENIE framework \cite{ref100} introduces a standardized approach to human evaluation, including the use of well-defined annotation interfaces and automated mechanisms for maintaining annotator quality. The Collaborative Evaluation (CoEval) pipeline \cite{ref66} combines the strengths of human and LLM-based evaluators, leveraging LLMs to provide initial assessments and then utilizing human scrutiny to refine the evaluation.

Other approaches, such as the use of ``checklist'' evaluation \cite{ref101} and the introduction of multi-dimensional evaluation frameworks \cite{ref102}, aim to capture a more comprehensive and nuanced understanding of the generated outputs, going beyond overall quality assessments.

Furthermore, researchers have explored the potential of using large language models (LLMs) as alternative evaluators for open-ended generation tasks \cite{ref103,ref104}. LLMs have shown promising results in terms of their ability to assess the quality of generated text, with the potential to overcome some of the limitations of human evaluation, such as inconsistency and scalability \cite{ref105}.

However, the use of LLMs as evaluators also comes with its own set of challenges, such as the need to ensure their alignment with human preferences and the potential for biases and limitations in their assessments \cite{ref106,ref107}. Ongoing research is exploring ways to address these challenges and further enhance the reliability and effectiveness of LLM-based evaluators \cite{ref98}.

In conclusion, human evaluation remains a crucial component in the assessment of open-ended NLG tasks, providing invaluable insights into the subjective and holistic aspects of generated outputs. However, the challenges and limitations associated with human evaluation have prompted the exploration of innovative approaches, such as collaborative evaluation frameworks and the use of LLMs as alternative evaluators. As the field of NLG continues to evolve, the need for reliable and comprehensive evaluation methods will only grow, underscoring the importance of ongoing research in this area.

\subsection{Automated Metrics}

Automated evaluation metrics have long been essential tools for assessing the quality of generated text across various natural language processing tasks, including machine translation, text summarization, and dialogue response generation. These metrics can be broadly categorized into three main groups: word-based, grammar-based, and neural-based.

Word-based metrics, such as BLEU \cite{ref108} and METEOR \cite{ref109}, focus on the lexical overlap between the generated text and the reference text(s). These metrics calculate a score based on the number of n-grams (sequences of n words) that are shared between the candidate and reference texts. While these metrics are simple to compute and widely used, they have been criticized for their inability to capture semantic similarity and their sensitivity to lexical variations \cite{ref110}.

Grammar-based metrics, on the other hand, take into account the grammatical structure of the generated text. Examples include GLEU \cite{ref111} and GrammarScore \cite{ref112}, which evaluate the grammatical correctness of the generated text by comparing it to a set of grammatically correct reference sentences. These metrics can be particularly useful for tasks where grammatical correctness is a crucial aspect, such as grammatical error correction.

The emergence of large language models (LLMs) \cite{ref90} has led to the development of a new class of neural-based evaluation metrics that leverage the semantic understanding and contextual awareness of these models. Examples include BERTScore \cite{ref113}, BLEURT \cite{ref114}, and COMET \cite{ref115}, which use pre-trained language models to assess the semantic similarity between the generated text and the reference text(s). These neural-based metrics have been shown to achieve higher correlation with human judgments compared to traditional word-based and grammar-based metrics \cite{ref116}.

One of the key advantages of neural-based metrics is their ability to capture semantic similarity beyond exact word matches. This is particularly important in tasks where paraphrasing and lexical variation are common, such as dialogue response generation and text summarization. Additionally, these metrics can be fine-tuned on specific datasets or tasks, allowing for more tailored and accurate evaluations \cite{ref90}.

However, the performance of neural-based metrics is not without its challenges. One major concern is their sensitivity to the specific training data and model architecture used, which can lead to biases and inconsistencies in their evaluations \cite{ref117}. Additionally, these metrics can be computationally expensive to compute, especially when evaluating large-scale datasets or models \cite{ref118}.

Moreover, there is an ongoing debate regarding the reliability and robustness of automated evaluation metrics, particularly in the face of adversarial attacks or perturbations \cite{ref119}. Studies have shown that existing metrics, both word-based and neural-based, can be fragile and inconsistent in their assessments, especially when faced with subtle changes to the generated text \cite{ref120}.

To address these challenges, researchers have explored ensemble and hybrid approaches that combine multiple evaluation metrics, leveraging the strengths of different methods to provide more comprehensive and reliable assessments \cite{ref116}. Additionally, there have been efforts to develop new evaluation frameworks that go beyond the traditional single-score approach and provide more fine-grained and interpretable assessments \cite{ref121,ref122}.

In summary, the landscape of automated evaluation metrics for natural language generation has evolved significantly, with the emergence of neural-based approaches that have demonstrated superior performance compared to traditional word-based and grammar-based metrics. However, the reliability and robustness of these metrics remain a pressing concern, and ongoing research is focused on addressing these challenges through the development of more comprehensive and contextualized evaluation frameworks.

\subsection{LLM-based Evaluators}

The use of large language models (LLMs) as alternative evaluators for open-ended generation tasks has emerged as a promising approach to address the limitations of traditional automatic metrics and human evaluations. LLMs, with their remarkable language understanding and generation capabilities, have the potential to provide more reliable and scalable evaluations compared to conventional methods.

One of the key benefits of using LLM-based evaluators is their ability to capture the nuances and complexities of open-ended text generation tasks, which often pose challenges for traditional automatic metrics. These metrics, such as BLEU, METEOR, and ROUGE, tend to focus on lexical and surface-level similarities, failing to account for the semantic and contextual aspects of the generated text \cite{ref123}. In contrast, LLM-based evaluators can leverage their deep understanding of language to assess the quality, coherence, and relevance of the generated text in a more holistic manner.

Furthermore, LLM-based evaluators can provide a more scalable and cost-effective solution compared to human evaluations, which can be time-consuming, labor-intensive, and subject to individual biases \cite{ref66}. By employing LLMs as evaluators, researchers and practitioners can collect large-scale feedback on generated content, enabling more comprehensive and data-driven assessments.

However, the use of LLM-based evaluators is not without its own set of challenges and considerations. One of the primary concerns is the potential for inherent biases within the LLMs themselves, which can be inherited from their training data and the way they are fine-tuned \cite{ref124}. These biases can lead to inconsistent or unreliable evaluations, particularly in cases where the generated text exhibits sensitive or controversial content.

Another challenge lies in the alignment between LLM-based evaluations and human preferences. LLMs, while highly capable, may not always capture the nuanced and subjective nature of human judgments, leading to discrepancies between model-based and human-based evaluations \cite{ref125}. Addressing this misalignment is crucial for ensuring that the use of LLM-based evaluators truly reflects the preferences and standards of the target audience.

To mitigate these challenges, researchers have explored various strategies and techniques. One approach is to calibrate and fine-tune LLM-based evaluators using human-provided feedback and carefully curated datasets \cite{ref126,ref107}. By incorporating human judgment and preferences into the evaluation process, these techniques aim to improve the reliability and alignment of LLM-based assessments.

Another approach is to leverage ensemble methods and multi-faceted evaluation frameworks that combine the strengths of LLM-based evaluators and human experts \cite{ref66,ref78}. These hybrid approaches can help address the limitations of individual evaluators and provide a more comprehensive assessment of the generated text.

Additionally, researchers have explored the use of probing techniques and interpretability methods to better understand the decision-making processes of LLM-based evaluators \cite{ref127}. By gaining insights into the internal workings of these models, researchers can identify and mitigate potential biases, leading to more reliable and transparent evaluations.

In conclusion, the use of LLM-based evaluators has the potential to revolutionize the way we assess open-ended generation tasks, offering scalable and cost-effective solutions. However, addressing the challenges of inherent biases, alignment with human preferences, and ensuring the reliability of these evaluators remains an active area of research. Continued efforts in this direction, leveraging techniques such as calibration, ensemble methods, and interpretability analysis, can unlock the full potential of LLM-based evaluators and foster more rigorous and reliable assessments of natural language generation tasks.

\subsection{Collaborative Evaluation Frameworks}

The rapid advancement of large language models (LLMs) has led to a growing need for reliable and comprehensive evaluation methods that can assess the quality of their generated outputs. Traditional human evaluation approaches, while valuable, can be time-consuming, expensive, and prone to inconsistencies \cite{ref128}. To address these limitations, researchers have explored the potential of using LLMs themselves as evaluators, leveraging their language understanding and generation capabilities to provide more scalable and cost-effective assessment \cite{ref129}.

However, relying solely on a single LLM as an evaluator often leads to biases, unreliable judgments, and a lack of comprehensive assessment \cite{ref130}. To overcome these challenges, the concept of collaborative evaluation has emerged as a promising approach, where human and LLM-based evaluators work together to assess the quality of generated outputs \cite{ref66}.

In a collaborative evaluation framework, the strengths of human evaluators, such as their domain expertise, contextual understanding, and nuanced judgment, are combined with the scalability and efficiency of LLM-based evaluators. By leveraging the complementary capabilities of both human and machine evaluators, collaborative frameworks aim to provide more reliable, comprehensive, and efficient assessments of LLM-generated outputs.

One key aspect of collaborative evaluation is the design of the interaction between human and LLM-based evaluators. \cite{ref131} proposes a multi-agent discussion framework that emulates the peer review process, where multiple LLMs act as evaluators and engage in a deliberative process to reach a consensus on the quality of a given text. This approach allows the evaluators to draw on their individual strengths, share their perspectives, and refine their assessments through iterative discussions.

Similarly, \cite{ref132} introduces a multi-agent referee team, called ChatEval, that autonomously discusses and evaluates the quality of generated responses by leveraging the diverse capabilities and expertise of its members. By harnessing the synergies of multiple LLMs, these multi-agent frameworks aim to transcend the limitations of single-agent approaches and provide more reliable and human-like evaluations.

Another key aspect of collaborative evaluation is the design of the evaluation criteria and the mechanisms for aligning them with human preferences. \cite{ref78} presents a framework that iteratively aligns LLM-based evaluators with human preference through a process of hierarchical criteria decomposition. By breaking down the evaluation task into finer-grained criteria and aggregating them according to estimated human preferences, HD-Eval aims to enhance the alignment and coverage of LLM-based evaluations.

Similarly, \cite{ref133} proposes a prompt-based approach called CoAScore, which utilizes multi-aspect knowledge through a Chain-of-Aspects prompting framework to assess the quality of natural language generation (NLG) outputs. By capturing the rich correlations between various aspects, CoAScore aims to provide more comprehensive and accurate evaluations compared to independent aspect assessments.

Collaborative evaluation frameworks also need to address the challenges of ensuring the reliability, consistency, and transparency of the evaluation process. \cite{ref130} highlights the potential biases in human and LLM-based evaluations, where factors like stylistic aspects and grammatical correctness may be overemphasized compared to factual accuracy. To address these biases, collaborative frameworks may need to incorporate mechanisms for detecting and mitigating potential sources of bias, as well as providing transparent explanations for their evaluations.

Furthermore, the design of collaborative evaluation frameworks must consider the scalability and efficiency of the evaluation process, particularly as the number of LLM models and the complexity of tasks continue to grow. \cite{ref76} presents a scalable and modular framework that integrates various evaluation methodologies, including dynamic evaluations, while also addressing concerns related to data contamination and efficient model inference.

In summary, collaborative evaluation frameworks that integrate human and LLM-based evaluators hold great promise in addressing the limitations of traditional evaluation approaches. By leveraging the complementary strengths of both human and machine evaluators, these frameworks aim to provide more reliable, comprehensive, and efficient assessments of LLM-generated outputs. However, the design of such frameworks must carefully consider the challenges of aligning evaluation criteria with human preferences, mitigating biases, and ensuring scalability and efficiency. Continued research and development in this area will be crucial for the effective and responsible deployment of LLMs in real-world applications.

\subsection{Robustness and Reliability of Evaluations}

The robustness and reliability of the evaluation methods used to assess the performance of large language models (LLMs) are of paramount importance, as these evaluations form the foundation for understanding the capabilities and limitations of these powerful models. While the existing literature on LLM evaluation has primarily focused on task-level and societal-level assessments, the impact of adversarial perturbations, model biases, and the consistency of evaluations across different settings have emerged as critical factors that warrant further investigation \cite{ref134}.

Adversarial perturbations, which are carefully crafted input modifications that can cause significant changes in model outputs, pose a significant threat to the reliability of LLM evaluations \cite{ref135}. These perturbations can exploit vulnerabilities in the models, leading to incorrect predictions that may not be readily apparent to human evaluators. Robust evaluation methods must, therefore, consider the impact of such adversarial inputs and develop strategies to identify and mitigate their influence.

Moreover, the biases inherent in the training data and the model architecture can also affect the reliability of LLM evaluations \cite{ref136}. These biases can lead to systematic errors or inconsistencies in the model's performance, particularly in domain-specific applications or when dealing with underrepresented groups. Comprehensive evaluation frameworks must, therefore, incorporate methods for detecting and addressing these biases, ensuring that the assessment of LLM performance is not unduly influenced by such factors.

The consistency of evaluations across different settings is another crucial aspect of ensuring the robustness and reliability of LLM assessments \cite{ref137}. Evaluations conducted on a single dataset or in a specific context may not accurately reflect the model's performance in other settings, particularly when dealing with domain shifts or changes in the input distribution. Robust evaluation frameworks should, therefore, consider the use of diverse benchmark datasets, simulated environments, and real-world scenarios to assess the model's performance and generalization capabilities across a range of conditions.

One approach to enhancing the robustness and reliability of LLM evaluations is the use of ensemble and hybrid evaluation methods \cite{ref138}. By combining the strengths of different evaluation techniques, such as human assessment and automated metrics, these approaches can provide a more comprehensive and reliable assessment of model performance. Additionally, the use of uncertainty quantification techniques, such as Monte Carlo dropout or Bayesian neural networks, can help to capture the inherent uncertainty in the model's predictions, providing a more nuanced understanding of its reliability \cite{ref139}.

Another promising direction for improving the robustness and reliability of LLM evaluations is the incorporation of causal reasoning and multi-modal evaluation \cite{ref136,ref140}. By considering the causal relationships between the input features and the model's predictions, as well as by incorporating diverse data modalities beyond text, such as images, audio, and video, researchers can develop more comprehensive and reliable evaluation frameworks that better capture the holistic capabilities of LLMs.

Furthermore, the open-source initiatives and collaborative efforts in the LLM evaluation community, such as the GPT-Fathom benchmarking suite and the PiCO peer-review-based evaluation approach, are important steps towards enhancing the robustness and reliability of LLM assessments \cite{ref83,ref107}. By leveraging the collective expertise of researchers, practitioners, and stakeholders, these initiatives can help to identify and address the key challenges in LLM evaluation, leading to more trustworthy and reliable assessments of model performance.

In conclusion, the robustness and reliability of LLM evaluations are critical for ensuring the responsible development and deployment of these powerful models. By addressing the impact of adversarial perturbations, model biases, and the consistency of evaluations across different settings, researchers can develop more comprehensive and reliable assessment frameworks that better capture the capabilities and limitations of LLMs. The integration of causal reasoning, multi-modal evaluation, and collaborative efforts in the LLM evaluation community hold promise for further enhancing the robustness and reliability of these assessments, ultimately paving the way for the safe and responsible use of LLMs in a wide range of applications.

\subsection{Ensemble and Hybrid Evaluation Approaches}

As the field of large language model (LLM) evaluation continues to evolve, researchers have recognized the potential of ensemble and hybrid approaches to enhance the reliability and comprehensiveness of assessments. These strategies leverage the complementary strengths of various evaluation methods, including both human and LLM-based evaluators, to provide a more holistic and robust evaluation of LLM performance.

One of the key advantages of ensemble and hybrid approaches is their ability to mitigate the limitations of individual evaluation methods. Human evaluations, while considered the gold standard, can be subject to biases, inconsistencies, and resource constraints \cite{ref141,ref142}. On the other hand, LLM-based evaluators, while offering scalability and efficiency, can also exhibit their own biases and lack of alignment with human preferences \cite{ref143,ref103}.

By combining multiple evaluation methods, ensemble and hybrid approaches aim to leverage the strengths of each approach while compensating for their weaknesses. For instance, \cite{ref144} proposes a Collaborative Evaluation (CoEval) pipeline that integrates the capabilities of both humans and LLMs, where LLMs generate initial evaluations, and humans then scrutinize and refine the results. This approach capitalizes on the speed and scalability of LLM-based evaluations while maintaining the contextual understanding and nuanced judgments of human evaluators.

Similarly, \cite{ref145} introduces the concept of ``ensemble interpretation,'' which integrates the results of various interpretation methods to achieve more stable and reliable explanations for LLM behavior. By combining multiple perspectives, ensemble interpretation can provide a more comprehensive understanding of the models' inner workings and decision-making processes, ultimately enhancing the trustworthiness and interpretability of LLM evaluations.

Another promising approach is the use of uncertainty-based ensemble methods, which leverage the diversity and disagreement among LLM-based evaluators to improve the reliability of assessments. \cite{ref143} presents Pairwise-preference Search (PairS), an uncertainty-guided search method that employs LLMs to conduct pairwise comparisons and efficiently rank candidate texts. By quantifying the uncertainty and transitivity of LLM evaluations, PairS demonstrates significant improvements over direct scoring approaches, better aligning with human preferences.

Beyond ensemble methods, researchers have also explored hybrid approaches that combine LLM-based evaluations with external knowledge sources or specialized reasoning capabilities. For example, \cite{ref146} utilizes the Shapley value to fairly quantify the contributions of individual prompts within an ensemble, helping to identify beneficial or detrimental prompts and guide prompt selection in data markets.

Additionally, \cite{ref147} proposes a framework that leverages a different LLM, along with human-in-the-loop verification, to automatically generate and apply auditing probes to a target LLM. This approach, which is generalizable to various LLMs, aims to enhance the reliability and transparency of LLM evaluation by introducing a layer of external validation and cross-model verification.

The integration of ensemble and hybrid approaches in LLM evaluation has demonstrated promising results in improving the reliability, robustness, and interpretability of assessments. By harnessing the complementary strengths of different evaluation methods, researchers can develop more comprehensive and trustworthy evaluation frameworks that better capture the nuanced capabilities and limitations of large language models.

As the field of LLM evaluation continues to advance, the adoption of ensemble and hybrid approaches will likely become increasingly crucial. These strategies can not only enhance the accuracy and consistency of LLM assessments but also provide valuable insights into the models' reasoning processes, biases, and alignment with human preferences. Continued research and development in this area will be essential to ensure the responsible and effective deployment of large language models across a wide range of applications.

\section{Assessing LLM Robustness, Reliability, and Alignment}

\subsection{Measuring and Quantifying Robustness and Reliability of LLMs}

The robustness and reliability of large language models (LLMs) have become increasingly crucial as these models are deployed in real-world applications. Quantifying these properties is essential for ensuring the safety and trustworthiness of LLMs. Researchers have proposed various techniques and metrics to assess the robustness and reliability of LLMs, addressing issues such as biases, inconsistencies, and safety concerns.

One prominent approach to measuring the robustness of LLMs is through the use of adversarial attacks \cite{ref148}. These attacks aim to identify vulnerabilities in the models by introducing small, carefully crafted perturbations to the input that can significantly alter the model's output. Metrics such as the attack success rate (ASR) have been used to quantify the robustness of LLMs to these adversarial perturbations \cite{ref149}. However, recent studies have highlighted the limitations of relying solely on ASR, as it may not capture the semantic quality of the model's responses \cite{ref150}. To address this, researchers have proposed incorporating additional metrics that assess the model's reasoning and the coherence of its outputs.

Another important aspect of quantifying the reliability of LLMs is the assessment of biases, both explicit and implicit \cite{ref151,ref152}. Explicit biases refer to the model's tendency to generate content that reflects known societal biases, such as gender or racial stereotypes. Metrics like the demographic parity score have been used to measure the extent of these biases \cite{ref153}. Implicit biases, on the other hand, are more subtle and can be harder to detect, as they are often reflected in the model's decision-making processes rather than its outputs. Researchers have proposed techniques like the Implicit Association Test (IAT) to uncover these implicit biases in LLMs \cite{ref152}.

Evaluating the consistency and safety of LLMs is another crucial aspect of quantifying their reliability. Inconsistencies can arise when an LLM generates different outputs for the same or similar inputs, which can undermine user trust and confidence in the model's capabilities \cite{ref154}. Metrics like the agreement between multiple model outputs or the perplexity of the model's predictions have been used to assess consistency \cite{ref155}. Safety concerns, on the other hand, relate to the potential for LLMs to generate harmful or inappropriate content, such as hate speech or misinformation. Researchers have proposed using detectors to identify and flag such content \cite{ref156}.

Developing robust and reliable LLMs is not without its challenges. One key challenge is the inherent biases present in the training data used to develop these models \cite{ref48}. These biases can be amplified and propagated through the model, leading to undesirable outputs. Addressing this challenge requires careful curation and cleaning of training data, as well as the development of debiasing techniques \cite{ref148}.

Another challenge is the difficulty in comprehensively evaluating the robustness and reliability of LLMs. The vast and diverse nature of language tasks and the complexity of natural language understanding make it challenging to develop a single, comprehensive evaluation framework \cite{ref157}. Researchers have proposed the use of multi-faceted evaluation approaches, combining various metrics and benchmarks to assess different aspects of an LLM's performance \cite{ref158}.

Despite these challenges, the ongoing research efforts to measure and quantify the robustness and reliability of LLMs are crucial for ensuring the safe and responsible development of these powerful models. As LLMs continue to be deployed in increasingly high-stakes applications, the need for robust and reliable evaluation techniques will only become more pressing. By addressing these challenges, researchers can help guide the evolution of LLMs towards more trustworthy and beneficial outcomes for society.

\subsection{Evaluating LLM Alignment with Human Values and Preferences}

Evaluating the alignment of large language models (LLMs) with human values, preferences, and societal norms is a critical challenge that has garnered significant attention in the research community. Given the growing prominence of LLMs in various applications, it is essential to ensure these models behave in accordance with human intentions and expectations.

A key aspect of this evaluation process is assessing the extent to which LLMs have internalized and can adhere to the diverse range of human values that exist across different cultures, demographics, and belief systems. This is a complex and multifaceted task, as human values are often subjective, context-dependent, and can evolve over time \cite{ref159}.

One approach to evaluating LLM alignment with human values involves leveraging existing frameworks from the field of moral psychology, such as Moral Foundation Theory \cite{ref160}. By mapping LLM outputs to the dimensions of moral values, such as care, fairness, loyalty, authority, and purity, researchers can gain insights into the models' underlying value systems and how they compare to those of human populations \cite{ref161}.

Additionally, researchers have explored the use of value-aligned benchmarks, where LLMs are tasked with generating responses that adhere to specific value principles, such as promoting social good, minimizing harm, or respecting individual rights \cite{ref162}. By carefully designing these benchmarks and evaluating LLM performance, it is possible to assess the models' ability to reason about and act in accordance with human values \cite{ref159}.

Another avenue of research involves studying the personalization and customization of LLM alignment to individual user preferences and values \cite{ref163}. This is particularly important given the diverse nature of human values and the potential for LLMs to be deployed in a wide range of contexts and applications \cite{ref164}.

Researchers have also explored the use of collaborative evaluation frameworks, where human and LLM-based evaluators work together to assess the alignment of generated outputs with human values and preferences \cite{ref78}. By leveraging the complementary strengths of these evaluation approaches, it is possible to achieve more comprehensive and reliable assessments of LLM alignment \cite{ref165}.

However, evaluating LLM alignment with human values is not without its challenges. One key issue is the inherent complexity and diversity of human values, which can make it difficult to define a universal set of alignment criteria \cite{ref166}. Additionally, the dynamic nature of human values and their sensitivity to contextual factors can pose significant challenges in ensuring LLMs remain aligned over time \cite{ref167}.

Furthermore, the potential for LLMs to exhibit biases and inconsistencies in their value-related reasoning and decision-making has been a subject of extensive research \cite{ref168}. Addressing these issues requires a multifaceted approach, including the development of advanced evaluation methodologies, the incorporation of causal reasoning capabilities, and the exploration of multi-modal evaluation frameworks \cite{ref169,ref140}.

In summary, the evaluation of LLM alignment with human values and preferences is a crucial and complex challenge that has gained significant attention in the research community. Leveraging insights from moral psychology, value-aligned benchmarking, and collaborative evaluation frameworks, researchers are working to develop robust and comprehensive assessment methodologies to ensure LLMs behave in accordance with human intentions and expectations. However, challenges remain, including the inherent complexity of human values, the dynamic nature of societal norms, and the potential for biases and inconsistencies in LLM reasoning. Addressing these challenges will be essential for the safe and responsible deployment of LLMs in real-world applications.

\subsection{Assessing LLM Hallucination and Mitigating Strategies}

The phenomenon of hallucination in large language models (LLMs) has become a pressing concern as these models become increasingly integrated into various applications and decision-making processes. Hallucination refers to the tendency of LLMs to generate factually incorrect but seemingly plausible outputs, posing significant risks to the reliability and trustworthiness of these models.

Researchers have made substantial efforts to understand and mitigate the hallucination problem in LLMs. \cite{ref170} provides a comprehensive analysis of hallucination, focusing on three key aspects: detection, source, and mitigation. The authors construct the HaluEval 2.0 benchmark to evaluate hallucination detection and propose a simple yet effective method for detecting hallucinations in LLM outputs.

Understanding the underlying causes of hallucination is crucial for developing effective mitigation strategies. \cite{ref171} introduces MIND, an unsupervised training framework that leverages the internal states of LLMs to detect hallucinations in real-time without the need for manual annotations. The authors also present HELM, a benchmark for evaluating hallucination detection across multiple LLMs, providing valuable insights into the inference process and the internal states of the models.

Beyond detection, researchers have also explored methods to directly mitigate hallucinations in LLMs. \cite{ref172} proposes FEWL, a hallucination metric that does not require gold-standard answers, addressing the challenge of costly and error-prone human annotations. FEWL leverages the answers from off-the-shelf LLMs as a proxy for gold-standard answers and quantifies the expertise of these reference models to provide more accurate hallucination measures.

Another approach to mitigating hallucinations is through the use of uncertainty quantification. \cite{ref173} introduces a novel fact-checking and hallucination detection pipeline based on token-level uncertainty quantification. The authors demonstrate that uncertainty scores can be used to identify unreliable predictions and detect hallucinations in LLM outputs.

While these efforts have made significant progress in addressing hallucination, the problem remains a challenging one. \cite{ref174} provides a comprehensive survey of the detection, explanation, and mitigation of hallucination, highlighting the unique challenges posed by LLMs and the need for further research.

One promising direction in mitigating hallucinations is the use of self-reflection and self-correction mechanisms. \cite{ref175} presents an interactive self-reflection methodology that incorporates knowledge acquisition and answer generation. By leveraging the interactivity and multitasking ability of LLMs, the approach steadily enhances the factuality, consistency, and entailment of the generated answers, effectively reducing hallucination.

Another avenue for hallucination mitigation is through the use of knowledge-based approaches. \cite{ref176} introduces a dataset, DelucionQA, that captures hallucinations made by retrieval-augmented LLMs in a domain-specific question-answering task. The authors also propose a set of hallucination detection methods to serve as baselines for future research.

The integration of causal reasoning into LLMs has also been explored as a potential strategy to mitigate hallucinations. \cite{ref177} proposes an association analysis approach that combines hallucination level quantification and hallucination reason investigation. The method reveals potential deficiencies in LLMs' commonsense memorization, relational reasoning, and instruction following, providing guidance for improving the models' capabilities and reducing hallucinations.

\subsection{Evaluating LLM Transparency and Explainability}

The transparency and explainability of large language models (LLMs) have become increasingly crucial as these models continue to demonstrate impressive capabilities across a wide range of tasks \cite{ref28,ref178}. Transparency in LLMs refers to the ability to understand the inner workings and decision-making processes of the models, while explainability focuses on the models' capacity to provide clear and understandable justifications for their outputs \cite{ref179}.

Evaluating the transparency and explainability of LLMs is a multifaceted challenge that requires a comprehensive approach \cite{ref180}. One important aspect is assessing the models' ability to provide transparent and understandable reasoning for their outputs, which can help build trust and facilitate the integration of LLMs into real-world applications, particularly in sensitive domains such as healthcare and finance \cite{ref181}.

Existing research has explored various methods for evaluating the transparency and explainability of LLMs. One approach involves leveraging interpretability techniques to analyze the models' inner workings and decision-making processes \cite{ref182,ref183}. These techniques can include examining the importance and contribution of individual model components, such as attention heads and neurons, in the overall prediction \cite{ref184,ref185}.

Another key aspect of evaluating LLM transparency and explainability is assessing the quality and reliability of the explanations generated by the models \cite{ref186}. Researchers have proposed various metrics and frameworks to measure the faithfulness, robustness, and alignment of these explanations with human understanding \cite{ref187,ref188}.

Additionally, the evaluation of LLM transparency and explainability should consider the context and intended use of the models \cite{ref189}. Different applications may require different levels of transparency and explainability, and the evaluation methods should be tailored accordingly \cite{ref190}.

One promising approach to evaluating LLM transparency and explainability is the use of collaborative evaluation frameworks, where both human and LLM-based evaluators work together to assess the quality and reliability of the models' outputs \cite{ref66}. This can help to leverage the complementary strengths of humans and LLMs, leading to more comprehensive and trustworthy evaluations \cite{ref191}.

Furthermore, the evaluation of LLM transparency and explainability should also consider the potential biases and limitations of the models \cite{ref192}. Researchers have highlighted the importance of evaluating the models' ability to provide transparent and unbiased explanations, particularly in domains where fairness and equity are crucial \cite{ref193}.

In summary, the evaluation of LLM transparency and explainability is a multifaceted challenge that requires a comprehensive approach. Existing research has explored various methods, including the analysis of model components, the assessment of explanation quality and reliability, and the consideration of context-specific requirements and potential biases. As LLMs continue to play an increasingly important role in various domains, the importance of evaluating their transparency and explainability will only grow, paving the way for more trustworthy and accountable AI systems.

\subsection{Leveraging Ensemble and Uncertainty-based Approaches}

The use of ensemble methods and uncertainty quantification techniques has emerged as a promising avenue for enhancing the robustness, reliability, and safety of large language models (LLMs). These approaches leverage the collective strengths of multiple models or model components to provide more reliable and trustworthy predictions, while also accounting for the inherent uncertainties associated with LLM outputs.

Ensemble methods, such as deep ensembles, have demonstrated their effectiveness in improving the overall performance and uncertainty estimation of deep learning models, including LLMs \cite{ref194,ref195,ref196}. By training multiple models with different initializations, architectures, or training data, ensemble methods can capture a more comprehensive range of possible outputs and their associated uncertainties. This diversity helps to mitigate the overconfidence and biases that often plague individual LLMs, leading to more reliable and well-calibrated uncertainty estimates.

Furthermore, ensemble-based approaches can be combined with uncertainty quantification techniques to provide a more holistic assessment of LLM reliability. Approaches such as Monte Carlo dropout \cite{ref194,ref197} and Bayesian neural networks \cite{ref198,ref195} can be incorporated into ensemble frameworks to capture both aleatoric (data-dependent) and epistemic (model-dependent) uncertainties. By considering these different sources of uncertainty, ensemble-based methods can more accurately identify the limitations and potential failure modes of LLMs, enabling more robust and reliable decision-making.

One key advantage of ensemble-based approaches is their ability to provide uncertainty estimates at both the individual output and the aggregate level. This granularity allows for more nuanced decision-making, where high-confidence outputs can be trusted, while uncertain outputs can be flagged for further scrutiny or human intervention. This is particularly important in safety-critical applications, where the consequences of LLM errors can be severe.

However, the adoption of ensemble-based approaches for LLM evaluation also comes with its own set of challenges. The increased computational and memory requirements associated with training and deploying multiple models can be prohibitive, especially for large-scale LLMs. To address this, recent work has explored more efficient ensemble-based techniques, such as the use of low-rank adaptation (LoRA) ensembles \cite{ref199} and early-exit cascades \cite{ref200}. These methods aim to strike a balance between the benefits of ensemble-based uncertainty estimation and the practical considerations of deploying LLMs in real-world settings.

Another challenge lies in the interpretability and transparency of ensemble-based uncertainty estimates. While these methods can provide valuable insights into the reliability of LLM outputs, the underlying mechanisms that drive the uncertainty estimates may not always be clear. This can hinder the ability of practitioners to understand and trust the uncertainty quantification, especially in high-stakes applications. Addressing this challenge requires further research into explainable and interpretable uncertainty quantification techniques for LLMs \cite{ref201,ref202}.

Despite these challenges, the integration of ensemble methods and uncertainty quantification techniques holds significant promise for enhancing the robustness, reliability, and safety of LLMs. By leveraging the complementary strengths of these approaches, researchers and practitioners can develop more trustworthy and accountable LLM systems that are better equipped to handle the complexities and uncertainties inherent in real-world language-based tasks. As the field of LLM development and evaluation continues to evolve, the continued exploration and advancement of ensemble-based and uncertainty-aware techniques will be crucial for realizing the full potential of these powerful language models.

\subsection{Personalization and Customization of LLM Alignment}

The personalization and customization of LLM alignment present both opportunities and challenges in ensuring the responsible and effective deployment of these powerful models. As LLMs become increasingly integrated into various applications, addressing the diversity of user preferences and values has emerged as a critical consideration \cite{ref163}.

Existing alignment techniques, such as Reinforcement Learning with Human Feedback (RLHF), often aim to optimize LLMs for a generic or aggregated set of preferences, overlooking the inherent heterogeneity in human values and perspectives. This one-size-fits-all approach risks undermining the unique preferences and needs of individual users, potentially leading to suboptimal or even harmful outcomes in personalized applications \cite{ref203}.

To address this challenge, researchers have explored more nuanced approaches that embrace the diversity of human preferences. One such approach is the concept of ``Group Preference Optimization'' (GPO), which steers LLMs to align with the preferences of specific user groups or demographics in a few-shot manner \cite{ref204}. By incorporating group-level preferences into the alignment process, GPO aims to balance the needs of diverse user segments, potentially unlocking more personalized and equitable LLM-powered experiences.

However, the personalization of LLM alignment introduces additional complexities and trade-offs. A key consideration is the balance between individual-level personalization and maintaining societal-level alignment with broader human values and norms \cite{ref163}. Overly personalized LLMs may risk diverging from societal consensus, potentially leading to harmful or unethical behavior in certain contexts.

To navigate this delicate balance, researchers have proposed frameworks that integrate both individual and societal-level preferences into the LLM alignment process. One such approach is the ``Heterogeneous Value Alignment Evaluation'' (HVAE) system, which assigns LLMs with different social values (e.g., individualistic vs.~collectivistic) and measures their ability to align with these varying value orientations \cite{ref162}. By evaluating LLMs' value rationality across diverse value systems, HVAE provides insights into the models' ability to balance individual and societal-level alignment.

Furthermore, the personalization of LLM alignment raises concerns about user privacy, model ownership, and the potential for manipulation. Recent work has explored approaches that empower users to own and customize their LLMs, such as ``One PEFT Per User'' (OPPU), which employs personalized parameter-efficient fine-tuning (PEFT) modules to store user-specific behavior patterns and preferences \cite{ref205}. By granting users control over their personal LLMs, OPPU aims to address the limitations of existing prompt-based personalization methods and ensure more transparent and user-centric customization.

Alongside these technical innovations, the personalization of LLM alignment also necessitates deeper consideration of the societal and ethical implications. Researchers have highlighted the importance of acknowledging the pluralistic nature of human values and the potential risks of LLMs being aligned with the preferences of specific individuals or groups \cite{ref163}. The concern is that overly personalized LLMs could exacerbate existing societal biases and inequalities, or even introduce new forms of discrimination.

To mitigate these risks, scholars have proposed policy frameworks that establish boundaries and guidelines for the personalization of LLM alignment. These frameworks emphasize the need to balance user-level customization with the preservation of societal-level norms and values, ensuring that personalized LLMs do not deviate too far from widely accepted ethical principles \cite{ref163}.

In summary, the personalization and customization of LLM alignment present both opportunities and challenges. While tailoring LLMs to individual user preferences can enhance user experience and engagement, it must be carefully balanced with the maintenance of societal-level alignment to avoid unintended negative consequences. Ongoing research explores innovative approaches, such as GPO, HVAE, and OPPU, to navigate this delicate balance, while also considering the broader societal and ethical implications of personalized LLM alignment. As LLMs become increasingly ubiquitous, addressing these complex issues will be crucial for ensuring the responsible and equitable deployment of these powerful models.

\subsection{Benchmarking and Evaluation Frameworks for Assessing LLM Robustness and Alignment}

The assessment of the robustness, reliability, and alignment of large language models (LLMs) is a critical task that requires comprehensive and multifaceted evaluation frameworks. These frameworks must not only capture the technical capabilities of LLMs but also address the broader societal implications of their deployment.

Existing evaluation platforms, such as the HuggingFace leaderboard, have primarily focused on task-level benchmarks, which provide a limited perspective on the overall performance and trustworthiness of LLMs \cite{ref65}. While these task-level assessments are valuable, they often overlook the nuanced aspects of LLM behavior, such as their ability to handle adversarial inputs, maintain consistent responses across contexts, and adhere to ethical principles.

To address this gap, researchers have proposed a range of specialized benchmarks and evaluation frameworks that delve deeper into the assessment of LLM robustness, reliability, and alignment. For instance, the TrustLLM benchmark \cite{ref206} encompasses multiple dimensions of trustworthiness, including truthfulness, safety, fairness, robustness, privacy, and machine ethics. By rigorously evaluating LLMs across these dimensions, the TrustLLM framework provides a more holistic understanding of their capabilities and limitations.

Similarly, the Flames benchmark \cite{ref207} focuses on the alignment of LLMs with Chinese-specific values and norms, going beyond the more general principles of being helpful, honest, and harmless (3H). By incorporating complex scenarios and adversarial prompts, Flames challenges LLMs to demonstrate genuine value alignment, rather than simply surface-level compliance.

The importance of comprehensive benchmarking and evaluation frameworks is further highlighted by the work of \cite{ref157}, which emphasizes the need for rigorous reliability testing to ensure the fairness, robustness, and transparency of NLP systems, including LLMs. The authors argue that reliability testing, complemented by interdisciplinary collaboration, is essential for the development of industry standards and the enactment of effective regulations.

Moreover, the assessment of LLM robustness and alignment must extend beyond narrow task-specific evaluations and address the models' behavior in more complex and realistic settings. The AgentBoard framework \cite{ref208} tackles this challenge by evaluating LLMs as general-purpose agents, rather than focusing solely on individual tasks. AgentBoard's fine-grained progress rate metric and comprehensive evaluation toolkit enable a deeper understanding of LLM capabilities and limitations in multi-turn interactions.

The need for collaborative and transparent evaluation approaches is further underscored by the work on \cite{ref80}, which introduces a novel framework that leverages a peer-review mechanism to automatically evaluate and rank LLMs. By harnessing the collective wisdom of multiple LLM evaluators, the PRE framework aims to mitigate the biases and limitations inherent in single-model evaluations.

Complementing these efforts, researchers have also explored the integration of uncertainty quantification \cite{ref209} and causal reasoning \cite{ref210} into LLM evaluation, recognizing the importance of these aspects in ensuring the reliability and safety of these models.

As the development of LLMs continues to accelerate, the need for comprehensive and rigorous evaluation frameworks becomes increasingly urgent. These frameworks must not only assess the technical capabilities of LLMs but also their alignment with societal values, their robustness to adversarial inputs, and their transparency in decision-making. By fostering collaborative efforts, interdisciplinary collaboration, and the continuous refinement of evaluation methodologies, the research community can play a pivotal role in shaping the responsible development and deployment of LLMs, ensuring their beneficial impact on society.

\section{Evaluating LLMs in Specialized Domains}

\subsection{Medical Applications}

The evaluation of large language models (LLMs) in medical applications has gained significant attention due to the potential of these models to revolutionize various aspects of healthcare delivery. LLMs have demonstrated remarkable capabilities in tasks such as clinical decision support, disease diagnosis, treatment recommendations, and medical workflow automation \cite{ref211,ref212}.

In the domain of clinical decision support, LLMs have shown promise in assisting clinicians by providing evidence-based recommendations and explanations for medical decisions. For instance, a study \cite{ref211} explored the effectiveness and reliability of LLM-generated explanations for diagnoses based on patient complaints. The researchers found that LLM explanations significantly increased doctors' agreement rates with given diagnoses and highlighted potential errors in LLM outputs, underscoring the need for careful integration and evaluation to ensure patient safety and optimal clinical utility.

Regarding disease diagnosis, LLMs have been explored for their potential to enhance the accuracy and efficiency of the diagnostic process. A recent framework \cite{ref87} proposed an automatic evaluation paradigm to assess the clinical capabilities of LLMs in delivering disease diagnosis and treatment recommendations. The framework leverages standardized patient scenarios and a multi-agent simulation environment to evaluate the LLM's adherence to clinical practice guidelines and its ability to engage in interactive clinical dialogues.

Furthermore, LLMs have shown promise in automating medical workflow tasks, such as clinical documentation and medication management. A study \cite{ref212} highlighted the potential of LLM agents to interact with stakeholders in open-ended conversations and influence clinical decision-making. The researchers emphasized the need for robust, real-world clinical evaluations to assess the performance of LLM agents within the clinical workflow.

Beyond these specific applications, LLMs have demonstrated their versatility in various medical domains, including oncology \cite{ref213}, mental health \cite{ref214}, and cardiovascular care \cite{ref215}. These studies have showcased the ability of LLMs to process and integrate diverse medical data, such as electronic health records, medical images, and scientific literature, to provide personalized recommendations and insights.

However, the integration of LLMs in medical settings is not without its challenges. Concerns have been raised regarding the reliability, transparency, and safety of LLM-based medical applications \cite{ref35}. Researchers have emphasized the importance of rigorous evaluation frameworks that assess the model's performance, alignment with clinical guidelines, and potential for harm \cite{ref216}.

To address these challenges, recent studies have explored various approaches to enhance the reliability and trustworthiness of LLMs in medical applications. These include incorporating domain-specific knowledge and reasoning capabilities \cite{ref217,ref33}, leveraging collaborative evaluation methods that combine human and LLM-based assessments \cite{ref66}, and developing interpretable and explainable models that can provide insights into their decision-making processes \cite{ref218,ref219}.

As the field of medical AI continues to evolve, the evaluation of LLMs in specialized domains, such as healthcare, will play a crucial role in ensuring the safe and effective deployment of these powerful language models. By addressing the challenges and leveraging the strengths of LLMs, researchers and healthcare professionals can unlock the immense potential of these models to transform medical practice and improve patient outcomes.

\subsection{Educational Applications}

The rapid advancements in large language models (LLMs) have garnered significant attention in the educational domain, with researchers and educators exploring the potential of these models to enhance various aspects of learning and teaching. One prominent area of focus has been the evaluation of LLMs in educational applications, particularly in the realm of medical education, content generation, and personalized learning.

In the domain of medical education, LLMs have shown promise in improving the performance of future medical professionals. \cite{ref85} presents the development of an open-source dataset and collection of medical conversational AI models, aimed at fine-tuning LLMs for effective medical applications. The authors emphasize the need for on-premises deployable models to safeguard patient privacy, a critical concern in the healthcare domain. By fine-tuning publicly accessible pre-trained LLMs on the curated medical dataset, the researchers demonstrate the models' improved performance in tasks related to the examinations that future medical doctors must pass to achieve certification.

Similarly, \cite{ref220} explores the integration of authoritative medical textbooks into the LLM framework to enhance their proficiency in the biomedical domain. The proposed system, called LLM-AMT, combines a Query Augmenter, a Hybrid Textbook Retriever, and a Knowledge Self-Refiner to effectively incorporate domain-specific knowledge. The authors' experimental results on medical question-answering tasks show that LLM-AMT significantly outperforms specialized models pre-trained on a massive medical corpus, highlighting the value of domain-specific knowledge curation for LLM-based medical applications.

Beyond medical education, the application of LLMs in content generation for educational purposes has also garnered attention. \cite{ref221} investigates the feasibility of leveraging LLMs to facilitate faster content creation for asynchronous adult learning courses. The authors implemented a robust human-in-the-loop process to ensure the accuracy and clarity of the generated content, and their findings suggest that this approach can indeed enable efficient content creation without compromising on quality.

The potential of LLMs in personalizing learning experiences has also been explored. \cite{ref222} presents a novel task of adaptive and personalized exercise generation for online language learning. The authors combine a knowledge tracing model, which estimates each student's evolving knowledge states from their learning history, with a controlled text generation model that generates exercise sentences based on the student's current estimated knowledge state and instructor requirements. The authors demonstrate the superiority of their model in generating personalized exercises that adapt to individual student abilities, highlighting the promise of LLMs in facilitating adaptive and personalized learning.

\cite{ref223} further explores the integration of LLMs in conversation-based personalized tutoring systems. The authors discuss the design considerations for such a system, which involves a student modeling component with diagnostic capabilities and a conversation-based tutor utilizing LLMs with prompt engineering to incorporate student assessment outcomes and various instructional strategies. The authors' proof-of-concept tutoring system showcases the potential of LLMs in enabling personalization and adaptivity in educational settings.

The utilization of LLMs in educational applications also raises important considerations regarding the ethical and practical implications of these models. \cite{ref224} systematically examines the current state of research on the use of LLMs to automate educational tasks, such as question generation, feedback provision, and essay grading. The authors identify various practical and ethical challenges, including low technological readiness, lack of replicability and transparency, and insufficient privacy and beneficence considerations. These findings underscore the need for a balanced and responsible approach in integrating LLMs into educational systems, ensuring that the benefits of these models are realized while addressing the potential risks and limitations.

\subsection{Natural and Social Sciences}

The use of large language models (LLMs) in natural and social science domains holds immense potential, as these powerful models can significantly enhance various aspects of scientific research and policy analysis. LLMs have demonstrated their capabilities in tasks such as scientific literature analysis, hypothesis generation, and policy recommendation, showcasing their ability to synthesize and generate high-quality content that can accelerate the advancement of these fields.

One of the key applications of LLMs in the natural sciences is their use in analyzing scientific literature. LLMs can be employed to summarize, extract key insights, and identify emerging trends from the vast and rapidly growing body of scientific publications \cite{ref225}. This is particularly valuable in fields where the volume of literature makes it challenging for researchers to keep up with the latest developments. By leveraging LLMs, scientists can quickly and efficiently identify relevant studies, synthesize findings, and uncover new research directions, ultimately accelerating the pace of scientific discovery.

Moreover, LLMs have shown promise in aiding the process of hypothesis generation, a crucial step in the scientific method \cite{ref226}. These models can draw upon their extensive knowledge to propose novel hypotheses that could lead to groundbreaking discoveries. By combining LLMs' ability to generate hypotheses with human expertise and domain-specific knowledge, researchers can explore a wider range of possibilities and identify promising avenues for further investigation.

In the social sciences, LLMs have demonstrated their potential in policy recommendation and analysis. These models can leverage their understanding of complex social, economic, and political dynamics to generate policy proposals, assess the potential impacts of policy decisions, and identify areas for policy intervention \cite{ref227}. This can be particularly valuable in fields such as economics, public policy, and political science, where LLMs can assist policymakers in making informed decisions that address societal challenges.

However, the use of LLMs in natural and social science domains is not without its challenges. One key concern is the potential for these models to perpetuate and amplify existing biases present in their training data \cite{ref228}. If the data used to train LLMs is skewed or unrepresentative of the diversity of scientific and social perspectives, the models may exhibit biases that could lead to flawed or biased outputs. Addressing these biases and ensuring the fairness and inclusivity of LLM-based applications is crucial for their responsible deployment in these domains.

Additionally, the interpretability and transparency of LLMs' decision-making processes are crucial in scientific and policy-related applications \cite{ref229}. Researchers and policymakers must be able to understand the reasoning behind the LLMs' outputs to ensure the validity and reliability of the information they provide. Efforts to enhance the explainability and interpretability of LLMs, as well as the development of evaluation frameworks that assess these qualities, are necessary to build trust and confidence in the use of these models in natural and social science domains.

Despite these challenges, the integration of LLMs into natural and social science research and policy analysis holds significant promise. By leveraging the extensive knowledge and powerful language generation capabilities of these models, researchers and policymakers can unlock new opportunities for scientific discovery, evidence-based decision-making, and the betterment of society. As the field of LLM development continues to evolve, it is crucial to address the issues of bias, interpretability, and transparency to ensure the responsible and impactful use of these technologies in the natural and social sciences.

\subsection{Agent-based Applications}

The rapid advancements in large language models (LLMs) have enabled their integration into a wide range of agent-based applications, where these powerful models are employed as intelligent agents to interact with users and complete tasks in simulated environments. These agent-based applications leverage the language understanding and generation capabilities of LLMs to create virtual entities that can engage in naturalistic interactions, make decisions, and execute complex sequences of actions.

One prominent example of agent-based applications is the use of LLMs in video game environments \cite{ref230,ref231,ref232}. These applications harness the versatility of LLMs to enable agent-driven gameplay, where the agents navigate the game world, engage in strategic planning, and collaborate with human players or other AI agents to accomplish shared objectives. For instance, the AgentBench \cite{ref230} framework presents a diverse set of gaming environments that challenge LLM agents to demonstrate their reasoning, decision-making, and collaboration abilities in scenarios such as tower defense, resource management, and multi-agent coordination.

Similarly, the LLMArena \cite{ref231} benchmark focuses on evaluating the capabilities of LLM agents in dynamic, multi-agent gaming environments, assessing their performance in tasks that require spatial reasoning, strategic planning, numerical reasoning, risk assessment, communication, opponent modeling, and team collaboration. These evaluations not only shed light on the current limitations of LLM agents but also provide valuable insights into the future development of more sophisticated and capable autonomous agents.

Beyond gaming, LLM agents have also found applications in simulated environments that mimic real-world scenarios, such as the AndroidArena \cite{ref233} framework, which evaluates the ability of LLM agents to navigate and complete tasks within a simulated Android operating system. The challenges faced by LLM agents in this environment, such as maintaining up-to-date understanding of the vast and dynamic action space, coordinating cross-application cooperation, and adhering to user constraints, highlight the need for advancements in areas like long-term reasoning, decision-making, and adaptive exploration.

Another notable application of LLM agents is in the context of multi-agent systems, where multiple LLM-powered agents collaborate or compete to accomplish complex tasks \cite{ref234,ref235}. These multi-agent systems leverage the unique capabilities of LLMs, such as their ability to engage in natural language communication, reason about social dynamics, and learn from past experiences, to create increasingly sophisticated and adaptable agent ecosystems. For instance, the MAgIC \cite{ref234} framework employs a Probabilistic Graphical Modeling (PGM) approach to enhance the LLM agents' capabilities in navigating complex social and cognitive dimensions, enabling them to exhibit behaviors like judgment, reasoning, deception, self-awareness, cooperation, coordination, and rationality.

Similarly, the MetaAgents \cite{ref235} framework introduces the concept of collaborative generative agents, where LLM-based agents are endowed with consistent behavior patterns and task-solving abilities, and situated in a simulated job fair environment to examine their coordination skills. These advancements in multi-agent systems highlight the potential of LLM agents to engage in meaningful, task-oriented social interactions, paving the way for their integration into a wide range of real-world applications.

The evaluation of LLM agents in these diverse agent-based applications has revealed both the strengths and limitations of current LLM-powered agents. While top commercial LLMs have demonstrated impressive capabilities in areas like reasoning, planning, and decision-making \cite{ref230,ref231}, open-source LLMs often struggle to match their performance, particularly in domains that require long-term reasoning, multi-agent coordination, and adherence to complex constraints \cite{ref233}.

These findings underscore the need for continued research and development in enhancing the agent capabilities of LLMs, with a focus on areas like causal reasoning, multi-modal interaction, and the integration of external knowledge and tools. By addressing these challenges, the research community can pave the way for the seamless integration of LLM agents into a wide range of real-world applications, from smart home management to autonomous transportation systems.

\section{Interpretability and Explainability of LLMs}

\subsection{The Importance of Interpretability and Explainability in Large Language Models (LLMs)}

The rapid advancements in large language models (LLMs) have transformed the landscape of natural language processing, enabling unprecedented breakthroughs in a wide range of tasks, from text generation to question answering \cite{ref236,ref237}. As LLMs become increasingly integral to diverse applications, their ``black-box'' nature has sparked significant concerns regarding transparency and ethical use \cite{ref238,ref55}.

Interpretability and explainability are crucial in the context of LLMs, as they play a pivotal role in addressing these concerns and fostering trust in these powerful models \cite{ref236,ref237}. Interpretability refers to the ability to understand the internal workings and decision-making processes of a model, while explainability focuses on the model's capacity to provide clear and understandable explanations for its outputs \cite{ref239}.

The need for interpretability and explainability in LLMs is particularly pressing as these models are increasingly deployed in high-stakes applications, such as healthcare, finance, and legal decision-making \cite{ref240,ref28}. In these domains, transparency and accountability are paramount, as the decisions made by LLMs can have significant implications for individuals and society \cite{ref35}.

Furthermore, the complexity and scale of LLMs pose unique challenges for interpretability and explainability. Unlike simpler machine learning models, the intricate architectures and vast parameter spaces of LLMs make it difficult to understand the underlying logic and reasoning behind their outputs \cite{ref237,ref236}. This opacity can lead to concerns about the reliability, safety, and fairness of LLM-based systems, as it becomes challenging to audit and validate their behaviors \cite{ref56}.

Addressing these challenges is crucial for ensuring the responsible and ethical deployment of LLMs. Interpretability and explainability can help users, developers, and policymakers better understand the limitations and potential biases of these models, enabling them to make informed decisions and take appropriate actions \cite{ref55,ref188}.

Moreover, interpretability and explainability are essential for fostering trust in LLMs. When users can understand the reasoning behind the models' outputs, they are more likely to rely on and accept the information or decisions provided by the system \cite{ref236,ref241}. This trust is particularly crucial in high-stakes domains, where the consequences of flawed or biased decisions can be severe \cite{ref240,ref28}.

In addition, interpretability and explainability can contribute to the ongoing development and refinement of LLMs. By providing insights into the models' internal workings and decision-making processes, researchers and developers can identify areas for improvement, address biases, and enhance the models' overall performance and reliability \cite{ref237,ref242}.

\subsection{Understanding LLM Architectures and Knowledge Composition}

The rapid advancements in large language models (LLMs) have prompted significant interest in understanding the mechanisms underlying their impressive performance. A key aspect of this understanding is elucidating the architectural composition of LLMs and how knowledge is encoded within their internal parameters.

Leveraging mechanistic interpretability techniques, researchers have made important strides in unraveling the inner workings of these complex models. One prominent approach is to employ activation maximization, a method that aims to synthesize inherently interpretable visual representations of the information encoded in individual neurons \cite{ref243}. This technique has been adapted to the natural language processing domain, with the introduction of ``feature textualization'' - a method that produces dense representations of neurons in the LLM word embedding space \cite{ref243}.

Through feature textualization, researchers have gained insights into the knowledge encoded in individual neurons within LLMs. However, a key finding from these studies is that individual neurons do not represent clearcut symbolic units of language, such as words, but rather encode a more nuanced and diffuse set of semantic concepts \cite{ref243}. This suggests that the architectural composition of knowledge within LLMs is not as straightforward as one might assume.

To further elucidate the knowledge composition within LLMs, researchers have explored the idea of ``competition of mechanisms'' - the interplay between multiple mechanisms within the model, and how one mechanism becomes dominant in the final prediction \cite{ref185}. By employing techniques such as logit inspection and attention modification, researchers have been able to uncover traces of these mechanisms and their competition across various model components, providing a more nuanced understanding of how LLMs leverage and reconcile different types of knowledge.

Building on these insights, researchers have also investigated the temporal evolution of the factual knowledge encoded in the latent representations of LLMs. By decoding the factual knowledge embedded in the latent space and representing its evolution across layers using a temporal knowledge graph, researchers have been able to expose different types of errors, ranging from representation-level to multi-hop reasoning errors \cite{ref244}.

This line of research highlights the importance of understanding the architectural composition of knowledge within LLMs, as it can shed light on the models' reasoning processes, their limitations, and potential avenues for improvement. By leveraging mechanistic interpretability techniques, researchers have been able to uncover the nuanced and multi-faceted nature of knowledge representation in these powerful models.

Furthermore, the insights gleaned from these studies have implications for enhancing the performance and reliability of LLMs. For instance, the understanding of how knowledge is encoded and retrieved within the models can inform strategies for model editing, pruning, and better alignment with human values \cite{ref237}. Additionally, the identification of specific mechanisms and their competition can enable more targeted interventions to improve the models' reasoning capabilities.

Looking ahead, the continued exploration of LLM architectures and knowledge composition remains a critical area of research. As these models become increasingly integral to a wide range of applications, understanding their inner workings is essential for ensuring their trustworthiness, safety, and alignment with human values. The insights gained from mechanistic interpretability studies will be pivotal in driving the responsible development and deployment of these transformative language models.

\subsection{Probing LLM Representations and Embeddings}

The rapid advancement of large language models (LLMs) has led to a growing interest in understanding how they encode and represent knowledge within their internal parameters and activations. Probing techniques have emerged as a powerful approach to shed light on the intricate mechanisms underlying LLM representations and embeddings \cite{ref237,ref245}.

One of the key insights gained from probing studies is the discovery that LLMs do not uniformly distribute knowledge across their layers and parameters. Instead, they tend to exhibit a hierarchical organization, where different layers and subsets of neurons specialize in encoding distinct types of linguistic and factual information \cite{ref246}.

Researchers have leveraged a variety of probing techniques to uncover the encoding of knowledge within LLM representations. One prominent approach is the use of linear probes, where a simple linear transformation is applied to the LLM's hidden representations to predict specific linguistic properties, such as part-of-speech, dependency relations, or semantic attributes \cite{ref247,ref248}. By analyzing the performance of these linear probes, researchers can gain insights into the extent to which the LLM's representations capture the targeted linguistic properties.

Another line of research has explored the use of more sophisticated probing methods, such as those based on information theory. These approaches aim to quantify the amount of information about a linguistic property that is encoded in the LLM's representations, going beyond simple classification accuracy \cite{ref249}. By leveraging measures like mutual information and entropy, researchers can gain a more nuanced understanding of how different layers and dimensions of the LLM's representations contribute to the encoding of specific linguistic phenomena.

In addition to probing the LLM's representations, researchers have also investigated the nature of the embeddings learned by these models. Embedding learning, a core component of representation learning, has been shown to be capable of modeling large-scale semantic knowledge graphs \cite{ref250}. By mapping the knowledge graph to a tensor representation, LLMs can learn to predict the entries of this tensor using their latent representations of entities and relations.

Interestingly, recent studies have found that the geometry of LLM representations can also provide insights into the encoding of truth and falsehood \cite{ref251}. Visualizations of LLM representations have revealed clear linear structures corresponding to the truthfulness of statements, suggesting that LLMs may linearly represent the truth or falsehood of factual statements.

Furthermore, probing techniques have been leveraged to understand the internal mechanisms by which LLMs extract and utilize factual knowledge. For instance, research has delved into the specific computational patterns and information flow within LLMs when answering queries that require accessing and reasoning about factual information \cite{ref252}. By tracing the propagation of information from relation to subject representations, these studies have uncovered a three-step internal mechanism for attribute extraction, where the subject representation is first enriched, followed by the integration of relational information and a targeted querying process.

The insights gained from probing LLM representations and embeddings have far-reaching implications for the development and deployment of these powerful models. By understanding how knowledge is encoded and organized within LLMs, researchers can work towards enhancing their interpretability, controllability, and alignment with human values and preferences. Moreover, the identification of specific knowledge structures and the ability to manipulate them through probing techniques open up new avenues for fine-tuning and editing LLMs to improve their performance, reliability, and safety in real-world applications.

\subsection{Analyzing LLM Training Dynamics and Emergent Behaviors}

The training dynamics of large language models (LLMs) are crucial for understanding how these powerful models acquire and represent knowledge, ultimately shedding light on the mechanisms behind their often surprising and emergent behaviors. Recent studies have delved into this topic, offering valuable insights through a mechanistic perspective.

One particularly intriguing phenomenon observed in LLMs is the process of ``grokking,'' where a model initially overfits the training data but then, after a prolonged period of training, undergoes a phase transition and begins to generalize remarkably well \cite{ref237,ref253,ref254,ref255,ref256,ref257,ref258,ref259}. Researchers have explored this phenomenon through the lens of mechanistic interpretability, reverse-engineering the algorithms and internal representations learned by the models.

One key insight from these studies is that grokking is not a sudden or discrete event, but rather a gradual process of amplifying structured mechanisms encoded in the model's weights \cite{ref260}. The authors of this work demonstrate that grokking arises from the gradual strengthening of specific circuits or sub-networks within the model, which are initially weak but then become dominant through the training process.

Similarly, the phenomenon of memorization in LLMs has been the subject of extensive research, with studies exploring the mechanisms behind this behavior \cite{ref261,ref262}. These studies have revealed that larger LLMs tend to memorize training data faster, but can also forget less throughout the training process, challenging the common assumption that larger models are more prone to overfitting.

The researchers found that models memorize certain parts of speech, such as nouns and numbers, more readily, as these entities can serve as unique identifiers for individual training examples. They also observed that the dynamics of memorization vary across different parts of speech, with nouns and numbers being memorized first, followed by other parts of speech.

These insights into the training dynamics and emergent behaviors of LLMs have important implications for understanding the inner workings of these models and developing more interpretable and reliable systems. By reverse-engineering the specific mechanisms and circuits responsible for grokking, memorization, and other emergent phenomena, researchers can gain a deeper understanding of how LLMs acquire and represent knowledge, and how this knowledge can be leveraged or modified to improve their performance, safety, and alignment with human values.

Moreover, the mechanistic perspective on LLM training dynamics can inform the design of new architectures, training regimes, and evaluation methods that are better suited to capture and harness the models' emergent capabilities. For example, the insights into the gradual amplification of structured representations during grokking could inspire the development of more targeted training strategies that actively encourage the formation of these beneficial circuits, or the design of evaluation frameworks that can better detect and quantify the emergence of these capabilities.

Overall, the research on analyzing the training dynamics and emergent behaviors of LLMs through a mechanistic lens represents a crucial step towards a deeper, more comprehensive understanding of these powerful models. By uncovering the internal mechanisms and processes that underlie their often surprising and complex behaviors, researchers can work towards developing the next generation of interpretable, controllable, and trustworthy large language models that can be safely and effectively deployed in a wide range of real-world applications.

\subsection{Enhancing LLM Performance and Alignment through Explainability}

The insights gained from the interpretability and explainability of large language models (LLMs) can be leveraged to enhance their performance and alignment with human values in various ways. Recent research has highlighted the critical importance of interpretability and explainability in LLMs as they become increasingly integral to diverse applications, raising concerns regarding transparency and ethical use \cite{ref263}.

One significant way in which explainability can enhance LLM performance is through model editing. Researchers have provided an in-depth exploration of the problems, methods, and opportunities related to model editing for LLMs, emphasizing that the methodology for maintaining the relevancy and rectifying errors of LLMs remains elusive \cite{ref264}. The surge in techniques for editing LLMs aims to efficiently alter the behavior of LLMs within a specific domain without negatively impacting performance across other inputs.

By understanding the internal mechanisms and knowledge composition of LLMs through various interpretability techniques, such as mechanistic interpretability and probing of representations and embeddings \cite{ref237}, researchers and practitioners can identify the precise location of knowledge in the model's neurons. This knowledge can then be leveraged to selectively edit the model's behavior, either by modifying the existing knowledge or introducing new knowledge, without disrupting the model's overall performance.

The application of model editing in the medical domain has been demonstrated, where researchers proposed a novel Layer-wise Scalable Adapter (MedLaSA) strategy to efficiently edit the factual medical knowledge and the explanatory ability of LLMs \cite{ref265}. By identifying the precise location of medical knowledge in the model's neurons and introducing scalable adapters, the authors were able to edit the model's knowledge without significantly affecting its performance on unrelated tasks.

Furthermore, the insights gained from the interpretability of LLMs can also be used to improve the efficiency of these models through techniques like pruning. Researchers have provided a comprehensive analysis of the impact of common compression techniques, such as pruning and quantization, on the parametric knowledge of LLMs \cite{ref266}. By understanding the location and structure of knowledge within the model, researchers can develop more targeted and effective pruning strategies that maintain the model's essential capabilities while significantly reducing its computational and memory footprint.

In addition to enhancing performance, the insights from LLM interpretability can also be leveraged to better align these models with human values and preferences. Researchers have demonstrated the use of post-hoc explanation methods to generate automated natural language rationales that provide corrective signals to LLMs, helping to improve their reliability, fluency, and citation accuracy \cite{ref267}. Similarly, a novel learning paradigm, ``Second Thought,'' has been proposed, which models the chain-of-edits between value-unaligned and value-aligned text, enabling LMs to re-align with human values \cite{ref268}.

The interpretability and explainability of LLMs can also play a crucial role in enhancing their transparency and trustworthiness, which is essential for their safe and ethical deployment in high-stakes applications. Researchers have proposed a ``metacognitive'' approach, dubbed CLEAR, to equip LLMs with capabilities for self-aware error identification and correction \cite{ref269}. By constructing concept-specific sparse subnetworks that illuminate transparent decision pathways, CLEAR provides a novel interface for model intervention after deployment, enabling the LLMs to self-consciously identify potential mispredictions and efficiently self-correct their errors.

In summary, the insights gained from the interpretability and explainability of LLMs can be leveraged to enhance their performance through targeted model editing, improve their efficiency through pruning, and better align them with human values and preferences. These advancements are crucial for the safe and ethical deployment of LLMs in a wide range of applications, especially in high-stakes domains where transparency and trustworthiness are of utmost importance.

\subsection{Evaluating the Quality and Reliability of LLM Explanations}

Evaluating the quality and reliability of LLM-generated explanations is crucial to ensuring their trustworthiness and usefulness in real-world applications. Existing research has proposed various approaches to quantify the faithfulness, plausibility, and robustness of model explanations, leveraging both human-centric and automated evaluation methods.

One common approach to evaluating explanation quality is through human evaluation, where study participants assess the coherence, clarity, and faithfulness of the explanations \cite{ref270,ref271}. However, human evaluation can be subjective and time-consuming, especially when scaling to large-scale evaluations.

To address the limitations of human evaluation, researchers have developed automated metrics to quantify the quality of LLM explanations. These include Attribution-Similarity, which measures the similarity between the input-output attribution of the model and the natural language explanation \cite{ref272}, as well as NLE-Sufficiency and NLE-Comprehensiveness, which evaluate the completeness and comprehensiveness of the explanations \cite{ref272}.

Additionally, researchers have proposed methods to quantify the uncertainty in LLM-generated explanations, such as Verbalized Uncertainty \cite{ref273} and Probing Uncertainty \cite{ref273}, which can provide valuable insights into the reliability of the explanations.

Beyond faithfulness and uncertainty, the robustness of explanations is another crucial aspect to consider. Researchers have introduced metrics to assess the stability of explanations to small perturbations in the input \cite{ref274} and the consistency of explanations across related inputs \cite{ref275}.

To further enhance the evaluation of LLM explanations, researchers have explored the use of synthetic datasets and causal reasoning. Approaches like Logical Satisfiability of Counterfactuals for Faithful Explanations in NLI \cite{ref276} and Faithfulness-through-Counterfactuals \cite{ref276} leverage counterfactual reasoning to provide a more rigorous assessment of the faithfulness of explanations.

The integration of causal reasoning in the evaluation of LLM explanations is a promising direction, as it can help uncover the underlying causal relationships and mechanisms driving the model's predictions \cite{ref277}. Additionally, the development of collaborative evaluation frameworks that leverage the complementary strengths of human and LLM-based evaluators can enhance the reliability and efficiency of explanation assessments \cite{ref278}.

Overall, the research on evaluating the quality and reliability of LLM explanations has made significant strides, but there is still ample room for further advancements. As LLMs become increasingly ubiquitous in high-stakes applications, the need for robust and trustworthy evaluation methodologies will only grow more pressing. Continued efforts to develop innovative evaluation approaches, leverage causal reasoning, and foster collaborative human-LLM evaluation frameworks will be crucial to ensuring the transparency and accountability of these powerful language models.

\subsection{Interpretable Interaction and Collaboration between Humans and LLMs}

The interpretability of large language models (LLMs) is critical for fostering effective collaboration and interaction between humans and these powerful AI systems. As LLMs continue to demonstrate remarkable capabilities across a wide range of tasks, enabling users to better understand, control, and trust the model's behavior and outputs becomes increasingly important \cite{ref241,ref279}.

Interpretability in this context goes beyond simply explaining the model's internal mechanisms or decision-making processes. It involves creating a shared understanding between the human user and the LLM, where the user can effectively engage with the model, provide feedback, and shape its behavior to align with their goals and preferences \cite{ref280,ref281}.

One key aspect of facilitating interpretable interaction is the ability to tailor the explanations and outputs provided by the LLM to the user's specific needs and mental models \cite{ref282}. This requires understanding the user's background, expertise, and cognitive biases, and adapting the explanations accordingly \cite{ref279}. For example, a domain expert may require more detailed and technical explanations of an LLM's decision-making process, while a non-expert user may benefit more from high-level, intuitive explanations that focus on the model's overall reasoning and the factors that influence its outputs \cite{ref283,ref284}. By catering to the user's unique needs and preferences, the LLM can foster a sense of trust and collaboration, enabling the user to better understand the model's capabilities and limitations, and to effectively guide its behavior \cite{ref188}.

Another crucial aspect of interpretable interaction is the ability for users to actively engage with and control the LLM's behavior. This involves providing mechanisms for users to inspect the model's internal representations, to edit or override its outputs, and to iteratively refine its behavior through feedback and guidance \cite{ref285,ref286}. For instance, users should be able to explore and understand the features and concepts that the LLM has learned, and to provide feedback on the relevance and appropriateness of these representations \cite{ref287}. Additionally, users should be able to intervene and correct the model's outputs, either by directly editing the generated text or by providing guidance on the desired behavior \cite{ref288,ref269}.

By empowering users to actively shape and refine the LLM's behavior, the interaction becomes more collaborative and interpretable, fostering a shared understanding and a sense of control that can lead to greater trust and better alignment between the human user and the AI system \cite{ref289,ref290}.

Furthermore, the interpretability of the collaboration between humans and LLMs can be enhanced by leveraging the complementary strengths of both parties. For example, LLMs can be used to generate explanations and insights that augment human understanding, while humans can provide feedback and oversight to ensure the LLM's behavior remains aligned with their goals and values \cite{ref291}. This synergistic relationship between human and LLM can lead to more robust and trustworthy decision-making, where the LLM's capabilities are harnessed in a way that is transparent and accountable to the user \cite{ref292,ref293}.

However, achieving truly interpretable interaction and collaboration between humans and LLMs is not without its challenges. The sheer complexity and scale of LLMs can make it difficult to fully capture and explain their inner workings, even with advanced interpretability techniques \cite{ref237,ref294}. Additionally, the rapidly evolving nature of LLM capabilities and the potential for unintended behaviors, such as hallucinations or biases, can make it challenging to maintain a consistent and reliable collaborative relationship \cite{ref178}.

To address these challenges, ongoing research efforts are exploring new approaches to enhance the interpretability and trustworthiness of LLMs, such as developing more transparent model architectures, leveraging causal reasoning, and incorporating human feedback and oversight into the model's learning and deployment processes \cite{ref295,ref296}. By continuously improving the interpretability and collaborative capabilities of LLMs, researchers and practitioners can unlock the full potential of these powerful AI systems, enabling more effective and trustworthy human-AI partnerships that can tackle complex problems and drive positive societal impact.

\subsection{Open Challenges and Future Directions in LLM Interpretability}

Despite the significant progress in understanding and interpreting the inner workings of large language models (LLMs), several key challenges and limitations remain. As LLMs continue to grow in scale and complexity, addressing these open challenges will be crucial to realizing their full potential and ensuring their safe and responsible deployment.

One of the primary open challenges in LLM interpretability is the integration of causal reasoning. While LLMs have demonstrated impressive performance on various tasks, their ability to understand and reason about causal relationships remains limited \cite{ref277,ref297}. Traditional approaches to interpretability have often focused on association-based explanations, failing to capture the underlying causal mechanisms that drive model behavior. Integrating causal reasoning into the interpretability of LLMs would enable a deeper understanding of the models' decision-making processes and enhance their reliability and robustness \cite{ref298}.

Another key challenge lies in the development of multimodal explanations for LLMs. As LLMs become increasingly adept at processing and integrating diverse data modalities, such as text, images, and even audio, the need for interpretability that spans these modalities becomes increasingly critical \cite{ref299}. Current interpretability methods have primarily focused on text-based explanations, leaving a significant gap in our understanding of how LLMs reason about and combine information from multiple modalities \cite{ref300}. Advancing the field of multimodal explanations for LLMs would not only enhance their transparency but also enable more effective human-AI collaboration in complex, multimodal tasks.

The role of human feedback and oversight in LLM interpretability also presents a significant challenge. While LLMs have demonstrated remarkable capabilities, their behavior can often be unpredictable, and their outputs may not align with human values and preferences \cite{ref178}. Developing effective mechanisms for incorporating human feedback and oversight into the interpretability of LLMs is crucial to ensuring their alignment with human values and preferences \cite{ref301}. This could involve the integration of human-in-the-loop approaches, where users can actively provide feedback and guidance to the models, or the development of interpretable mechanisms that allow for the rapid and intuitive evaluation of model outputs by human experts \cite{ref241}.

Furthermore, the scalability and generalizability of interpretability methods for LLMs remain significant challenges. As LLMs grow in size and complexity, the computational and memory requirements for interpreting their behavior can become prohibitive. Developing scalable and efficient interpretability techniques that can handle the growing scale of LLMs is crucial to ensuring their widespread adoption and deployment \cite{ref242}. Additionally, ensuring that interpretability methods can generalize across different types of LLMs and applications remains an open challenge, as the unique characteristics and biases of each model may require tailored interpretability approaches \cite{ref263}.

Finally, the integration of LLM interpretability with other AI capabilities, such as causal reasoning, multimodal understanding, and human-AI collaboration, presents a significant opportunity for future research \cite{ref302,ref303}. By leveraging the synergies between these different areas, researchers can develop more comprehensive and effective approaches to LLM interpretability, ultimately enhancing the reliability, safety, and trustworthiness of these powerful models \cite{ref304,ref299}.

\section{Open-source Resources and Initiatives}

\subsection{Open-source LLM Evaluation Resources and Initiatives}

In the rapidly evolving landscape of large language models (LLMs), the evaluation of these powerful AI systems has become a critical focus for the research community. As LLMs continue to advance, the need for comprehensive and standardized evaluation frameworks has become increasingly pressing to understand their capabilities, limitations, and potential impacts.

One prominent open-source initiative that has emerged to address this challenge is the GPT-Fathom benchmark \cite{ref83}. GPT-Fathom provides a systematic and reproducible approach to evaluating over 10 leading LLMs, including OpenAI's legacy models, across a diverse set of more than 20 curated benchmarks spanning 7 capability categories. By aligning the evaluation settings and prompts, GPT-Fathom offers valuable insights into the evolutionary trajectory of these models, addressing the issues of inconsistent reporting and potential cherry-picking that have plagued existing LLM leaderboards.

Complementing these broader evaluation efforts, the research community has also developed specialized open-source resources targeted towards specific domains and applications. For instance, the EquityMedQA dataset \cite{ref81} aims to surface biases and equity-related harms in LLM-generated medical content, while the MedAlpaca dataset and training framework \cite{ref85} provides a valuable resource for advancing the development and deployment of medical-domain LLMs.

Beyond standalone initiatives, the community has also seen the emergence of novel evaluation approaches, such as the PiCO framework \cite{ref84}, which explores an unsupervised peer-review mechanism for assessing the ability hierarchy among LLMs. These open-source efforts collectively demonstrate a growing emphasis on transparency, reproducibility, and responsible development in the LLM research ecosystem.

The comprehensive evaluation of LLMs has also been the focus of broader survey efforts, such as the work presented in ``Evaluating Large Language Models: A Comprehensive Survey'' \cite{ref11}. This survey provides a panoramic perspective on the current state of LLM evaluation methodologies, benchmarks, and specialized domain-specific assessments, underscoring the importance of developing comprehensive evaluation platforms.

The open-source nature of these initiatives, characterized by the availability of benchmarks, datasets, and evaluation frameworks, is a significant advantage. It enables the broader research community to collaborate, replicate, and build upon these resources, driving progress in the field and addressing the pressing challenges associated with the widespread adoption of LLMs. As the capabilities of these models continue to evolve, the role of open-source evaluation resources in guiding their development and deployment will only become more crucial.

\subsection{GPT-Fathom: Benchmarking LLMs Towards GPT-4 and Beyond}

The rapid advancement of large language models (LLMs) has driven the need for a comprehensive evaluation suite to assess their capabilities and limitations. Existing LLM leaderboards often reference scores reported in other papers without consistent settings and prompts, which may inadvertently encourage cherry-picking favored settings and prompts for better results \cite{ref83}. To address this issue, researchers have introduced GPT-Fathom, an open-source and reproducible LLM evaluation suite built on top of OpenAI Evals.

GPT-Fathom systematically evaluates over 10 leading LLMs, including OpenAI's legacy models, on more than 20 curated benchmarks across 7 capability categories, all under aligned settings \cite{ref83}. This comprehensive approach provides valuable insights into the evolutionary path from GPT-3 to GPT-4, addressing questions that the community is eager to understand.

One of the key aspects of GPT-Fathom is its ability to shed light on the technical details of LLM advancements. For instance, the evaluation suite explores whether adding code data improves an LLM's reasoning capability, which aspects of LLM capability can be improved by supervised fine-tuning (SFT) and reinforcement learning from human feedback (RLHF), and how much the ``alignment tax'' impacts LLM performance \cite{ref83}. By analyzing the retrospective performance of OpenAI's earlier models, GPT-Fathom offers valuable insights into the technical evolution of these models, helping to improve the transparency of advanced LLMs.

Moreover, GPT-Fathom's systematic approach to evaluating LLMs across diverse benchmarks and capability categories provides a more holistic understanding of their strengths and limitations. The evaluation suite covers a wide range of tasks, including language understanding, generation, reasoning, and task-specific applications, allowing researchers and practitioners to gain a comprehensive view of LLM performance.

One of the key findings from the GPT-Fathom evaluation is the identification of specific areas where LLMs excel and where they still struggle. For example, the study may reveal that a particular LLM performs exceptionally well on language generation tasks but falters on complex reasoning problems. This granular understanding of LLM capabilities can inform the development of more targeted training strategies and prompt engineering techniques to address their weaknesses and further enhance their performance.

Furthermore, the open-source nature of GPT-Fathom is a significant advantage, as it enables the broader research community to collaborate, replicate, and build upon the evaluation framework. By making the code, datasets, and evaluation scripts publicly available, GPT-Fathom promotes transparency and encourages the development of more robust and reliable LLM evaluation methodologies.

The systematic benchmarking approach of GPT-Fathom also extends beyond the current state-of-the-art models. By including OpenAI's legacy models in the evaluation, the suite provides insights into the evolutionary trajectory of these models, shedding light on the advancements and the challenges encountered along the way \cite{ref83}. This retrospective analysis not only offers valuable lessons for the future development of LLMs but also helps the research community better understand the technical foundations upon which the latest models are built.

In addition to the benchmarking capabilities, GPT-Fathom's open-source nature allows researchers and practitioners to customize the evaluation suite to meet their specific needs. The modular design of the framework enables the integration of new benchmarks, the incorporation of additional evaluation metrics, and the exploration of alternative evaluation strategies tailored to various domains and use cases.

The comprehensive evaluation approach and the open-source nature of GPT-Fathom make it a valuable resource for the LLM research community. By providing a standardized and reproducible evaluation framework, GPT-Fathom helps to address the challenges of inconsistent reporting and cherry-picking in existing LLM leaderboards. Moreover, the insights generated by GPT-Fathom can inform the ongoing development of LLMs, guiding researchers and engineers towards more robust and reliable models that can better serve the needs of diverse applications and users.

\subsection{Comprehensive Evaluation Frameworks for LLMs}

The rapid development of large language models (LLMs) has led to the emergence of various open-source platforms aimed at facilitating the comprehensive evaluation of these powerful models. These platforms seek to address the multifaceted aspects of LLM assessment, including their capabilities, alignment with human values, and safety considerations.

One prominent example is the OpenEval platform \cite{ref305}, which introduces a comprehensive evaluation testbed for Chinese LLMs. OpenEval encompasses three key dimensions: capability assessment, alignment evaluation, and safety evaluation. The capability assessment component includes 12 benchmark datasets spanning various tasks, such as NLP challenges, disciplinary knowledge, commonsense reasoning, and mathematical reasoning. This multifaceted approach ensures a thorough examination of the LLMs' language understanding and generation abilities across diverse domains.

In the domain of alignment evaluation, OpenEval incorporates 7 datasets that assess the potential biases, offensiveness, and illegality in the outputs generated by Chinese LLMs. This aligns with the growing concern for the responsible development of LLMs, ensuring that they adhere to ethical principles and societal norms. The safety evaluation component of OpenEval focuses on assessing potential risks, such as power-seeking behaviors and self-awareness, in advanced LLMs. By including 6 dedicated datasets for safety evaluation, OpenEval aims to uncover and mitigate the anticipated risks associated with the proliferation of these models.

Another comprehensive evaluation framework is the GPT-Fathom suite \cite{ref83}, which systematically evaluates over 10 leading LLMs, including OpenAI's legacy models, across 20+ curated benchmarks spanning 7 capability categories. The authors of GPT-Fathom emphasize the importance of providing a consistent and reproducible evaluation environment, addressing the limitations of existing leaderboards that often rely on inconsistent settings and prompts. By conducting a retrospective study on OpenAI's earlier models, GPT-Fathom offers valuable insights into the evolutionary path from GPT-3 to GPT-4, shedding light on the technical details and factors that contribute to the progressive improvements in LLM capabilities.

The FreeEval framework \cite{ref76} aims to address the challenges of integrating diverse evaluation methodologies and ensuring the reliability and efficiency of LLM assessments. FreeEval provides a modular and scalable platform that simplifies the integration of various evaluation approaches, including dynamic evaluation techniques that require sophisticated LLM interactions. Additionally, FreeEval incorporates meta-evaluation techniques, such as human evaluation and data contamination detection, to enhance the fairness and reliability of the evaluation outcomes. The framework's high-performance infrastructure, including distributed computation and caching strategies, enables extensive evaluations across multi-node, multi-GPU clusters for both open-source and proprietary LLMs.

The Prometheus platform \cite{ref306} introduces an open-source LLM evaluator that aims to match the capabilities of proprietary models, such as GPT-4, in assessing long-form responses. Prometheus is trained on the Feedback Collection dataset, which consists of fine-grained score rubrics, instructions, and responses generated by GPT-4. This approach allows Prometheus to evaluate any given long-form text based on customized score rubrics provided by the user, achieving performance on par with GPT-4 while being fully open-source and accessible.

The PiCO framework \cite{ref84} explores a novel unsupervised approach to LLM evaluation, utilizing a peer-review mechanism where LLMs evaluate each other's responses. PiCO formalizes this as a constrained optimization problem, aiming to maximize the consistency of each LLM's capabilities and scores. This approach aims to determine the ability hierarchy among LLMs without relying on human annotations or closed-environment benchmarks.

These comprehensive evaluation frameworks, along with others like AgentSims \cite{ref307}, AgentBoard \cite{ref208}, and Chatbot Arena \cite{ref308}, demonstrate the concerted efforts to develop robust and versatile platforms for assessing LLMs across various dimensions. By providing open-source resources, standardized benchmarks, and advanced evaluation methodologies, these initiatives pave the way for more transparent, reproducible, and meaningful assessments of LLMs, ultimately guiding their responsible development and deployment.

\subsection{Toolkits for Surfacing Biases and Equity Harms in LLMs}

The growing integration of large language models (LLMs) in various applications has led to concerns regarding their potential to perpetuate societal biases and equity-related harms. To address this critical issue, the research community has developed several open-source resources and methodologies aimed at surfacing and evaluating these biases.

One prominent example is the EquityMedQA dataset \cite{ref81}, which was curated to identify biases in LLM-generated answers to medical questions. The dataset was designed through an iterative participatory approach, involving diverse raters to assess potential biases across multiple demographic factors. This comprehensive approach helps to surface biases that may have been missed by narrower evaluation methods.

Beyond EquityMedQA, the research community has developed other open-source toolkits and frameworks to surface biases and equity-related harms in LLMs. The ROBBIE framework \cite{ref153} provides a systematic approach to benchmarking bias and toxicity across multiple demographic axes and LLM families. Similarly, the NBIAS framework \cite{ref309} offers a comprehensive and robust approach to detecting biases in textual data, including that generated by LLMs.

The Eagle dataset \cite{ref310} aims to better reflect the types of biases and harms that may arise in real-world LLM applications, by using prompts that users would actually provide in everyday contexts, rather than relying solely on artificially created datasets.

The research community has also developed open-source frameworks to address the challenge of evaluating the fairness and equity of LLMs in specific domains. For instance, the FairMonitor framework \cite{ref311} proposes a four-stage approach to evaluating stereotypes and biases in LLM-generated content, focusing on the education sector.

Furthermore, the Bias in Language Models: Beyond Trick Tests and Toward RUTEd Evaluation paper \cite{ref312} highlights the limitations of decontextualized ``trick tests'' for evaluating bias in LLMs, and proposes a more grounded approach, called RUTEd (Realistic Use and Tangible Effects) evaluation, which assesses biases in the context of realistic, long-form content generation tasks.

These open-source resources and methodologies demonstrate the research community's commitment to addressing the critical issue of biases and equity-related harms in LLMs. By developing comprehensive evaluation frameworks, curating diverse datasets, and fostering collaboration between researchers, practitioners, and domain experts, these initiatives pave the way for the responsible development and deployment of LLMs that promote fairness and equity across various applications.

\subsection{Peer-review-based LLM Evaluation Approaches}

Peer review has long been a cornerstone of academic quality control, and the same principles are now being applied to the evaluation of large language models (LLMs) in an open-source setting. Initiatives like PiCO \cite{ref84} have leveraged the concept of peer review to automatically evaluate and rank LLMs in an unsupervised manner.

The key idea behind PiCO is to create an environment where multiple LLMs can evaluate each other's responses, with the underlying assumption that higher-performing LLMs will be able to more accurately assess the capabilities of lower-performing ones. By assigning each LLM a learnable capability parameter, PiCO aims to determine the ability hierarchy among the models through a constrained optimization process that maximizes the consistency of each LLM's capabilities and scores.

This approach is particularly compelling as it addresses some of the inherent limitations of existing LLM evaluation methods. Traditional approaches often rely on human annotations or the use of a single ``strongest'' LLM as the evaluator, which can introduce biases and inconsistencies. In contrast, the peer review-based approach of PiCO allows for a more objective and comprehensive assessment of LLM performance, without the need for ground truth data or human oversight.

Another peer review-based initiative, PRE \cite{ref80}, takes a similar approach but with a slightly different implementation. PRE constructs a small qualification exam to select ``reviewer'' LLMs from a pool of powerful models, which are then used to rate or compare the ``submissions'' of different candidate LLMs. The final ranking of the evaluatee LLMs is generated based on the results provided by all the reviewer models.

The key advantage of these peer review-based approaches is their ability to leverage the collective knowledge and capabilities of multiple LLMs, rather than relying on a single model or human evaluators. By creating a self-contained ecosystem where LLMs assess each other, these initiatives aim to overcome the biases and limitations inherent in traditional evaluation methods.

Moreover, the unsupervised nature of these peer review-based approaches is particularly appealing, as it eliminates the need for costly and time-consuming human annotations or the development of specialized evaluation datasets. This makes the evaluation process more scalable and accessible, allowing for the rapid and efficient assessment of a wide range of LLMs.

However, it is important to note that the success of these peer review-based approaches relies on several critical assumptions. Firstly, it is crucial that the LLMs participating in the evaluation process have a sufficiently diverse range of capabilities, ensuring that higher-performing models can accurately assess the weaknesses of their lower-performing counterparts. Secondly, the models must exhibit a consistent and reliable ability to evaluate each other's responses, without being susceptible to self-enhancement or positional biases.

To address these challenges, recent research has explored techniques like the incorporation of peer discussion \cite{ref313} and the use of ordered weighted averages \cite{ref314} to improve the robustness and reliability of peer review-based LLM evaluations.

Furthermore, the advent of large-scale open-source LLM ecosystems, such as the H2O Open Ecosystem \cite{ref315}, has the potential to significantly enhance the viability and impact of peer review-based evaluation approaches. By providing a curated collection of high-quality, open-source LLMs, these initiatives can create a fertile ground for the development and testing of novel peer review-based evaluation frameworks.

Looking ahead, the integration of peer review-based evaluation methods with other open-source initiatives, such as the GPT-Fathom benchmark \cite{ref83} and the Chatbot Arena platform \cite{ref308}, could further strengthen the robustness and transparency of LLM evaluation. By combining the strengths of peer review, comprehensive benchmarking, and human-in-the-loop assessments, the research community can work towards a more holistic and reliable approach to evaluating the capabilities and limitations of large language models.

\subsection{Open-source LLM Ecosystems for Healthcare}

The rapid advancements in large language models (LLMs) have sparked a growing interest in their potential applications within the healthcare domain. However, the integration of these powerful models into clinical settings is not without its challenges, particularly concerning data privacy, computational constraints, and the need for domain-specific knowledge. To address these issues, the research community has introduced several open-source datasets, models, and frameworks that aim to facilitate the development and evaluation of medical LLMs.

One such initiative is the Hippocrates framework, which is specifically designed to advance the use of LLMs in healthcare \cite{ref86}. Hippocrates offers unrestricted access to its training datasets, codebase, checkpoints, and evaluation protocols, enabling the research community to build upon, refine, and rigorously evaluate medical LLMs within a transparent ecosystem. The framework includes the Hippo family of 7B models, which are tailored for the medical domain and fine-tuned from Mistral and LLaMA2 through continual pre-training, instruction tuning, and reinforcement learning from human and AI feedback. These Hippo models have demonstrated superior performance compared to existing open-source medical LLMs, even surpassing models with 70B parameters \cite{ref86}.

Another initiative, MedAlpaca, provides an open-source collection of medical conversational AI models and training data \cite{ref85}. This dataset, consisting of over 160,000 entries, is specifically crafted to fine-tune LLMs for effective medical applications. The authors investigate the impact of fine-tuning these datasets on publicly accessible pre-trained LLMs and compare the performance of pre-trained-only models against the fine-tuned models in medical certification exams.

The MedAgents framework, on the other hand, leverages LLM-based agents in a role-playing setting to enhance their proficiency and reasoning capabilities in the medical domain \cite{ref316}. This training-free framework engages multiple agents in a collaborative multi-round discussion, allowing them to mine and harness the medical expertise within LLMs and extend their reasoning abilities. Experimental results on various medical datasets demonstrate the effectiveness of this approach in improving LLM performance on zero-shot medical reasoning tasks.

Building upon the advancements in medical LLMs, the Clinical Camel model, fine-tuned from LLaMA-2 using QLoRA, achieves state-of-the-art performance across medical benchmarks among openly available medical LLMs \cite{ref317}. Leveraging efficient single-GPU training, Clinical Camel surpasses GPT-3.5 in five-shot evaluations on various medical benchmarks, including the USMLE Sample Exam, PubMedQA, MedQA, and MedMCQA.

Another notable open-source initiative is the PMC-LLaMA, which systematically investigates the process of adapting a general-purpose foundation language model towards the medical domain \cite{ref318}. This model is fine-tuned using a comprehensive dataset of 202M tokens, encompassing medical question-answering, reasoning rationales, and conversational dialogues. PMC-LLaMA exhibits superior performance across multiple public medical question-answering benchmarks, even surpassing ChatGPT.

The MEDITRON-70B, a suite of open-source LLMs with 7B and 70B parameters, extends pretraining on a curated medical corpus, including PubMed articles, abstracts, and medical guidelines, to enhance their performance on healthcare-specific tasks \cite{ref319}. Evaluations on four major medical benchmarks show significant performance gains over several state-of-the-art baselines, both before and after task-specific finetuning.

The open-source EquityMedQA dataset and associated frameworks \cite{ref81} aim to surface biases and equity-related harms in LLM-generated content for medical questions. This initiative, grounded in a participatory approach and diverse rater groups, highlights the importance of using varied assessment methodologies to identify biases that may be missed by narrower evaluation approaches.

Furthermore, the open-source CLUE benchmark \cite{ref320} provides a standardized approach to evaluating and developing LLMs in healthcare, aligning future model development with the real-world needs of clinical application. This benchmark includes novel datasets derived from MIMIC IV discharge letters and existing clinical tasks, enabling a comprehensive comparison between biomedical and general-domain LLMs.

These open-source initiatives, along with the release of datasets, models, and evaluation frameworks, demonstrate the research community's commitment to advancing the development and responsible deployment of medical LLMs. By fostering transparency, collaboration, and rigorous evaluation, these open-source resources aim to unlock the full potential of LLMs in enhancing healthcare outcomes while addressing the unique challenges posed by the integration of these powerful models in clinical settings.

\subsection{LLMs as Research Tools}

The rapid advancements in large language models (LLMs) have unlocked new opportunities for researchers across various domains, positioning these powerful models as invaluable research assistants. LLMs, with their remarkable capabilities in natural language processing and generation, have become instrumental in accelerating scientific discoveries and enhancing the efficiency of research workflows.

One prominent example of leveraging LLMs as research tools is in the field of qualitative analysis. Researchers have started to utilize the language understanding and generation capabilities of LLMs to assist in tasks such as text summarization, sentiment analysis, and thematic coding \cite{ref6}. By integrating LLMs into their workflows, researchers can efficiently process and analyze large volumes of qualitative data, saving time and resources while gaining valuable insights.

However, the use of LLMs as research tools is not without its challenges. One of the primary concerns is the issue of prompt tuning, where the specific phrasing and structure of the input prompt can significantly influence the model's output \cite{ref6}. Researchers must carefully craft prompts that elicit the desired response from the LLM, requiring a deep understanding of the model's capabilities and limitations.

Another challenge is the potential for biases and subjectivity in LLM-generated outputs. LLMs are trained on large, diverse datasets, which can inadvertently incorporate biases present in the data. Researchers must be cautious when relying on LLM outputs, as they may not always align with human judgment or expertise \cite{ref6}. To address this issue, researchers have explored strategies such as prompt optimization techniques and the integration of human expertise to mitigate the impact of these biases.

Despite these challenges, the potential of LLMs as research tools continues to be explored and expanded. Researchers have begun to leverage LLMs to assist in various research domains, such as scientific literature analysis, hypothesis generation, and policy recommendations \cite{ref321}. By automating the process of sifting through vast amounts of literature, LLMs can free up researchers' time, allowing them to focus on more high-level tasks and theorization. Furthermore, LLMs have shown promise in aiding researchers in hypothesis generation and idea exploration, as well as supporting policymakers and researchers in making more informed decisions.

While the use of LLMs as research assistants offers numerous benefits, it is crucial to acknowledge the limitations and potential risks associated with their deployment. Researchers must be cognizant of the model's tendencies to hallucinate or generate plausible-sounding but inaccurate information, as well as the potential for the amplification of existing biases \cite{ref6}. Careful validation, cross-checking, and the integration of human expertise are essential in mitigating these risks.

As the research community continues to explore the use of LLMs as assistants, the development of comprehensive guidelines, best practices, and ethical frameworks will be crucial. This will help ensure that the integration of LLMs into research workflows is done in a responsible and accountable manner, maximizing the benefits while minimizing the risks.

\subsection{Automatic Evaluation of LLMs' Clinical Capabilities}

The rapid advancements in large language models (LLMs) have unlocked new opportunities for healthcare, from information synthesis to clinical decision support. These new LLMs are not just capable of modeling language, but can also act as intelligent ``agents'' that interact with stakeholders in open-ended conversations and even influence clinical decision-making. However, rather than relying on benchmarks that measure a model's ability to process clinical data or answer standardized test questions, LLM agents should be assessed for their performance on real-world clinical tasks.

To address this need, researchers have proposed comprehensive evaluation frameworks, which they call ``Artificial-intelligence Structured Clinical Examinations'' (``AI-SCI''), to assess the performance of LLM agents in realistic clinical settings. These new evaluation frameworks draw from comparable technologies where machines operate with varying degrees of self-governance, such as self-driving cars. High-fidelity simulations may also be used to evaluate interactions between users and LLMs within a clinical workflow, or to model the dynamic interactions of multiple LLMs \cite{ref212}.

One such framework, introduced in ``Towards Automatic Evaluation for LLMs' Clinical Capabilities,'' proposes an automatic evaluation paradigm tailored to assess the LLMs' capabilities in delivering clinical services, such as disease diagnosis and treatment. The evaluation paradigm contains three key elements: metric, data, and algorithm. Inspired by professional clinical practice pathways, the researchers formulate a LLM-specific clinical pathway (LCP) to define the clinical capabilities that a doctor agent should possess. They then introduce Standardized Patients (SPs) from medical education as the guideline for collecting medical data for evaluation, which can ensure the completeness of the evaluation procedure \cite{ref87}.

Another approach, introduced in ``Automatic Interactive Evaluation for Large Language Models with State Aware Patient Simulator,'' focuses on the dynamic, realistic assessment of LLMs through multi-turn doctor-patient simulations. The researchers present the Automated Interactive Evaluation (AIE) framework and the State-Aware Patient Simulator (SAPS), which provide a closer approximation to real clinical scenarios and allow for a detailed analysis of LLM behaviors in response to complex patient interactions. Unlike prior methods that rely on static medical knowledge assessments, AIE and SAPS offer a dynamic platform for evaluating LLMs, enabling a more comprehensive assessment of their performance in clinical practice \cite{ref322}.

The ``Emulating Human Cognitive Processes for Expert-Level Medical Question-Answering with Large Language Models'' study introduces BooksMed, a novel framework based on a Large Language Model (LLM) that uniquely emulates human cognitive processes to deliver evidence-based and reliable responses, utilizing the GRADE (Grading of Recommendations, Assessment, Development, and Evaluations) framework to effectively quantify evidence strength. To ensure that clinical decision-making is appropriately assessed, the researchers present ExpertMedQA, a multispecialty clinical benchmark comprised of open-ended, expert-level clinical questions, and validated by a diverse group of medical professionals. By demanding an in-depth understanding and critical appraisal of up-to-date clinical literature, ExpertMedQA rigorously evaluates LLM performance \cite{ref82}.

The ``An Automatic Evaluation Framework for Multi-turn Medical Consultations Capabilities of Large Language Models'' study introduces an automated evaluation framework that assesses the practical capabilities of LLMs as virtual doctors during multi-turn consultations. The researchers design consultation tasks that require LLMs to be aware of what they do not know, to inquire about missing medical information from patients, and to ultimately make diagnoses. To evaluate the performance of LLMs for these tasks, a benchmark is proposed by reformulating medical multiple-choice questions from the United States Medical Licensing Examinations (USMLE), and comprehensive evaluation metrics are developed and evaluated on three constructed test sets \cite{ref323}.

These open-source frameworks and evaluation approaches demonstrate the growing efforts to develop comprehensive and systematic methods for assessing the clinical capabilities of LLMs. By moving beyond traditional benchmarks and static assessments, these initiatives aim to provide a more realistic and reliable evaluation of LLMs in healthcare settings, enabling the responsible and effective integration of these powerful models into clinical practice.

\subsection{Open-source Benchmarks and Frameworks for Long-context LLMs}

The rapid advancement of large language models (LLMs) has led to a growing demand for comprehensive evaluation frameworks that can effectively assess their performance, particularly in the context of long-sequence processing. To address this need, various open-source initiatives have emerged, providing researchers and practitioners with the necessary tools and resources to evaluate the capabilities of long-context LLMs.

One such prominent open-source benchmark is L-Eval, introduced in the paper ``L-Eval: Instituting Standardized Evaluation for Long Context Language Models'' \cite{ref89}. L-Eval is designed to institute a more standardized evaluation approach for long-context language models (LCLMs), addressing two key aspects: dataset construction and evaluation metrics. The benchmark encompasses 20 subtasks, 508 long documents, and over 2,000 human-labeled query-response pairs, covering diverse question styles, domains, and input lengths ranging from 3,000 to 200,000 tokens. The authors emphasize the importance of evaluating LCLMs using appropriate metrics, as they find that popular n-gram matching metrics generally fail to correlate well with human judgment. Instead, they advocate for the use of length-instruction-enhanced (LIE) evaluation and the employment of LLM judges to provide a more accurate assessment of long-context understanding.

Building upon the insights from L-Eval and Ada-LEval, the authors of ``LV-Eval: A Balanced Long-Context Benchmark with 5 Length Levels Up to 256K'' \cite{ref324} propose LV-Eval, a comprehensive long-context benchmark with five length levels (16,000, 32,000, 64,000, 128,000, and 256,000 words) reaching up to 256,000 words. LV-Eval features two main tasks, single-hop question answering and multi-hop question answering, comprising 11 bilingual datasets. The design of LV-Eval incorporates techniques such as confusing facts insertion, keyword and phrase replacement, and keyword-recall-based metric design to mitigate potential knowledge leakage and ensure more objective evaluations. The results from evaluating 10 LLMs on LV-Eval reveal that commercial LLMs generally outperform open-source models within their claimed context length, but their overall performance is surpassed by open-source LLMs with longer context lengths. Additionally, the authors highlight the significant performance degradation of LLMs in the presence of confusing information and the importance of addressing issues related to knowledge leakage and inaccurate metrics.

In a similar vein, the authors of ``LooGLE: Can Long-Context Language Models Understand Long Contexts'' \cite{ref12} introduced LooGLE, a long-context generic language evaluation benchmark for LLMs' long-context understanding. LooGLE features relatively new documents post-2022, with over 24,000 tokens per document and 6,000 newly generated questions spanning diverse domains. The benchmark required human annotators to meticulously craft more than 1,100 high-quality question-answer pairs to meet the long-dependency requirements. The evaluation of eight state-of-the-art LLMs on LooGLE revealed that commercial models outperformed open-source models, and LLMs struggled with more intricate long-dependency tasks, despite their success in short-dependency tasks.

\subsection{Integrating External Knowledge into Open-source LLMs}

The open-source large language models (LLMs), such as LLaMA and Mistral, have made significant strides in advancing the field of natural language processing. However, despite their impressive capabilities, these models still face challenges in certain specialized domains that require more targeted knowledge. To address this limitation, open-source initiatives have emerged that aim to enhance the performance of these LLMs by incorporating external knowledge sources.

One such initiative is the UMLS-augmented LLMs, which leverage the Unified Medical Language System (UMLS) to boost the medical knowledge and reasoning capabilities of open-source LLMs \cite{ref210}. UMLS is a comprehensive biomedical ontology that integrates various medical vocabularies, terminologies, and classifications, making it a valuable resource for enhancing the domain-specific knowledge of LLMs.

The UMLS-augmented LLM approach involves pre-training the base LLM, such as LLaMA or Mistral, on the UMLS knowledge base, in addition to the standard pre-training corpus. This additional pre-training step allows the LLM to acquire a deeper understanding of medical concepts, terminology, and relationships, which can be crucial for tasks like medical question answering, disease diagnosis, and treatment recommendations.

The integration of UMLS knowledge has been shown to significantly improve the performance of open-source LLMs on various medical tasks \cite{ref210}. For example, when tested on the LiveQA dataset, which consists of medical questions, the UMLS-augmented LLaMA-2-13b model exhibited a notable improvement in the ROUGE Score and BERTScore metrics compared to the base LLaMA-2-13b model. Furthermore, the UMLS-augmented model also achieved higher scores on physician-evaluated dimensions, such as factuality, completeness, readability, and relevancy, demonstrating its enhanced capability in generating medically accurate and relevant responses.

Another open-source initiative that aims to boost the medical capabilities of LLMs is the Hippocrates framework \cite{ref86}. Hippocrates is designed to provide a comprehensive ecosystem for the development, evaluation, and deployment of medical LLMs, addressing the challenges faced by proprietary models in terms of accessibility, transparency, and reproducibility.

The Hippocrates framework includes a suite of open-source medical LLMs, such as the Hippo family of models, which are fine-tuned from the base LLaMA and Mistral models using a variety of techniques, including continual pre-training, instruction tuning, and reinforcement learning from human and AI feedback. These Hippo models have demonstrated significant performance improvements over existing open-source medical LLMs, and in some cases, have even surpassed the capabilities of larger proprietary models, such as GPT-4 and Med-PaLM.

By making the training datasets, codebase, checkpoints, and evaluation protocols openly available, Hippocrates aims to stimulate collaborative research and facilitate advancements in the field of medical LLMs. This open approach encourages the broader scientific community to build upon, refine, and rigorously evaluate the medical LLMs within a transparent ecosystem, ultimately driving the progress of healthcare-focused AI research and applications.

In addition to the UMLS-augmented LLMs and the Hippocrates framework, there are other open-source initiatives that focus on integrating external knowledge into LLMs. For instance, the Knowledge Graph-based Retrieval Augmented Generation (KG-RAG) framework \cite{ref325} leverages the SPOKE biomedical knowledge graph to enhance the performance of LLMs, such as LLaMA-2-13b, GPT-3.5-Turbo, and GPT-4, on various biomedical tasks, including drug repurposing, true/false questions, and multiple-choice questions.

The KG-RAG framework generates meaningful biomedical text by combining the LLM's intrinsic knowledge with the structured knowledge from the SPOKE knowledge graph. This approach has been shown to consistently improve the performance of LLMs across different prompt types, with the Llama-2 model exhibiting a remarkable 71\% boost on the challenging multiple-choice question dataset when using the KG-RAG framework.

Furthermore, the Parametric Knowledge Guiding (PKG) framework \cite{ref326} aims to enhance the performance of open-source LLMs on domain-specific tasks by providing them with access to relevant external knowledge, without altering the LLMs' parameters. By leveraging a knowledge-guiding module, PKG allows LLMs to utilize domain-specific information, such as factual, tabular, medical, and multimodal knowledge, to improve their performance on a range of knowledge-intensive tasks.

These open-source initiatives, along with others that focus on integrating external knowledge into LLMs, demonstrate the growing recognition of the importance of domain-specific knowledge for enhancing the capabilities of these powerful language models. By tapping into specialized knowledge sources and making them accessible to open-source LLMs, these efforts help bridge the gap between the general-purpose knowledge of LLMs and the more targeted requirements of real-world applications, particularly in specialized domains like healthcare and biomedicine.

As the field of open-source LLMs continues to evolve, the integration of external knowledge is likely to become an increasingly crucial aspect of their development and deployment, enabling these models to better serve the diverse needs of users and applications across various industries and domains.

\section{Challenges and Future Directions}

\subsection{Integrating Causal Reasoning in LLMs}

Causal reasoning is a fundamental aspect of human intelligence, and its integration into Large Language Models (LLMs) has become a pressing concern in the field of artificial intelligence. The ability to understand and reason about cause-and-effect relationships is crucial for tasks such as decision-making, counterfactual analysis, and explanatory power, which are increasingly important as LLMs are deployed in real-world applications \cite{ref277}.

The emergence of LLMs has opened up new opportunities for advancing causal reasoning in AI systems. These models have demonstrated impressive capabilities in language understanding and generation, which can be leveraged to enhance their causal reasoning abilities. For example, LLMs can be used to extract and encode causal knowledge from large corpora of text, which can then be used to inform causal discovery and counterfactual reasoning tasks \cite{ref327,ref298}.

However, the current state of causal reasoning in LLMs is limited, and there are significant challenges that need to be addressed. One key challenge is the ability of LLMs to perform causal discovery, which involves identifying the underlying causal structure from observational data \cite{ref328}. While LLMs have shown promise in this area, their performance is often constrained by the quality and quantity of the training data, as well as the specific prompts or tasks used to elicit causal reasoning \cite{ref329}.

Another important aspect of causal reasoning is counterfactual reasoning, which involves the ability to imagine hypothetical scenarios and reason about their consequences \cite{ref330}. LLMs have demonstrated some success in this domain, but their performance is often limited by their reliance on statistical patterns in the training data, rather than a deeper understanding of causal mechanisms \cite{ref331}.

To address these challenges, researchers have proposed various approaches for integrating causal reasoning into LLMs. One promising approach is the use of structured causal models, which can provide a more explicit representation of causal relationships and enable more robust reasoning \cite{ref332}. These models can be used to guide the training of LLMs, ensuring that they learn to reason about causal relationships in a more principled way \cite{ref333}.

Another approach is to leverage the complementary strengths of LLMs and existing causal inference methods. For example, LLMs can be used to extract and encode causal knowledge from text, which can then be combined with statistical causal discovery techniques to improve the overall performance of causal reasoning \cite{ref334}.

Furthermore, researchers have explored the use of causal proxy models, which can be trained to mimic the behavior of LLMs while also providing more explicit causal explanations \cite{ref335}. These models can be used to enhance the interpretability and reliability of LLM-based systems, particularly in high-stakes applications.

Despite these promising approaches, there are still significant challenges in integrating causal reasoning into LLMs. One key challenge is the need to develop more robust and scalable causal reasoning algorithms that can handle the complexity and scale of modern language models \cite{ref336}. Additionally, there is a need to address the limitations of LLMs in terms of their ability to handle logical reasoning and symbolic inference \cite{ref337,ref338}.

Looking to the future, the integration of causal reasoning into LLMs holds great promise for advancing the field of artificial intelligence. By combining the powerful language understanding and generation capabilities of LLMs with more principled causal reasoning, researchers can develop AI systems that are more robust, reliable, and explainable, with applications across a wide range of domains \cite{ref339,ref340}. However, realizing this potential will require continued research and innovation, as well as a deeper understanding of the strengths and limitations of both LLMs and causal reasoning approaches.

\subsection{Addressing Limitations in Multi-Agent Reasoning}

Existing multi-agent systems that leverage large language models (LLMs) often face significant limitations, including narrow optimization objectives, restricted resource constraints, and challenges in interpreting and debugging the system's behavior. These limitations hinder the development of truly autonomous and adaptive multi-agent systems that can effectively navigate complex, real-world scenarios.

One of the primary limitations of current multi-agent systems is their narrow optimization objectives. Most existing approaches focus on optimizing for a single, specific task or metric, such as task completion rate or reward maximization \cite{ref341}. While these objectives may be suitable for well-defined, constrained problems, they often fail to capture the nuances and complexities of real-world situations, where agents may need to balance multiple, potentially conflicting objectives simultaneously.

Moreover, multi-agent systems operate within restricted resource constraints, such as limited computational power, memory, or communication bandwidth. These constraints can significantly impact the system's performance and limit its ability to engage in sophisticated reasoning and decision-making processes. For example, when resources are scarce, agents may be forced to make suboptimal decisions or resort to heuristics, compromising the overall effectiveness of the system \cite{ref233}.

Another significant challenge in current multi-agent systems is the difficulty in interpreting and debugging the system's behavior. As the complexity of the system increases, with multiple agents interacting and communicating in dynamic environments, the causal relationships between actions, decisions, and outcomes become increasingly opaque. This lack of transparency can make it challenging to identify the root causes of system failures or to fine-tune the system's performance \cite{ref342}.

To address these limitations, we introduce the concept of ``reasoning capacity'' as a unifying criterion for the design and evaluation of multi-agent systems. Reasoning capacity refers to the ability of an agent (or a system of agents) to engage in logical, causal, and strategic reasoning to address complex problems \cite{ref342}. By focusing on reasoning capacity as a central design principle, we can develop multi-agent systems that are more adaptable, flexible, and effective in navigating real-world scenarios.

This approach involves several key elements, including holistic objective formulation, adaptive resource management, transparent and interpretable reasoning, and self-reflection and meta-cognition. By incorporating these elements, we can create multi-agent systems that possess the necessary cognitive capabilities to tackle complex, real-world problems effectively.

Furthermore, the integration of reasoning capacity into multi-agent systems can also have broader implications for the field of artificial general intelligence (AGI). As researchers strive to develop AGI systems that can match or exceed human-level intelligence, the ability to engage in flexible, context-aware reasoning will be a crucial component. By advancing the state of the art in multi-agent reasoning capacity, we can contribute to the ongoing efforts towards the realization of AGI \cite{ref343}.

\subsection{Enhancing Social Reasoning in LLMs}

Evaluating the Theory-of-Mind (ToM) reasoning capabilities of large language models (LLMs) has been a topic of increasing interest and debate in the field of artificial intelligence. While recent research has shed light on the potential emergence of ToM-like abilities in LLMs, the findings have been inconsistent and the validity of the evaluation approaches has been a subject of concern \cite{ref344,ref345}.

One of the key challenges in assessing the social reasoning capabilities of LLMs lies in the complex and multifaceted nature of ToM. ToM encompasses the ability to understand and reason about the mental states of others, including their beliefs, desires, intentions, and emotions. Evaluating this aspect of cognition in LLMs requires a comprehensive and systematic approach that goes beyond simple test scenarios or anecdotal evidence.

To address these challenges, researchers have proposed a novel framework for procedurally generating social reasoning benchmarks using causal templates \cite{ref346,ref347}. This approach aims to create a more robust and ecologically valid evaluation framework that captures the nuances of ToM reasoning.

The key idea behind this framework is to leverage causal templates to generate a diverse set of social reasoning scenarios that require LLMs to reason about the mental states of multiple agents, including their beliefs, desires, and intentions. By using causal templates, the researchers can systematically control the complexity and structure of the scenarios, ensuring that the evaluation tasks are both challenging and informative.

One of the advantages of this approach is that it can address the inconsistent results and validity concerns that have plagued previous ToM evaluations. By generating a large and diverse set of scenarios, researchers can assess the LLMs' ToM reasoning capabilities across a wide range of contexts, reducing the risk of overfit to specific test cases or data leakage.

Moreover, the use of causal templates allows for the creation of scenarios that are grounded in established psychological theories and empirical findings, enhancing the ecological validity of the evaluation. This is particularly important when assessing the social reasoning capabilities of LLMs, which are expected to interact with humans in complex, real-world settings.

In addition to the procedurally generated scenarios, the proposed framework also incorporates the use of multi-agent settings, where LLMs are required to reason about the mental states of multiple characters simultaneously. This added complexity further challenges the LLMs' ToM reasoning abilities, as they must integrate and reconcile the perspectives of multiple agents to arrive at a coherent understanding of the social dynamics.

One example of such a multi-agent scenario could be a negotiation setting, where the LLM is tasked with understanding the motivations, goals, and strategies of both negotiating parties \cite{ref348}. By successfully navigating these complex social interactions, the LLM would demonstrate a more sophisticated grasp of ToM reasoning.

Furthermore, the proposed framework also explores the integration of causal reasoning and counterfactual thinking into the social reasoning tasks. This addition is motivated by the idea that ToM is not solely about recognizing mental states, but also about understanding the causal relationships and counterfactual implications that underlie social interactions \cite{ref347}.

By incorporating causal and counterfactual reasoning, the framework aims to capture a more holistic and nuanced understanding of ToM, as humans often rely on these cognitive capabilities when interpreting and predicting the behavior of others.

The development of this novel framework for generating social reasoning benchmarks is a significant step towards a more comprehensive and reliable evaluation of LLMs' ToM capabilities. By using causal templates to create diverse and challenging scenarios, researchers can gain a deeper understanding of the strengths and limitations of LLMs in this critical aspect of social intelligence.

Moreover, the framework's emphasis on multi-agent settings and the integration of causal and counterfactual reasoning can lead to the development of more sophisticated and versatile social reasoning models, better equipped to navigate the complexities of human social interactions.

As the field of artificial intelligence continues to advance, the need for LLMs with robust social reasoning capabilities becomes increasingly important. The proposed framework for generating social reasoning benchmarks represents a promising approach to address this challenge, paving the way for the development of LLMs that can truly understand and navigate the nuances of human social cognition.

\subsection{Advancing Multimodal Evaluation}

The rapid advancements in artificial intelligence, particularly the emergence of large language models (LLMs) \cite{ref61}, have paved the way for the development of large multimodal models (LMMs) capable of processing complex tasks involving joint reasoning over text and visual content. However, the evaluation of these LMMs remains a crucial challenge that needs to be addressed to ensure their reliable and responsible deployment.

Traditionally, the evaluation of LLMs has primarily focused on text-based tasks, neglecting the crucial aspect of multimodal reasoning. As LMMs gain prominence in diverse applications, such as navigating maps in public places \cite{ref349}, it becomes imperative to move beyond text-only assessments and integrate multimodal data and tasks in the evaluation process.

One key aspect that must be considered in the evaluation of LMMs is the importance of causal-aware reasoning. Causal inference plays a vital role in understanding the underlying mechanisms driving the relationships between textual and visual elements, which is crucial for tasks that demand a deeper understanding of the context-sensitive interactions between these modalities \cite{ref277,ref297}.

Recent studies have highlighted the potential synergy between LLMs and causal inference frameworks \cite{ref333,ref328,ref350}. By incorporating causal reasoning capabilities, LMMs can better capture the complex cause-effect relationships within multimodal data, leading to more robust and reliable performance.

The integration of causal inference into the evaluation of LMMs can take various forms. For instance, researchers can develop benchmark datasets that explicitly test the models' ability to reason about causal relationships between textual and visual elements \cite{ref330,ref351}. These datasets could include tasks that require the models to predict the outcomes of interventions or to identify confounding factors that may influence the observed relationships.

Moreover, the evaluation of LMMs can be enhanced by incorporating techniques from the field of causal inference, such as counterfactual reasoning and structural causal models \cite{ref352,ref327,ref329}. These approaches can provide insights into the models' understanding of causal mechanisms and their ability to generalize beyond the training data.

In addition to causal-aware reasoning, the evaluation of LMMs should also consider the models' ability to handle multimodal tasks that involve complex interactions between text and visual information. This could include tasks such as visual question answering, image captioning, and multimodal reasoning \cite{ref353,ref354}.

Furthermore, the evaluation of LMMs should take into account the models' performance not only on individual tasks but also on their ability to collaborate and integrate information across multiple modalities \cite{ref66,ref349}. This would require the development of comprehensive benchmarks and evaluation frameworks that assess the models' capacity to reason holistically and provide coherent and meaningful responses to complex, multi-faceted queries.

In conclusion, the evaluation of LMMs must evolve beyond text-only assessments to incorporate multimodal data and tasks, with a strong emphasis on causal-aware reasoning. By leveraging the synergies between LLMs and causal inference frameworks, researchers and practitioners can develop more robust and reliable LMMs that can navigate the complexities of the real world with greater accuracy and interpretability. This holistic approach to the evaluation of LMMs will be crucial in ensuring their responsible and effective deployment in a wide range of applications.

\subsection{Collaborative Evaluation Approaches}

Existing approaches to evaluating the quality of text generated by large language models (LLMs) often face significant limitations. Human evaluation, widely considered the gold standard, can be costly, time-consuming, and prone to inconsistencies \cite{ref66}. While the use of LLMs as alternative evaluators has shown promise, these single-agent-based approaches suffer from inherent biases and instabilities \cite{ref131,ref129}.

To address these limitations, researchers have explored the potential of leveraging the complementary strengths of humans and LLMs through collaborative evaluation approaches. The Collaborative Evaluation (CoEval) pipeline \cite{ref66} represents a notable advancement in this direction. CoEval combines the human ability to adapt evaluation criteria to diverse task requirements with the scalability and cost-effectiveness of LLM-based assessments. The framework involves a two-stage process, where an LLM first provides an initial evaluation, which is then scrutinized and revised by human evaluators.

The key advantage of the CoEval approach lies in its ability to harness the synergies between human and LLM-based evaluators. LLMs can efficiently assess lengthy texts, saving significant time and reducing human evaluation outliers. At the same time, human scrutiny plays a crucial role in revising around 20\% of LLM evaluation scores, ensuring the ultimate reliability of the assessment. This collaborative approach helps address the shortcomings of both human and LLM-based evaluators when working in isolation.

However, the CoEval pipeline is not without its own limitations. The framework still relies on the availability of human evaluators, which can be a scarce and expensive resource, particularly for large-scale evaluations. Additionally, the specific division of labor between humans and LLMs, as well as the optimal design of the collaboration, remain open questions that warrant further investigation.

Building on the insights gained from CoEval, researchers have explored alternative collaborative evaluation frameworks that further leverage the diverse capabilities of LLMs. The MATEval approach \cite{ref131} introduces a multi-agent referee team, where a group of LLMs engage in autonomous discussions and debates to collectively evaluate the quality of generated responses. This multi-agent approach enables the models to synergize their distinct capabilities and expertise, leading to more reliable and comprehensive assessments.

Similarly, the ChatEval framework \cite{ref132} employs a multi-agent debate approach, where a team of LLMs engage in discussions to evaluate open-ended responses. By harnessing the diverse perspectives and expertise of the agent team, ChatEval transcends mere textual scoring, offering a human-mimicking evaluation process for reliable assessments.

These multi-agent collaborative evaluation approaches hold promise in addressing the limitations of single-agent LLM-based evaluators. By leveraging the complementary strengths of multiple LLMs, they can enhance the efficiency, effectiveness, and reliability of the evaluation process. Additionally, the multi-agent debate format can provide valuable insights into the decision-making process of the LLMs, allowing for a more nuanced understanding of their evaluation capabilities.

Looking ahead, the field of collaborative evaluation for LLMs is ripe for further exploration and innovation. One potential direction is the development of more advanced coordination mechanisms among the LLM agents, such as the incorporation of self-reflection, feedback loops, and meta-learning strategies. This could enable the agent team to collectively refine their evaluation criteria and decision-making processes, leading to even more robust and aligned assessments.

Another promising avenue is the exploration of human-LLM collaborative evaluation frameworks that go beyond the simple division of labor. Researchers could investigate approaches where humans and LLMs engage in interactive dialogues, with both parties leveraging their unique strengths to jointly arrive at more comprehensive and insightful evaluations. This could involve humans providing high-level guidance and feedback to the LLMs, while the LLMs offer real-time analysis and suggestions to enhance the evaluation process.

Overall, the pursuit of collaborative evaluation approaches for LLMs represents a vital step in addressing the limitations of existing evaluation methods. By harnessing the complementary capabilities of humans and LLMs, researchers can develop more reliable, efficient, and comprehensive evaluation frameworks that keep pace with the rapid advancements in large language models. Continued progress in this direction will be crucial for ensuring the trustworthiness, transparency, and responsible development of these powerful AI systems.

\subsection{Addressing the Socio-Technical Gap}

The rapid advancement of artificial intelligence (AI) and the growing adoption of large language models (LLMs) have brought about a significant transformation in various domains, from healthcare to education and beyond. However, the evaluation of these powerful AI systems has often been confined to narrow technical capabilities, overlooking the crucial socio-technical considerations that are essential for their effective and responsible deployment \cite{ref64,ref355}.

The socio-technical gap refers to the disconnect between the technical capabilities of AI systems and the complex social, cultural, and environmental contexts in which they are expected to operate \cite{ref356}. This gap has become increasingly evident as the impacts of AI extend beyond the confines of individual tasks or applications, affecting entire socio-technical systems and the well-being of diverse stakeholders \cite{ref357}. To address this gap, evaluation methods must move beyond the traditional focus on technical metrics and incorporate a more holistic understanding of the societal implications and constraints within which these AI systems function \cite{ref358}.

This shift in perspective requires a fundamental rethinking of the evaluation process, one that prioritizes the alignment of AI systems with societal values, ethical principles, and the diverse needs of all affected stakeholders \cite{ref359,ref360}. One key aspect of this shift is the acknowledgment of the trade-offs that often arise between the technical optimization of AI systems and their social impact \cite{ref361}. Evaluation methods must be designed to navigate these complex trade-offs, finding the right balance between technical efficiency and societal considerations \cite{ref355}.

Moreover, the socio-technical gap highlights the need for diverse evaluation approaches that capture the multifaceted nature of AI systems \cite{ref362}. Traditional benchmarks and standardized metrics, while valuable, may fail to capture the nuanced and context-dependent impacts of AI on various social and environmental domains \cite{ref363}. Embracing a range of evaluation methods, including qualitative assessments, user studies, and participatory approaches, can provide a more comprehensive understanding of the societal implications of AI \cite{ref364}.

Addressing the socio-technical gap also requires acknowledging the pragmatic constraints and trade-offs inherent in the evaluation process itself \cite{ref365}. Conducting comprehensive socio-technical evaluations can be resource-intensive, posing challenges in terms of time, cost, and feasibility. Evaluation frameworks must balance the need for rigorous, context-sensitive assessments with the practical realities of AI development and deployment \cite{ref366}.

By embracing a socio-technical perspective in AI evaluation, researchers and practitioners can ensure that the development and deployment of these powerful technologies are aligned with societal values and the diverse needs of all affected stakeholders. This holistic approach not only mitigates the risks of AI but also unlocks the transformative potential of these technologies to positively impact individuals, communities, and the broader socio-technical landscape \cite{ref367,ref368}.


\begin{thebibliography}{368}
\addcontentsline{toc}{section}{References}
\bibitem{ref1} A Comprehensive Overview of Large Language Models. arXiv:2307.06435v9, 2023. \url{https://arxiv.org/abs/2307.06435v9}
\bibitem{ref2} A Survey of GPT-3 Family Large Language Models Including ChatGPT and GPT-4. arXiv:2310.12321v1, 2023. \url{https://arxiv.org/abs/2310.12321v1}
\bibitem{ref3} History, Development, and Principles of Large Language Models-An Introductory Survey. arXiv:2402.06853v1, 2024. \url{https://arxiv.org/abs/2402.06853v1}
\bibitem{ref4} Large Language Models for Telecom Forthcoming Impact on the Industry. arXiv:2308.06013v2, 2023. \url{https://arxiv.org/abs/2308.06013v2}
\bibitem{ref5} Large Language Models for Education A Survey and Outlook. arXiv:2403.18105v2, 2024. \url{https://arxiv.org/abs/2403.18105v2}
\bibitem{ref6} Apprentices to Research Assistants Advancing Research with Large Language Models. arXiv:2404.06404v1, 2024. \url{https://arxiv.org/abs/2404.06404v1}
\bibitem{ref7} People's Perceptions Toward Bias and Related Concepts in Large Language Models A Systematic Review. arXiv:2309.14504v2, 2023. \url{https://arxiv.org/abs/2309.14504v2}
\bibitem{ref8} Middleware for LLMs Tools Are Instrumental for Language Agents in Complex Environments. arXiv:2402.14672v1, 2024. \url{https://arxiv.org/abs/2402.14672v1}
\bibitem{ref9} Large Multimodal Agents A Survey. arXiv:2402.15116v1, 2024. \url{https://arxiv.org/abs/2402.15116v1}
\bibitem{ref10} Exploring Advanced Methodologies in Security Evaluation for LLMs. arXiv:2402.17970v2, 2024. \url{https://arxiv.org/abs/2402.17970v2}
\bibitem{ref11} Evaluating Large Language Models A Comprehensive Survey. arXiv:2310.19736v3, 2023. \url{https://arxiv.org/abs/2310.19736v3}
\bibitem{ref12} LooGLE Can Long-Context Language Models Understand Long Contexts. arXiv:2311.04939v1, 2023. \url{https://arxiv.org/abs/2311.04939v1}
\bibitem{ref13} A User-Centric Benchmark for Evaluating Large Language Models. arXiv:2404.13940v2, 2024. \url{https://arxiv.org/abs/2404.13940v2}
\bibitem{ref14} Benchmarking the Text-to-SQL Capability of Large Language Models A Comprehensive Evaluation. arXiv:2403.02951v2, 2024. \url{https://arxiv.org/abs/2403.02951v2}
\bibitem{ref15} LLMMaps -- A Visual Metaphor for Stratified Evaluation of Large Language Models. arXiv:2304.00457v3, 2023. \url{https://arxiv.org/abs/2304.00457v3}
\bibitem{ref16} Opportunities and Challenges of Applying Large Language Models in Building Energy Efficiency and Decarbonization Studies An Exploratory Overview. arXiv:2312.11701v1, 2023. \url{https://arxiv.org/abs/2312.11701v1}
\bibitem{ref17} LLMs for Cyber Security New Opportunities. arXiv:2404.11338v1, 2024. \url{https://arxiv.org/abs/2404.11338v1}
\bibitem{ref18} Large Language Models A Survey. arXiv:2402.06196v2, 2024. \url{https://arxiv.org/abs/2402.06196v2}
\bibitem{ref19} A survey on feature weighting based K-Means algorithms. arXiv:1601.03483v1, 2015. \url{https://arxiv.org/abs/1601.03483v1}
\bibitem{ref20} Reinforced Co-Training. arXiv:1804.06035v1, 2018. \url{https://arxiv.org/abs/1804.06035v1}
\bibitem{ref21} TransGPT Multi-modal Generative Pre-trained Transformer for Transportation. arXiv:2402.07233v1, 2024. \url{https://arxiv.org/abs/2402.07233v1}
\bibitem{ref22} Eight Things to Know about Large Language Models. arXiv:2304.00612v1, 2023. \url{https://arxiv.org/abs/2304.00612v1}
\bibitem{ref23} Datasets for Large Language Models A Comprehensive Survey. arXiv:2402.18041v1, 2024. \url{https://arxiv.org/abs/2402.18041v1}
\bibitem{ref24} GLaM Efficient Scaling of Language Models with Mixture-of-Experts. arXiv:2112.06905v2, 2021. \url{https://arxiv.org/abs/2112.06905v2}
\bibitem{ref25} Revealing the structure of language model capabilities. arXiv:2306.10062v1, 2023. \url{https://arxiv.org/abs/2306.10062v1}
\bibitem{ref26} The Cost of Down-Scaling Language Models Fact Recall Deteriorates before In-Context Learning. arXiv:2310.04680v1, 2023. \url{https://arxiv.org/abs/2310.04680v1}
\bibitem{ref27} Large Language Models The Need for Nuance in Current Debates and a Pragmatic Perspective on Understanding. arXiv:2310.19671v2, 2023. \url{https://arxiv.org/abs/2310.19671v2}
\bibitem{ref28} Rethinking Large Language Models in Mental Health Applications. arXiv:2311.11267v2, 2023. \url{https://arxiv.org/abs/2311.11267v2}
\bibitem{ref29} Large Language Models and Explainable Law a Hybrid Methodology. arXiv:2311.11811v1, 2023. \url{https://arxiv.org/abs/2311.11811v1}
\bibitem{ref30} Maximizing User Experience with LLMOps-Driven Personalized Recommendation Systems. arXiv:2404.00903v1, 2024. \url{https://arxiv.org/abs/2404.00903v1}
\bibitem{ref31} The Dark Side of ChatGPT Legal and Ethical Challenges from Stochastic Parrots and Hallucination. arXiv:2304.14347v1, 2023. \url{https://arxiv.org/abs/2304.14347v1}
\bibitem{ref32} Securing Large Language Models Threats, Vulnerabilities and Responsible Practices. arXiv:2403.12503v1, 2024. \url{https://arxiv.org/abs/2403.12503v1}
\bibitem{ref33} LLMs-Healthcare Current Applications and Challenges of Large Language Models in various Medical Specialties. arXiv:2311.12882v3, 2023. \url{https://arxiv.org/abs/2311.12882v3}
\bibitem{ref34} Creating Trustworthy LLMs Dealing with Hallucinations in Healthcare AI. arXiv:2311.01463v1, 2023. \url{https://arxiv.org/abs/2311.01463v1}
\bibitem{ref35} The Ethics of ChatGPT in Medicine and Healthcare A Systematic Review on Large Language Models (LLMs). arXiv:2403.14473v1, 2024. \url{https://arxiv.org/abs/2403.14473v1}
\bibitem{ref36} Exploring the Nexus of Large Language Models and Legal Systems A Short Survey. arXiv:2404.00990v1, 2024. \url{https://arxiv.org/abs/2404.00990v1}
\bibitem{ref37} Large Legal Fictions Profiling Legal Hallucinations in Large Language Models. arXiv:2401.01301v1, 2024. \url{https://arxiv.org/abs/2401.01301v1}
\bibitem{ref38} A Short Survey of Viewing Large Language Models in Legal Aspect. arXiv:2303.09136v1, 2023. \url{https://arxiv.org/abs/2303.09136v1}
\bibitem{ref39} Revolutionizing Finance with LLMs An Overview of Applications and Insights. arXiv:2401.11641v1, 2024. \url{https://arxiv.org/abs/2401.11641v1}
\bibitem{ref40} Vox Populi, Vox ChatGPT Large Language Models, Education and Democracy. arXiv:2311.06207v1, 2023. \url{https://arxiv.org/abs/2311.06207v1}
\bibitem{ref41} TransportationGames Benchmarking Transportation Knowledge of (Multimodal) Large Language Models. arXiv:2401.04471v1, 2024. \url{https://arxiv.org/abs/2401.04471v1}
\bibitem{ref42} LLeMpower Understanding Disparities in the Control and Access of Large Language Models. arXiv:2404.09356v1, 2024. \url{https://arxiv.org/abs/2404.09356v1}
\bibitem{ref43} Mapping LLM Security Landscapes A Comprehensive Stakeholder Risk Assessment Proposal. arXiv:2403.13309v1, 2024. \url{https://arxiv.org/abs/2403.13309v1}
\bibitem{ref44} Unveiling the Misuse Potential of Base Large Language Models via In-Context Learning. arXiv:2404.10552v1, 2024. \url{https://arxiv.org/abs/2404.10552v1}
\bibitem{ref45} Security and Privacy Challenges of Large Language Models A Survey. arXiv:2402.00888v1, 2024. \url{https://arxiv.org/abs/2402.00888v1}
\bibitem{ref46} The opportunities and risks of large language models in mental health. arXiv:2403.14814v2, 2024. \url{https://arxiv.org/abs/2403.14814v2}
\bibitem{ref47} From Representational Harms to Quality-of-Service Harms A Case Study on Llama 2 Safety Safeguards. arXiv:2403.13213v2, 2024. \url{https://arxiv.org/abs/2403.13213v2}
\bibitem{ref48} Towards Trustable Language Models Investigating Information Quality of Large Language Models. arXiv:2401.13086v1, 2024. \url{https://arxiv.org/abs/2401.13086v1}
\bibitem{ref49} Reliability Check An Analysis of GPT-3's Response to Sensitive Topics and Prompt Wording. arXiv:2306.06199v1, 2023. \url{https://arxiv.org/abs/2306.06199v1}
\bibitem{ref50} RAmBLA A Framework for Evaluating the Reliability of LLMs as Assistants in the Biomedical Domain. arXiv:2403.14578v1, 2024. \url{https://arxiv.org/abs/2403.14578v1}
\bibitem{ref51} Fairness of ChatGPT. arXiv:2305.18569v1, 2023. \url{https://arxiv.org/abs/2305.18569v1}
\bibitem{ref52} A Survey on Fairness in Large Language Models. arXiv:2308.10149v2, 2023. \url{https://arxiv.org/abs/2308.10149v2}
\bibitem{ref53} Use large language models to promote equity. arXiv:2312.14804v1, 2023. \url{https://arxiv.org/abs/2312.14804v1}
\bibitem{ref54} Tackling Bias in Pre-trained Language Models Current Trends and Under-represented Societies. arXiv:2312.01509v1, 2023. \url{https://arxiv.org/abs/2312.01509v1}
\bibitem{ref55} AI Transparency in the Age of LLMs A Human-Centered Research Roadmap. arXiv:2306.01941v2, 2023. \url{https://arxiv.org/abs/2306.01941v2}
\bibitem{ref56} Auditing large language models a three-layered approach. arXiv:2302.08500v2, 2023. \url{https://arxiv.org/abs/2302.08500v2}
\bibitem{ref57} Trustworthy Large Models in Vision A Survey. arXiv:2311.09680v5, 2023. \url{https://arxiv.org/abs/2311.09680v5}
\bibitem{ref58} Large Language Model Supply Chain A Research Agenda. arXiv:2404.12736v1, 2024. \url{https://arxiv.org/abs/2404.12736v1}
\bibitem{ref59} Scaling Dimension. arXiv:2302.09101v1, 2023. \url{https://arxiv.org/abs/2302.09101v1}
\bibitem{ref60} A Comprehensive Evaluation Framework for Deep Model Robustness. arXiv:2101.09617v2, 2021. \url{https://arxiv.org/abs/2101.09617v2}
\bibitem{ref61} A Survey on Evaluation of Large Language Models. arXiv:2307.03109v9, 2023. \url{https://arxiv.org/abs/2307.03109v9}
\bibitem{ref62} Surveying Attitudinal Alignment Between Large Language Models Vs. Humans Towards 17 Sustainable Development Goals. arXiv:2404.13885v1, 2024. \url{https://arxiv.org/abs/2404.13885v1}
\bibitem{ref63} Comprehensive Reassessment of Large-Scale Evaluation Outcomes in LLMs A Multifaceted Statistical Approach. arXiv:2403.15250v1, 2024. \url{https://arxiv.org/abs/2403.15250v1}
\bibitem{ref64} Rethinking Model Evaluation as Narrowing the Socio-Technical Gap. arXiv:2306.03100v3, 2023. \url{https://arxiv.org/abs/2306.03100v3}
\bibitem{ref65} Post Turing Mapping the landscape of LLM Evaluation. arXiv:2311.02049v1, 2023. \url{https://arxiv.org/abs/2311.02049v1}
\bibitem{ref66} Collaborative Evaluation Exploring the Synergy of Large Language Models and Humans for Open-ended Generation Evaluation. arXiv:2310.19740v1, 2023. \url{https://arxiv.org/abs/2310.19740v1}
\bibitem{ref67} Assessment of Multimodal Large Language Models in Alignment with Human Values. arXiv:2403.17830v1, 2024. \url{https://arxiv.org/abs/2403.17830v1}
\bibitem{ref68} TencentLLMEval A Hierarchical Evaluation of Real-World Capabilities for Human-Aligned LLMs. arXiv:2311.05374v1, 2023. \url{https://arxiv.org/abs/2311.05374v1}
\bibitem{ref69} Are LLM-based Evaluators Confusing NLG Quality Criteria. arXiv:2402.12055v1, 2024. \url{https://arxiv.org/abs/2402.12055v1}
\bibitem{ref70} Towards Optimizing with Large Language Models. arXiv:2310.05204v2, 2023. \url{https://arxiv.org/abs/2310.05204v2}
\bibitem{ref71} Who is ChatGPT Benchmarking LLMs' Psychological Portrayal Using PsychoBench. arXiv:2310.01386v2, 2023. \url{https://arxiv.org/abs/2310.01386v2}
\bibitem{ref72} CriticBench Evaluating Large Language Models as Critic. arXiv:2402.13764v3, 2024. \url{https://arxiv.org/abs/2402.13764v3}
\bibitem{ref73} Beyond Human Norms Unveiling Unique Values of Large Language Models through Interdisciplinary Approaches. arXiv:2404.12744v1, 2024. \url{https://arxiv.org/abs/2404.12744v1}
\bibitem{ref74} Pragmatic Competence Evaluation of Large Language Models for Korean. arXiv:2403.12675v1, 2024. \url{https://arxiv.org/abs/2403.12675v1}
\bibitem{ref75} LawBench Benchmarking Legal Knowledge of Large Language Models. arXiv:2309.16289v1, 2023. \url{https://arxiv.org/abs/2309.16289v1}
\bibitem{ref76} FreeEval A Modular Framework for Trustworthy and Efficient Evaluation of Large Language Models. arXiv:2404.06003v1, 2024. \url{https://arxiv.org/abs/2404.06003v1}
\bibitem{ref77} CheckEval Robust Evaluation Framework using Large Language Model via Checklist. arXiv:2403.18771v1, 2024. \url{https://arxiv.org/abs/2403.18771v1}
\bibitem{ref78} HD-Eval Aligning Large Language Model Evaluators Through Hierarchical Criteria Decomposition. arXiv:2402.15754v1, 2024. \url{https://arxiv.org/abs/2402.15754v1}
\bibitem{ref79} UltraEval A Lightweight Platform for Flexible and Comprehensive Evaluation for LLMs. arXiv:2404.07584v1, 2024. \url{https://arxiv.org/abs/2404.07584v1}
\bibitem{ref80} PRE A Peer Review Based Large Language Model Evaluator. arXiv:2401.15641v1, 2024. \url{https://arxiv.org/abs/2401.15641v1}
\bibitem{ref81} A Toolbox for Surfacing Health Equity Harms and Biases in Large Language Models. arXiv:2403.12025v1, 2024. \url{https://arxiv.org/abs/2403.12025v1}
\bibitem{ref82} Emulating Human Cognitive Processes for Expert-Level Medical Question-Answering with Large Language Models. arXiv:2310.11266v1, 2023. \url{https://arxiv.org/abs/2310.11266v1}
\bibitem{ref83} GPT-Fathom Benchmarking Large Language Models to Decipher the Evolutionary Path towards GPT-4 and Beyond. arXiv:2309.16583v6, 2023. \url{https://arxiv.org/abs/2309.16583v6}
\bibitem{ref84} PiCO Peer Review in LLMs based on the Consistency Optimization. arXiv:2402.01830v2, 2024. \url{https://arxiv.org/abs/2402.01830v2}
\bibitem{ref85} MedAlpaca -- An Open-Source Collection of Medical Conversational AI Models and Training Data. arXiv:2304.08247v2, 2023. \url{https://arxiv.org/abs/2304.08247v2}
\bibitem{ref86} Hippocrates An Open-Source Framework for Advancing Large Language Models in Healthcare. arXiv:2404.16621v1, 2024. \url{https://arxiv.org/abs/2404.16621v1}
\bibitem{ref87} Towards Automatic Evaluation for LLMs' Clinical Capabilities Metric, Data, and Algorithm. arXiv:2403.16446v1, 2024. \url{https://arxiv.org/abs/2403.16446v1}
\bibitem{ref88} GPT Models in Construction Industry Opportunities, Limitations, and a Use Case Validation. arXiv:2305.18997v1, 2023. \url{https://arxiv.org/abs/2305.18997v1}
\bibitem{ref89} L-Eval Instituting Standardized Evaluation for Long Context Language Models. arXiv:2307.11088v3, 2023. \url{https://arxiv.org/abs/2307.11088v3}
\bibitem{ref90} Perception Score, A Learned Metric for Open-ended Text Generation Evaluation. arXiv:2008.03082v2, 2020. \url{https://arxiv.org/abs/2008.03082v2}
\bibitem{ref91} The Perils of Using Mechanical Turk to Evaluate Open-Ended Text Generation. arXiv:2109.06835v1, 2021. \url{https://arxiv.org/abs/2109.06835v1}
\bibitem{ref92} Deconstructing NLG Evaluation Evaluation Practices, Assumptions, and Their Implications. arXiv:2205.06828v1, 2022. \url{https://arxiv.org/abs/2205.06828v1}
\bibitem{ref93} Unifying Human and Statistical Evaluation for Natural Language Generation. arXiv:1904.02792v1, 2019. \url{https://arxiv.org/abs/1904.02792v1}
\bibitem{ref94} Perceived Trustworthiness of Natural Language Generators. arXiv:2305.18176v1, 2023. \url{https://arxiv.org/abs/2305.18176v1}
\bibitem{ref95} Judge the Judges A Large-Scale Evaluation Study of Neural Language Models for Online Review Generation. arXiv:1901.00398v2, 2019. \url{https://arxiv.org/abs/1901.00398v2}
\bibitem{ref96} Artificial Intelligence versus Maya Angelou Experimental evidence that people cannot differentiate AI-generated from human-written poetry. arXiv:2005.09980v2, 2020. \url{https://arxiv.org/abs/2005.09980v2}
\bibitem{ref97} The Authenticity Gap in Human Evaluation. arXiv:2205.11930v2, 2022. \url{https://arxiv.org/abs/2205.11930v2}
\bibitem{ref98} Toward Verifiable and Reproducible Human Evaluation for Text-to-Image Generation. arXiv:2304.01816v1, 2023. \url{https://arxiv.org/abs/2304.01816v1}
\bibitem{ref99} Evaluating Human-Language Model Interaction. arXiv:2212.09746v5, 2022. \url{https://arxiv.org/abs/2212.09746v5}
\bibitem{ref100} GENIE Toward Reproducible and Standardized Human Evaluation for Text Generation. arXiv:2101.06561v4, 2021. \url{https://arxiv.org/abs/2101.06561v4}
\bibitem{ref101} Perturbation CheckLists for Evaluating NLG Evaluation Metrics. arXiv:2109.05771v1, 2021. \url{https://arxiv.org/abs/2109.05771v1}
\bibitem{ref102} Multi-Dimensional Evaluation of Text Summarization with In-Context Learning. arXiv:2306.01200v1, 2023. \url{https://arxiv.org/abs/2306.01200v1}
\bibitem{ref103} Can Large Language Models Be an Alternative to Human Evaluations. arXiv:2305.01937v1, 2023. \url{https://arxiv.org/abs/2305.01937v1}
\bibitem{ref104} Is ChatGPT a Good NLG Evaluator A Preliminary Study. arXiv:2303.04048v3, 2023. \url{https://arxiv.org/abs/2303.04048v3}
\bibitem{ref105} Estimating Subjective Crowd-Evaluations as an Additional Objective to Improve Natural Language Generation. arXiv:2104.05224v1, 2021. \url{https://arxiv.org/abs/2104.05224v1}
\bibitem{ref106} Human or Machine Automating Human Likeliness Evaluation of NLG Texts. arXiv:2006.03189v1, 2020. \url{https://arxiv.org/abs/2006.03189v1}
\bibitem{ref107} Calibrating LLM-Based Evaluator. arXiv:2309.13308v1, 2023. \url{https://arxiv.org/abs/2309.13308v1}
\bibitem{ref108} Tangled up in BLEU Reevaluating the Evaluation of Automatic Machine Translation Evaluation Metrics. arXiv:2006.06264v2, 2020. \url{https://arxiv.org/abs/2006.06264v2}
\bibitem{ref109} MT Metrics Correlate with Human Ratings of Simultaneous Speech Translation. arXiv:2211.08633v2, 2022. \url{https://arxiv.org/abs/2211.08633v2}
\bibitem{ref110} Why We Need New Evaluation Metrics for NLG. arXiv:1707.06875v1, 2017. \url{https://arxiv.org/abs/1707.06875v1}
\bibitem{ref111} Revisiting Grammatical Error Correction Evaluation and Beyond. arXiv:2211.01635v1, 2022. \url{https://arxiv.org/abs/2211.01635v1}
\bibitem{ref112} There's No Comparison Reference-less Evaluation Metrics in Grammatical Error Correction. arXiv:1610.02124v1, 2016. \url{https://arxiv.org/abs/1610.02124v1}
\bibitem{ref113} BERTScore Evaluating Text Generation with BERT. arXiv:1904.09675v3, 2019. \url{https://arxiv.org/abs/1904.09675v3}
\bibitem{ref114} BLEURT Learning Robust Metrics for Text Generation. arXiv:2004.04696v5, 2020. \url{https://arxiv.org/abs/2004.04696v5}
\bibitem{ref115} COMET A Neural Framework for MT Evaluation. arXiv:2009.09025v2, 2020. \url{https://arxiv.org/abs/2009.09025v2}
\bibitem{ref116} BLEU Meets COMET Combining Lexical and Neural Metrics Towards Robust Machine Translation Evaluation. arXiv:2305.19144v1, 2023. \url{https://arxiv.org/abs/2305.19144v1}
\bibitem{ref117} Layer or Representation Space What makes BERT-based Evaluation Metrics Robust. arXiv:2209.02317v2, 2022. \url{https://arxiv.org/abs/2209.02317v2}
\bibitem{ref118} Towards Interpretable and Efficient Automatic Reference-Based Summarization Evaluation. arXiv:2303.03608v2, 2023. \url{https://arxiv.org/abs/2303.03608v2}
\bibitem{ref119} Robustness Tests for Automatic Machine Translation Metrics with Adversarial Attacks. arXiv:2311.00508v1, 2023. \url{https://arxiv.org/abs/2311.00508v1}
\bibitem{ref120} Curious Case of Language Generation Evaluation Metrics A Cautionary Tale. arXiv:2010.13588v1, 2020. \url{https://arxiv.org/abs/2010.13588v1}
\bibitem{ref121} SMART Sentences as Basic Units for Text Evaluation. arXiv:2208.01030v1, 2022. \url{https://arxiv.org/abs/2208.01030v1}
\bibitem{ref122} MISMATCH Fine-grained Evaluation of Machine-generated Text with Mismatch Error Types. arXiv:2306.10452v1, 2023. \url{https://arxiv.org/abs/2306.10452v1}
\bibitem{ref123} Aligning language models with human preferences. arXiv:2404.12150v1, 2024. \url{https://arxiv.org/abs/2404.12150v1}
\bibitem{ref124} Bias patterns in the application of LLMs for clinical decision support A comprehensive study. arXiv:2404.15149v1, 2024. \url{https://arxiv.org/abs/2404.15149v1}
\bibitem{ref125} Large Language Models are Not Yet Human-Level Evaluators for Abstractive Summarization. arXiv:2305.13091v2, 2023. \url{https://arxiv.org/abs/2305.13091v2}
\bibitem{ref126} Likelihood-based Mitigation of Evaluation Bias in Large Language Models. arXiv:2402.15987v2, 2024. \url{https://arxiv.org/abs/2402.15987v2}
\bibitem{ref127} Benchmarking Cognitive Biases in Large Language Models as Evaluators. arXiv:2309.17012v1, 2023. \url{https://arxiv.org/abs/2309.17012v1}
\bibitem{ref128} Experts, Errors, and Context A Large-Scale Study of Human Evaluation for Machine Translation. arXiv:2104.14478v1, 2021. \url{https://arxiv.org/abs/2104.14478v1}
\bibitem{ref129} Can Large Language Models be Trusted for Evaluation Scalable Meta-Evaluation of LLMs as Evaluators via Agent Debate. arXiv:2401.16788v1, 2024. \url{https://arxiv.org/abs/2401.16788v1}
\bibitem{ref130} Style Over Substance Evaluation Biases for Large Language Models. arXiv:2307.03025v3, 2023. \url{https://arxiv.org/abs/2307.03025v3}
\bibitem{ref131} MATEval A Multi-Agent Discussion Framework for Advancing Open-Ended Text Evaluation. arXiv:2403.19305v2, 2024. \url{https://arxiv.org/abs/2403.19305v2}
\bibitem{ref132} ChatEval Towards Better LLM-based Evaluators through Multi-Agent Debate. arXiv:2308.07201v1, 2023. \url{https://arxiv.org/abs/2308.07201v1}
\bibitem{ref133} CoAScore Chain-of-Aspects Prompting for NLG Evaluation. arXiv:2312.10355v1, 2023. \url{https://arxiv.org/abs/2312.10355v1}
\bibitem{ref134} Evaluation Gaps in Machine Learning Practice. arXiv:2205.05256v1, 2022. \url{https://arxiv.org/abs/2205.05256v1}
\bibitem{ref135} Adversarial Tradeoffs in Robust State Estimation. arXiv:2111.08864v3, 2021. \url{https://arxiv.org/abs/2111.08864v3}
\bibitem{ref136} Causal Fair Metric Bridging Causality, Individual Fairness, and Adversarial Robustness. arXiv:2310.19391v2, 2023. \url{https://arxiv.org/abs/2310.19391v2}
\bibitem{ref137} On the Sensitivity of Adversarial Robustness to Input Data Distributions. arXiv:1902.08336v1, 2019. \url{https://arxiv.org/abs/1902.08336v1}
\bibitem{ref138} Revisiting Ensembles in an Adversarial Context Improving Natural Accuracy. arXiv:2002.11572v1, 2020. \url{https://arxiv.org/abs/2002.11572v1}
\bibitem{ref139} Efficient Estimation of Average-Case Robustness for Multi-Class Classification. arXiv:2307.13885v4, 2023. \url{https://arxiv.org/abs/2307.13885v4}
\bibitem{ref140} Evaluating Multimodal Interactive Agents. arXiv:2205.13274v2, 2022. \url{https://arxiv.org/abs/2205.13274v2}
\bibitem{ref141} Substance over Style Document-Level Targeted Content Transfer. arXiv:2010.08618v1, 2020. \url{https://arxiv.org/abs/2010.08618v1}
\bibitem{ref142} Humans or LLMs as the Judge A Study on Judgement Biases. arXiv:2402.10669v3, 2024. \url{https://arxiv.org/abs/2402.10669v3}
\bibitem{ref143} Aligning Robot and Human Representations. arXiv:2302.01928v2, 2023. \url{https://arxiv.org/abs/2302.01928v2}
\bibitem{ref144} Re-evaluating Evaluation. arXiv:1806.02643v2, 2018. \url{https://arxiv.org/abs/1806.02643v2}
\bibitem{ref145} Ensemble Sampling. arXiv:1705.07347v4, 2017. \url{https://arxiv.org/abs/1705.07347v4}
\bibitem{ref146} Prompt Valuation Based on Shapley Values. arXiv:2312.15395v1, 2023. \url{https://arxiv.org/abs/2312.15395v1}
\bibitem{ref147} Developing a Framework for Auditing Large Language Models Using Human-in-the-Loop. arXiv:2402.09346v2, 2024. \url{https://arxiv.org/abs/2402.09346v2}
\bibitem{ref148} Towards Understanding and Mitigating Social Biases in Language Models. arXiv:2106.13219v1, 2021. \url{https://arxiv.org/abs/2106.13219v1}
\bibitem{ref149} Towards Precise Observations of Neural Model Robustness in Classification. arXiv:2404.16457v1, 2024. \url{https://arxiv.org/abs/2404.16457v1}
\bibitem{ref150} Think Twice Measuring the Efficiency of Eliminating Prediction Shortcuts of Question Answering Models. arXiv:2305.06841v2, 2023. \url{https://arxiv.org/abs/2305.06841v2}
\bibitem{ref151} Bias and Fairness in Large Language Models A Survey. arXiv:2309.00770v2, 2023. \url{https://arxiv.org/abs/2309.00770v2}
\bibitem{ref152} Measuring Implicit Bias in Explicitly Unbiased Large Language Models. arXiv:2402.04105v1, 2024. \url{https://arxiv.org/abs/2402.04105v1}
\bibitem{ref153} ROBBIE Robust Bias Evaluation of Large Generative Language Models. arXiv:2311.18140v1, 2023. \url{https://arxiv.org/abs/2311.18140v1}
\bibitem{ref154} LLMs as Factual Reasoners Insights from Existing Benchmarks and Beyond. arXiv:2305.14540v1, 2023. \url{https://arxiv.org/abs/2305.14540v1}
\bibitem{ref155} Generating with Confidence Uncertainty Quantification for Black-box Large Language Models. arXiv:2305.19187v2, 2023. \url{https://arxiv.org/abs/2305.19187v2}
\bibitem{ref156} Detectors for Safe and Reliable LLMs Implementations, Uses, and Limitations. arXiv:2403.06009v1, 2024. \url{https://arxiv.org/abs/2403.06009v1}
\bibitem{ref157} Reliability Testing for Natural Language Processing Systems. arXiv:2105.02590v3, 2021. \url{https://arxiv.org/abs/2105.02590v3}
\bibitem{ref158} A Holistic Assessment of the Reliability of Machine Learning Systems. arXiv:2307.10586v2, 2023. \url{https://arxiv.org/abs/2307.10586v2}
\bibitem{ref159} Modelling Human Values for AI Reasoning. arXiv:2402.06359v1, 2024. \url{https://arxiv.org/abs/2402.06359v1}
\bibitem{ref160} Denevil Towards Deciphering and Navigating the Ethical Values of Large Language Models via Instruction Learning. arXiv:2310.11053v3, 2023. \url{https://arxiv.org/abs/2310.11053v3}
\bibitem{ref161} MoCa Measuring Human-Language Model Alignment on Causal and Moral Judgment Tasks. arXiv:2310.19677v2, 2023. \url{https://arxiv.org/abs/2310.19677v2}
\bibitem{ref162} Heterogeneous Value Alignment Evaluation for Large Language Models. arXiv:2305.17147v3, 2023. \url{https://arxiv.org/abs/2305.17147v3}
\bibitem{ref163} Personalisation within bounds A risk taxonomy and policy framework for the alignment of large language models with personalised feedback. arXiv:2303.05453v1, 2023. \url{https://arxiv.org/abs/2303.05453v1}
\bibitem{ref164} Personalized Soups Personalized Large Language Model Alignment via Post-hoc Parameter Merging. arXiv:2310.11564v1, 2023. \url{https://arxiv.org/abs/2310.11564v1}
\bibitem{ref165} Bridging the Gulf of Envisioning Cognitive Design Challenges in LLM Interfaces. arXiv:2309.14459v2, 2023. \url{https://arxiv.org/abs/2309.14459v2}
\bibitem{ref166} The Empty Signifier Problem Towards Clearer Paradigms for Operationalising Alignment in Large Language Models. arXiv:2310.02457v2, 2023. \url{https://arxiv.org/abs/2310.02457v2}
\bibitem{ref167} A Moral Imperative The Need for Continual Superalignment of Large Language Models. arXiv:2403.14683v1, 2024. \url{https://arxiv.org/abs/2403.14683v1}
\bibitem{ref168} Exploring the psychology of LLMs' Moral and Legal Reasoning. arXiv:2308.01264v2, 2023. \url{https://arxiv.org/abs/2308.01264v2}
\bibitem{ref169} Relying on the Unreliable The Impact of Language Models' Reluctance to Express Uncertainty. arXiv:2401.06730v1, 2024. \url{https://arxiv.org/abs/2401.06730v1}
\bibitem{ref170} The Dawn After the Dark An Empirical Study on Factuality Hallucination in Large Language Models. arXiv:2401.03205v1, 2024. \url{https://arxiv.org/abs/2401.03205v1}
\bibitem{ref171} Unsupervised Real-Time Hallucination Detection based on the Internal States of Large Language Models. arXiv:2403.06448v1, 2024. \url{https://arxiv.org/abs/2403.06448v1}
\bibitem{ref172} Measuring and Reducing LLM Hallucination without Gold-Standard Answers via Expertise-Weighting. arXiv:2402.10412v1, 2024. \url{https://arxiv.org/abs/2402.10412v1}
\bibitem{ref173} Fact-Checking the Output of Large Language Models via Token-Level Uncertainty Quantification. arXiv:2403.04696v1, 2024. \url{https://arxiv.org/abs/2403.04696v1}
\bibitem{ref174} Siren's Song in the AI Ocean A Survey on Hallucination in Large Language Models. arXiv:2309.01219v2, 2023. \url{https://arxiv.org/abs/2309.01219v2}
\bibitem{ref175} Towards Mitigating Hallucination in Large Language Models via Self-Reflection. arXiv:2310.06271v1, 2023. \url{https://arxiv.org/abs/2310.06271v1}
\bibitem{ref176} DelucionQA Detecting Hallucinations in Domain-specific Question Answering. arXiv:2312.05200v1, 2023. \url{https://arxiv.org/abs/2312.05200v1}
\bibitem{ref177} Quantifying and Attributing the Hallucination of Large Language Models via Association Analysis. arXiv:2309.05217v1, 2023. \url{https://arxiv.org/abs/2309.05217v1}
\bibitem{ref178} Despite super-human performance, current LLMs are unsuited for decisions about ethics and safety. arXiv:2212.06295v1, 2022. \url{https://arxiv.org/abs/2212.06295v1}
\bibitem{ref179} Explainability Is in the Mind of the Beholder Establishing the Foundations of Explainable Artificial Intelligence. arXiv:2112.14466v2, 2021. \url{https://arxiv.org/abs/2112.14466v2}
\bibitem{ref180} Local Interpretability of Calibrated Prediction Models A Case of Type 2 Diabetes Mellitus Screening Test. arXiv:2006.13815v1, 2020. \url{https://arxiv.org/abs/2006.13815v1}
\bibitem{ref181} What Makes a Good Explanation A Harmonized View of Properties of Explanations. arXiv:2211.05667v2, 2022. \url{https://arxiv.org/abs/2211.05667v2}
\bibitem{ref182} Towards Explainability and Fairness in Swiss Judgement Prediction Benchmarking on a Multilingual Dataset. arXiv:2402.17013v1, 2024. \url{https://arxiv.org/abs/2402.17013v1}
\bibitem{ref183} Explaining black boxes with a SMILE Statistical Model-agnostic Interpretability with Local Explanations. arXiv:2311.07286v1, 2023. \url{https://arxiv.org/abs/2311.07286v1}
\bibitem{ref184} LM Transparency Tool Interactive Tool for Analyzing Transformer Language Models. arXiv:2404.07004v1, 2024. \url{https://arxiv.org/abs/2404.07004v1}
\bibitem{ref185} Competition of Mechanisms Tracing How Language Models Handle Facts and Counterfactuals. arXiv:2402.11655v1, 2024. \url{https://arxiv.org/abs/2402.11655v1}
\bibitem{ref186} From Anecdotal Evidence to Quantitative Evaluation Methods A Systematic Review on Evaluating Explainable AI. arXiv:2201.08164v3, 2022. \url{https://arxiv.org/abs/2201.08164v3}
\bibitem{ref187} Quantifying Interpretability and Trust in Machine Learning Systems. arXiv:1901.08558v1, 2019. \url{https://arxiv.org/abs/1901.08558v1}
\bibitem{ref188} The Promise and Peril of Human Evaluation for Model Interpretability. arXiv:1711.07414v2, 2017. \url{https://arxiv.org/abs/1711.07414v2}
\bibitem{ref189} Human Factors in Model Interpretability Industry Practices, Challenges, and Needs. arXiv:2004.11440v2, 2020. \url{https://arxiv.org/abs/2004.11440v2}
\bibitem{ref190} Model Transparency and Interpretability Survey and Application to the Insurance Industry. arXiv:2209.00562v1, 2022. \url{https://arxiv.org/abs/2209.00562v1}
\bibitem{ref191} Explaining Explanations in AI. arXiv:1811.01439v1, 2018. \url{https://arxiv.org/abs/1811.01439v1}
\bibitem{ref192} Paragraph-level Rationale Extraction through Regularization A case study on European Court of Human Rights Cases. arXiv:2103.13084v1, 2021. \url{https://arxiv.org/abs/2103.13084v1}
\bibitem{ref193} Intuitively Assessing ML Model Reliability through Example-Based Explanations and Editing Model Inputs. arXiv:2102.08540v2, 2021. \url{https://arxiv.org/abs/2102.08540v2}
\bibitem{ref194} Uncertainty in Gradient Boosting via Ensembles. arXiv:2006.10562v4, 2020. \url{https://arxiv.org/abs/2006.10562v4}
\bibitem{ref195} Ensembles for Uncertainty Estimation Benefits of Prior Functions and Bootstrapping. arXiv:2206.03633v1, 2022. \url{https://arxiv.org/abs/2206.03633v1}
\bibitem{ref196} Calibrating Ensembles for Scalable Uncertainty Quantification in Deep Learning-based Medical Segmentation. arXiv:2209.09563v1, 2022. \url{https://arxiv.org/abs/2209.09563v1}
\bibitem{ref197} Uncertainty Quantification in Machine Learning for Engineering Design and Health Prognostics A Tutorial. arXiv:2305.04933v2, 2023. \url{https://arxiv.org/abs/2305.04933v2}
\bibitem{ref198} Deep Ensembles from a Bayesian Perspective. arXiv:2105.13283v2, 2021. \url{https://arxiv.org/abs/2105.13283v2}
\bibitem{ref199} LoRA ensembles for large language model fine-tuning. arXiv:2310.00035v2, 2023. \url{https://arxiv.org/abs/2310.00035v2}
\bibitem{ref200} Window-Based Early-Exit Cascades for Uncertainty Estimation When Deep Ensembles are More Efficient than Single Models. arXiv:2303.08010v3, 2023. \url{https://arxiv.org/abs/2303.08010v3}
\bibitem{ref201} Quantifying Uncertainties in Natural Language Processing Tasks. arXiv:1811.07253v1, 2018. \url{https://arxiv.org/abs/1811.07253v1}
\bibitem{ref202} Towards Clear Expectations for Uncertainty Estimation. arXiv:2207.13341v1, 2022. \url{https://arxiv.org/abs/2207.13341v1}
\bibitem{ref203} Aligning Language Models to User Opinions. arXiv:2305.14929v1, 2023. \url{https://arxiv.org/abs/2305.14929v1}
\bibitem{ref204} Group Preference Optimization Few-Shot Alignment of Large Language Models. arXiv:2310.11523v1, 2023. \url{https://arxiv.org/abs/2310.11523v1}
\bibitem{ref205} Democratizing Large Language Models via Personalized Parameter-Efficient Fine-tuning. arXiv:2402.04401v1, 2024. \url{https://arxiv.org/abs/2402.04401v1}
\bibitem{ref206} TrustLLM Trustworthiness in Large Language Models. arXiv:2401.05561v4, 2024. \url{https://arxiv.org/abs/2401.05561v4}
\bibitem{ref207} Flames Benchmarking Value Alignment of LLMs in Chinese. arXiv:2311.06899v4, 2023. \url{https://arxiv.org/abs/2311.06899v4}
\bibitem{ref208} AgentBoard An Analytical Evaluation Board of Multi-turn LLM Agents. arXiv:2401.13178v1, 2024. \url{https://arxiv.org/abs/2401.13178v1}
\bibitem{ref209} Benchmarking LLMs via Uncertainty Quantification. arXiv:2401.12794v2, 2024. \url{https://arxiv.org/abs/2401.12794v2}
\bibitem{ref210} Integrating UMLS Knowledge into Large Language Models for Medical Question Answering. arXiv:2310.02778v2, 2023. \url{https://arxiv.org/abs/2310.02778v2}
\bibitem{ref211} Deciphering Diagnoses How Large Language Models Explanations Influence Clinical Decision Making. arXiv:2310.01708v1, 2023. \url{https://arxiv.org/abs/2310.01708v1}
\bibitem{ref212} Large Language Models as Agents in the Clinic. arXiv:2309.10895v1, 2023. \url{https://arxiv.org/abs/2309.10895v1}
\bibitem{ref213} Autonomous Artificial Intelligence Agents for Clinical Decision Making in Oncology. arXiv:2404.04667v1, 2024. \url{https://arxiv.org/abs/2404.04667v1}
\bibitem{ref214} Bespoke Large Language Models for Digital Triage Assistance in Mental Health Care. arXiv:2403.19790v1, 2024. \url{https://arxiv.org/abs/2403.19790v1}
\bibitem{ref215} A Visual Analytics System for Multi-model Comparison on Clinical Data Predictions. arXiv:2002.10998v2, 2020. \url{https://arxiv.org/abs/2002.10998v2}
\bibitem{ref216} Updating the Minimum Information about CLinical Artificial Intelligence (MI-CLAIM) checklist for generative modeling research. arXiv:2403.02558v1, 2024. \url{https://arxiv.org/abs/2403.02558v1}
\bibitem{ref217} Large Language Models Encode Clinical Knowledge. arXiv:2212.13138v1, 2022. \url{https://arxiv.org/abs/2212.13138v1}
\bibitem{ref218} Diagnostic Reasoning Prompts Reveal the Potential for Large Language Model Interpretability in Medicine. arXiv:2308.06834v1, 2023. \url{https://arxiv.org/abs/2308.06834v1}
\bibitem{ref219} Explainable Deep Learning in Healthcare A Methodological Survey from an Attribution View. arXiv:2112.02625v1, 2021. \url{https://arxiv.org/abs/2112.02625v1}
\bibitem{ref220} Augmenting Black-box LLMs with Medical Textbooks for Clinical Question Answering. arXiv:2309.02233v2, 2023. \url{https://arxiv.org/abs/2309.02233v2}
\bibitem{ref221} Prototyping the use of Large Language Models (LLMs) for adult learning content creation at scale. arXiv:2306.01815v1, 2023. \url{https://arxiv.org/abs/2306.01815v1}
\bibitem{ref222} Adaptive and Personalized Exercise Generation for Online Language Learning. arXiv:2306.02457v1, 2023. \url{https://arxiv.org/abs/2306.02457v1}
\bibitem{ref223} Empowering Personalized Learning through a Conversation-based Tutoring System with Student Modeling. arXiv:2403.14071v1, 2024. \url{https://arxiv.org/abs/2403.14071v1}
\bibitem{ref224} Practical and Ethical Challenges of Large Language Models in Education A Systematic Scoping Review. arXiv:2303.13379v2, 2023. \url{https://arxiv.org/abs/2303.13379v2}
\bibitem{ref225} SciAssess Benchmarking LLM Proficiency in Scientific Literature Analysis. arXiv:2403.01976v2, 2024. \url{https://arxiv.org/abs/2403.01976v2}
\bibitem{ref226} Large Language Models are Zero Shot Hypothesis Proposers. arXiv:2311.05965v1, 2023. \url{https://arxiv.org/abs/2311.05965v1}
\bibitem{ref227} Can Large Language Models Transform Computational Social Science. arXiv:2305.03514v3, 2023. \url{https://arxiv.org/abs/2305.03514v3}
\bibitem{ref228} Large Language Models are Geographically Biased. arXiv:2402.02680v1, 2024. \url{https://arxiv.org/abs/2402.02680v1}
\bibitem{ref229} Rethinking STS and NLI in Large Language Models. arXiv:2309.08969v2, 2023. \url{https://arxiv.org/abs/2309.08969v2}
\bibitem{ref230} AgentBench Evaluating LLMs as Agents. arXiv:2308.03688v2, 2023. \url{https://arxiv.org/abs/2308.03688v2}
\bibitem{ref231} LLM As DBA. arXiv:2308.05481v2, 2023. \url{https://arxiv.org/abs/2308.05481v2}
\bibitem{ref232} Sketched MinDist. arXiv:1907.02171v2, 2019. \url{https://arxiv.org/abs/1907.02171v2}
\bibitem{ref233} Understanding the Weakness of Large Language Model Agents within a Complex Android Environment. arXiv:2402.06596v1, 2024. \url{https://arxiv.org/abs/2402.06596v1}
\bibitem{ref234} Formal Methods with a Touch of Magic. arXiv:2005.12175v2, 2020. \url{https://arxiv.org/abs/2005.12175v2}
\bibitem{ref235} Meta-Neighborhoods. arXiv:1909.09140v3, 2019. \url{https://arxiv.org/abs/1909.09140v3}
\bibitem{ref236} Rethinking Interpretability in the Era of Large Language Models. arXiv:2402.01761v1, 2024. \url{https://arxiv.org/abs/2402.01761v1}
\bibitem{ref237} Towards Uncovering How Large Language Model Works An Explainability Perspective. arXiv:2402.10688v2, 2024. \url{https://arxiv.org/abs/2402.10688v2}
\bibitem{ref238} The Ethics of Interaction Mitigating Security Threats in LLMs. arXiv:2401.12273v1, 2024. \url{https://arxiv.org/abs/2401.12273v1}
\bibitem{ref239} Model-Agnostic Interpretation Framework in Machine Learning A Comparative Study in NBA Sports. arXiv:2401.02630v1, 2024. \url{https://arxiv.org/abs/2401.02630v1}
\bibitem{ref240} Caveat Lector Large Language Models in Legal Practice. arXiv:2403.09163v1, 2024. \url{https://arxiv.org/abs/2403.09163v1}
\bibitem{ref241} Just Like Me The Role of Opinions and Personal Experiences in The Perception of Explanations in Subjective Decision-Making. arXiv:2404.12558v1, 2024. \url{https://arxiv.org/abs/2404.12558v1}
\bibitem{ref242} Sparsity-Guided Holistic Explanation for LLMs with Interpretable Inference-Time Intervention. arXiv:2312.15033v1, 2023. \url{https://arxiv.org/abs/2312.15033v1}
\bibitem{ref243} Investigating the Encoding of Words in BERT's Neurons using Feature Textualization. arXiv:2311.08240v1, 2023. \url{https://arxiv.org/abs/2311.08240v1}
\bibitem{ref244} Unveiling LLMs The Evolution of Latent Representations in a Temporal Knowledge Graph. arXiv:2404.03623v1, 2024. \url{https://arxiv.org/abs/2404.03623v1}
\bibitem{ref245} Patchscopes A Unifying Framework for Inspecting Hidden Representations of Language Models. arXiv:2401.06102v2, 2024. \url{https://arxiv.org/abs/2401.06102v2}
\bibitem{ref246} Not All Models Localize Linguistic Knowledge in the Same Place A Layer-wise Probing on BERToids' Representations. arXiv:2109.05958v2, 2021. \url{https://arxiv.org/abs/2109.05958v2}
\bibitem{ref247} Introducing Orthogonal Constraint in Structural Probes. arXiv:2012.15228v2, 2020. \url{https://arxiv.org/abs/2012.15228v2}
\bibitem{ref248} Probing the Probing Paradigm Does Probing Accuracy Entail Task Relevance. arXiv:2005.00719v3, 2020. \url{https://arxiv.org/abs/2005.00719v3}
\bibitem{ref249} Conditional probing measuring usable information beyond a baseline. arXiv:2109.09234v1, 2021. \url{https://arxiv.org/abs/2109.09234v1}
\bibitem{ref250} Learning with Memory Embeddings. arXiv:1511.07972v9, 2015. \url{https://arxiv.org/abs/1511.07972v9}
\bibitem{ref251} The Geometry of Truth Emergent Linear Structure in Large Language Model Representations of True False Datasets. arXiv:2310.06824v2, 2023. \url{https://arxiv.org/abs/2310.06824v2}
\bibitem{ref252} Dissecting Recall of Factual Associations in Auto-Regressive Language Models. arXiv:2304.14767v3, 2023. \url{https://arxiv.org/abs/2304.14767v3}
\bibitem{ref253} Grokking in Linear Estimators -- A Solvable Model that Groks without Understanding. arXiv:2310.16441v1, 2023. \url{https://arxiv.org/abs/2310.16441v1}
\bibitem{ref254} In-Context Learning Dynamics with Random Binary Sequences. arXiv:2310.17639v3, 2023. \url{https://arxiv.org/abs/2310.17639v3}
\bibitem{ref255} Unified View of Grokking, Double Descent and Emergent Abilities A Perspective from Circuits Competition. arXiv:2402.15175v2, 2024. \url{https://arxiv.org/abs/2402.15175v2}
\bibitem{ref256} Towards Understanding Grokking An Effective Theory of Representation Learning. arXiv:2205.10343v2, 2022. \url{https://arxiv.org/abs/2205.10343v2}
\bibitem{ref257} Omnigrok Grokking Beyond Algorithmic Data. arXiv:2210.01117v2, 2022. \url{https://arxiv.org/abs/2210.01117v2}
\bibitem{ref258} Unifying Grokking and Double Descent. arXiv:2303.06173v1, 2023. \url{https://arxiv.org/abs/2303.06173v1}
\bibitem{ref259} Droplets of Good Representations Grokking as a First Order Phase Transition in Two Layer Networks. arXiv:2310.03789v2, 2023. \url{https://arxiv.org/abs/2310.03789v2}
\bibitem{ref260} Progress measures for grokking via mechanistic interpretability. arXiv:2301.05217v3, 2023. \url{https://arxiv.org/abs/2301.05217v3}
\bibitem{ref261} Exploring Memorization in Fine-tuned Language Models. arXiv:2310.06714v2, 2023. \url{https://arxiv.org/abs/2310.06714v2}
\bibitem{ref262} Memorization Without Overfitting Analyzing the Training Dynamics of Large Language Models. arXiv:2205.10770v2, 2022. \url{https://arxiv.org/abs/2205.10770v2}
\bibitem{ref263} From Understanding to Utilization A Survey on Explainability for Large Language Models. arXiv:2401.12874v2, 2024. \url{https://arxiv.org/abs/2401.12874v2}
\bibitem{ref264} Editing Large Language Models Problems, Methods, and Opportunities. arXiv:2305.13172v3, 2023. \url{https://arxiv.org/abs/2305.13172v3}
\bibitem{ref265} Editing Factual Knowledge and Explanatory Ability of Medical Large Language Models. arXiv:2402.18099v1, 2024. \url{https://arxiv.org/abs/2402.18099v1}
\bibitem{ref266} The Cost of Compression Investigating the Impact of Compression on Parametric Knowledge in Language Models. arXiv:2312.00960v1, 2023. \url{https://arxiv.org/abs/2312.00960v1}
\bibitem{ref267} Towards Reliable and Fluent Large Language Models Incorporating Feedback Learning Loops in QA Systems. arXiv:2309.06384v1, 2023. \url{https://arxiv.org/abs/2309.06384v1}
\bibitem{ref268} Second Thoughts are Best Learning to Re-Align With Human Values from Text Edits. arXiv:2301.00355v2, 2023. \url{https://arxiv.org/abs/2301.00355v2}
\bibitem{ref269} Tuning-Free Accountable Intervention for LLM Deployment -- A Metacognitive Approach. arXiv:2403.05636v1, 2024. \url{https://arxiv.org/abs/2403.05636v1}
\bibitem{ref270} Towards Faithful and Robust LLM Specialists for Evidence-Based Question-Answering. arXiv:2402.08277v3, 2024. \url{https://arxiv.org/abs/2402.08277v3}
\bibitem{ref271} The Probabilities Also Matter A More Faithful Metric for Faithfulness of Free-Text Explanations in Large Language Models. arXiv:2404.03189v1, 2024. \url{https://arxiv.org/abs/2404.03189v1}
\bibitem{ref272} Benchmarking Faithfulness Towards Accurate Natural Language Explanations in Vision-Language Tasks. arXiv:2304.08174v1, 2023. \url{https://arxiv.org/abs/2304.08174v1}
\bibitem{ref273} Quantifying Uncertainty in Natural Language Explanations of Large Language Models. arXiv:2311.03533v1, 2023. \url{https://arxiv.org/abs/2311.03533v1}
\bibitem{ref274} On the Robustness of Interpretability Methods. arXiv:1806.08049v1, 2018. \url{https://arxiv.org/abs/1806.08049v1}
\bibitem{ref275} On Measuring Faithfulness or Self-consistency of Natural Language Explanations. arXiv:2311.07466v2, 2023. \url{https://arxiv.org/abs/2311.07466v2}
\bibitem{ref276} Logical Satisfiability of Counterfactuals for Faithful Explanations in NLI. arXiv:2205.12469v1, 2022. \url{https://arxiv.org/abs/2205.12469v1}
\bibitem{ref277} Causal Reasoning and Large Language Models Opening a New Frontier for Causality. arXiv:2305.00050v2, 2023. \url{https://arxiv.org/abs/2305.00050v2}
\bibitem{ref278} QualEval Qualitative Evaluation for Model Improvement. arXiv:2311.02807v1, 2023. \url{https://arxiv.org/abs/2311.02807v1}
\bibitem{ref279} Sensible AI Re-imagining Interpretability and Explainability using Sensemaking Theory. arXiv:2205.05057v1, 2022. \url{https://arxiv.org/abs/2205.05057v1}
\bibitem{ref280} Sequential Explanations with Mental Model-Based Policies. arXiv:2007.09028v1, 2020. \url{https://arxiv.org/abs/2007.09028v1}
\bibitem{ref281} The Definitions of Interpretability and Learning of Interpretable Models. arXiv:2105.14171v1, 2021. \url{https://arxiv.org/abs/2105.14171v1}
\bibitem{ref282} The Role of Individual User Differences in Interpretable and Explainable Machine Learning Systems. arXiv:2009.06675v1, 2020. \url{https://arxiv.org/abs/2009.06675v1}
\bibitem{ref283} Interpreting Blackbox Models via Model Extraction. arXiv:1705.08504v6, 2017. \url{https://arxiv.org/abs/1705.08504v6}
\bibitem{ref284} What is Interpretable Using Machine Learning to Design Interpretable Decision-Support Systems. arXiv:1811.10799v2, 2018. \url{https://arxiv.org/abs/1811.10799v2}
\bibitem{ref285} GAM Changer Editing Generalized Additive Models with Interactive Visualization. arXiv:2112.03245v1, 2021. \url{https://arxiv.org/abs/2112.03245v1}
\bibitem{ref286} Policy Regularization for Legible Behavior. arXiv:2203.04303v1, 2022. \url{https://arxiv.org/abs/2203.04303v1}
\bibitem{ref287} Proto-lm A Prototypical Network-Based Framework for Built-in Interpretability in Large Language Models. arXiv:2311.01732v2, 2023. \url{https://arxiv.org/abs/2311.01732v2}
\bibitem{ref288} Outcome-Explorer A Causality Guided Interactive Visual Interface for Interpretable Algorithmic Decision Making. arXiv:2101.00633v5, 2021. \url{https://arxiv.org/abs/2101.00633v5}
\bibitem{ref289} Minimum Levels of Interpretability for Artificial Moral Agents. arXiv:2307.00660v1, 2023. \url{https://arxiv.org/abs/2307.00660v1}
\bibitem{ref290} Interpretable Artificial Intelligence through the Lens of Feature Interaction. arXiv:2103.03103v1, 2021. \url{https://arxiv.org/abs/2103.03103v1}
\bibitem{ref291} A Concept and Argumentation based Interpretable Model in High Risk Domains. arXiv:2208.08149v1, 2022. \url{https://arxiv.org/abs/2208.08149v1}
\bibitem{ref292} Debating with More Persuasive LLMs Leads to More Truthful Answers. arXiv:2402.06782v2, 2024. \url{https://arxiv.org/abs/2402.06782v2}
\bibitem{ref293} A Survey on Explainable Artificial Intelligence (XAI) Towards Medical XAI. arXiv:1907.07374v5, 2019. \url{https://arxiv.org/abs/1907.07374v5}
\bibitem{ref294} The Calibration Gap between Model and Human Confidence in Large Language Models. arXiv:2401.13835v1, 2024. \url{https://arxiv.org/abs/2401.13835v1}
\bibitem{ref295} Causal Rule Learning Enhancing the Understanding of Heterogeneous Treatment Effect via Weighted Causal Rules. arXiv:2310.06746v1, 2023. \url{https://arxiv.org/abs/2310.06746v1}
\bibitem{ref296} Meaningful Models Utilizing Conceptual Structure to Improve Machine Learning Interpretability. arXiv:1607.00279v1, 2016. \url{https://arxiv.org/abs/1607.00279v1}
\bibitem{ref297} Large Language Models and Causal Inference in Collaboration A Comprehensive Survey. arXiv:2403.09606v1, 2024. \url{https://arxiv.org/abs/2403.09606v1}
\bibitem{ref298} Is Knowledge All Large Language Models Needed for Causal Reasoning. arXiv:2401.00139v1, 2023. \url{https://arxiv.org/abs/2401.00139v1}
\bibitem{ref299} Towards Interpretable Mental Health Analysis with Large Language Models. arXiv:2304.03347v4, 2023. \url{https://arxiv.org/abs/2304.03347v4}
\bibitem{ref300} DDCoT Duty-Distinct Chain-of-Thought Prompting for Multimodal Reasoning in Language Models. arXiv:2310.16436v2, 2023. \url{https://arxiv.org/abs/2310.16436v2}
\bibitem{ref301} Steering LLMs Towards Unbiased Responses A Causality-Guided Debiasing Framework. arXiv:2403.08743v1, 2024. \url{https://arxiv.org/abs/2403.08743v1}
\bibitem{ref302} Causality Learning A New Perspective for Interpretable Machine Learning. arXiv:2006.16789v2, 2020. \url{https://arxiv.org/abs/2006.16789v2}
\bibitem{ref303} Interpretable Machine Learning Fundamental Principles and 10 Grand Challenges. arXiv:2103.11251v2, 2021. \url{https://arxiv.org/abs/2103.11251v2}
\bibitem{ref304} Towards Reasoning in Large Language Models via Multi-Agent Peer Review Collaboration. arXiv:2311.08152v2, 2023. \url{https://arxiv.org/abs/2311.08152v2}
\bibitem{ref305} OpenEval Benchmarking Chinese LLMs across Capability, Alignment and Safety. arXiv:2403.12316v1, 2024. \url{https://arxiv.org/abs/2403.12316v1}
\bibitem{ref306} Prometheus Inducing Fine-grained Evaluation Capability in Language Models. arXiv:2310.08491v2, 2023. \url{https://arxiv.org/abs/2310.08491v2}
\bibitem{ref307} AgentSims An Open-Source Sandbox for Large Language Model Evaluation. arXiv:2308.04026v1, 2023. \url{https://arxiv.org/abs/2308.04026v1}
\bibitem{ref308} Chatbot Arena An Open Platform for Evaluating LLMs by Human Preference. arXiv:2403.04132v1, 2024. \url{https://arxiv.org/abs/2403.04132v1}
\bibitem{ref309} NBIAS A Natural Language Processing Framework for Bias Identification in Text. arXiv:2308.01681v3, 2023. \url{https://arxiv.org/abs/2308.01681v3}
\bibitem{ref310} Eagle Ethical Dataset Given from Real Interactions. arXiv:2402.14258v1, 2024. \url{https://arxiv.org/abs/2402.14258v1}
\bibitem{ref311} FairMonitor A Four-Stage Automatic Framework for Detecting Stereotypes and Biases in Large Language Models. arXiv:2308.10397v2, 2023. \url{https://arxiv.org/abs/2308.10397v2}
\bibitem{ref312} Bias in Language Models Beyond Trick Tests and Toward RUTEd Evaluation. arXiv:2402.12649v1, 2024. \url{https://arxiv.org/abs/2402.12649v1}
\bibitem{ref313} PRD Peer Rank and Discussion Improve Large Language Model based Evaluations. arXiv:2307.02762v1, 2023. \url{https://arxiv.org/abs/2307.02762v1}
\bibitem{ref314} Learning Fair Ranking Policies via Differentiable Optimization of Ordered Weighted Averages. arXiv:2402.05252v1, 2024. \url{https://arxiv.org/abs/2402.05252v1}
\bibitem{ref315} H2O Open Ecosystem for State-of-the-art Large Language Models. arXiv:2310.13012v2, 2023. \url{https://arxiv.org/abs/2310.13012v2}
\bibitem{ref316} MedAgents Large Language Models as Collaborators for Zero-shot Medical Reasoning. arXiv:2311.10537v3, 2023. \url{https://arxiv.org/abs/2311.10537v3}
\bibitem{ref317} Clinical Camel An Open Expert-Level Medical Language Model with Dialogue-Based Knowledge Encoding. arXiv:2305.12031v2, 2023. \url{https://arxiv.org/abs/2305.12031v2}
\bibitem{ref318} PMC-LLaMA Towards Building Open-source Language Models for Medicine. arXiv:2304.14454v3, 2023. \url{https://arxiv.org/abs/2304.14454v3}
\bibitem{ref319} MEDITRON-70B Scaling Medical Pretraining for Large Language Models. arXiv:2311.16079v1, 2023. \url{https://arxiv.org/abs/2311.16079v1}
\bibitem{ref320} CLUE A Clinical Language Understanding Evaluation for LLMs. arXiv:2404.04067v2, 2024. \url{https://arxiv.org/abs/2404.04067v2}
\bibitem{ref321} Friend or Foe Exploring the Implications of Large Language Models on the Science System. arXiv:2306.09928v1, 2023. \url{https://arxiv.org/abs/2306.09928v1}
\bibitem{ref322} Automatic Interactive Evaluation for Large Language Models with State Aware Patient Simulator. arXiv:2403.08495v2, 2024. \url{https://arxiv.org/abs/2403.08495v2}
\bibitem{ref323} An Automatic Evaluation Framework for Multi-turn Medical Consultations Capabilities of Large Language Models. arXiv:2309.02077v1, 2023. \url{https://arxiv.org/abs/2309.02077v1}
\bibitem{ref324} LV-Eval A Balanced Long-Context Benchmark with 5 Length Levels Up to 256K. arXiv:2402.05136v1, 2024. \url{https://arxiv.org/abs/2402.05136v1}
\bibitem{ref325} Biomedical knowledge graph-enhanced prompt generation for large language models. arXiv:2311.17330v1, 2023. \url{https://arxiv.org/abs/2311.17330v1}
\bibitem{ref326} Augmented Large Language Models with Parametric Knowledge Guiding. arXiv:2305.04757v2, 2023. \url{https://arxiv.org/abs/2305.04757v2}
\bibitem{ref327} The Magic of IF Investigating Causal Reasoning Abilities in Large Language Models of Code. arXiv:2305.19213v1, 2023. \url{https://arxiv.org/abs/2305.19213v1}
\bibitem{ref328} Can We Utilize Pre-trained Language Models within Causal Discovery Algorithms. arXiv:2311.11212v1, 2023. \url{https://arxiv.org/abs/2311.11212v1}
\bibitem{ref329} Cause and Effect Can Large Language Models Truly Understand Causality. arXiv:2402.18139v2, 2024. \url{https://arxiv.org/abs/2402.18139v2}
\bibitem{ref330} Evaluating Interventional Reasoning Capabilities of Large Language Models. arXiv:2404.05545v1, 2024. \url{https://arxiv.org/abs/2404.05545v1}
\bibitem{ref331} Estimating Categorical Counterfactuals via Deep Twin Networks. arXiv:2109.01904v6, 2021. \url{https://arxiv.org/abs/2109.01904v6}
\bibitem{ref332} CLadder Assessing Causal Reasoning in Language Models. arXiv:2312.04350v3, 2023. \url{https://arxiv.org/abs/2312.04350v3}
\bibitem{ref333} From Query Tools to Causal Architects Harnessing Large Language Models for Advanced Causal Discovery from Data. arXiv:2306.16902v1, 2023. \url{https://arxiv.org/abs/2306.16902v1}
\bibitem{ref334} Bridging Causal Discovery and Large Language Models A Comprehensive Survey of Integrative Approaches and Future Directions. arXiv:2402.11068v1, 2024. \url{https://arxiv.org/abs/2402.11068v1}
\bibitem{ref335} Causal Proxy Models for Concept-Based Model Explanations. arXiv:2209.14279v1, 2022. \url{https://arxiv.org/abs/2209.14279v1}
\bibitem{ref336} Causal Deep Learning. arXiv:2303.02186v2, 2023. \url{https://arxiv.org/abs/2303.02186v2}
\bibitem{ref337} Can Large Language Models Infer Causation from Correlation. arXiv:2306.05836v3, 2023. \url{https://arxiv.org/abs/2306.05836v3}
\bibitem{ref338} Do Large Language Models Understand Logic or Just Mimick Context. arXiv:2402.12091v1, 2024. \url{https://arxiv.org/abs/2402.12091v1}
\bibitem{ref339} Causal Inference Using LLM-Guided Discovery. arXiv:2310.15117v1, 2023. \url{https://arxiv.org/abs/2310.15117v1}
\bibitem{ref340} Discovery of the Hidden World with Large Language Models. arXiv:2402.03941v1, 2024. \url{https://arxiv.org/abs/2402.03941v1}
\bibitem{ref341} Navigating Complexity Orchestrated Problem Solving with Multi-Agent LLMs. arXiv:2402.16713v1, 2024. \url{https://arxiv.org/abs/2402.16713v1}
\bibitem{ref342} Reasoning Capacity in Multi-Agent Systems Limitations, Challenges and Human-Centered Solutions. arXiv:2402.01108v1, 2024. \url{https://arxiv.org/abs/2402.01108v1}
\bibitem{ref343} Puzzle Solving using Reasoning of Large Language Models A Survey. arXiv:2402.11291v2, 2024. \url{https://arxiv.org/abs/2402.11291v2}
\bibitem{ref344} Neural Theory-of-Mind On the Limits of Social Intelligence in Large LMs. arXiv:2210.13312v2, 2022. \url{https://arxiv.org/abs/2210.13312v2}
\bibitem{ref345} Views Are My Own, But Also Yours Benchmarking Theory of Mind using Common Ground. arXiv:2403.02451v1, 2024. \url{https://arxiv.org/abs/2403.02451v1}
\bibitem{ref346} ToMChallenges A Principle-Guided Dataset and Diverse Evaluation Tasks for Exploring Theory of Mind. arXiv:2305.15068v2, 2023. \url{https://arxiv.org/abs/2305.15068v2}
\bibitem{ref347} MindGames Targeting Theory of Mind in Large Language Models with Dynamic Epistemic Modal Logic. arXiv:2305.03353v2, 2023. \url{https://arxiv.org/abs/2305.03353v2}
\bibitem{ref348} NegotiationToM A Benchmark for Stress-testing Machine Theory of Mind on Negotiation Surrounding. arXiv:2404.13627v1, 2024. \url{https://arxiv.org/abs/2404.13627v1}
\bibitem{ref349} ConTextual Evaluating Context-Sensitive Text-Rich Visual Reasoning in Large Multimodal Models. arXiv:2401.13311v1, 2024. \url{https://arxiv.org/abs/2401.13311v1}
\bibitem{ref350} Large Language Model for Causal Decision Making. arXiv:2312.17122v3, 2023. \url{https://arxiv.org/abs/2312.17122v3}
\bibitem{ref351} CausalBench A Comprehensive Benchmark for Causal Learning Capability of Large Language Models. arXiv:2404.06349v1, 2024. \url{https://arxiv.org/abs/2404.06349v1}
\bibitem{ref352} Understanding Causality with Large Language Models Feasibility and Opportunities. arXiv:2304.05524v1, 2023. \url{https://arxiv.org/abs/2304.05524v1}
\bibitem{ref353} Quantifying and Mitigating Unimodal Biases in Multimodal Large Language Models A Causal Perspective. arXiv:2403.18346v3, 2024. \url{https://arxiv.org/abs/2403.18346v3}
\bibitem{ref354} Causal-CoG A Causal-Effect Look at Context Generation for Boosting Multi-modal Language Models. arXiv:2312.06685v1, 2023. \url{https://arxiv.org/abs/2312.06685v1}
\bibitem{ref355} Resolving Ethics Trade-offs in Implementing Responsible AI. arXiv:2401.08103v3, 2024. \url{https://arxiv.org/abs/2401.08103v3}
\bibitem{ref356} Designing for Human Rights in AI. arXiv:2005.04949v2, 2020. \url{https://arxiv.org/abs/2005.04949v2}
\bibitem{ref357} Socially Cognizant Robotics for a Technology Enhanced Society. arXiv:2310.18303v1, 2023. \url{https://arxiv.org/abs/2310.18303v1}
\bibitem{ref358} Evaluating the Social Impact of Generative AI Systems in Systems and Society. arXiv:2306.05949v2, 2023. \url{https://arxiv.org/abs/2306.05949v2}
\bibitem{ref359} Inclusive Artificial Intelligence. arXiv:2212.12633v2, 2022. \url{https://arxiv.org/abs/2212.12633v2}
\bibitem{ref360} Responsible and Inclusive Technology Framework A Formative Framework to Promote Societal Considerations in Information Technology Contexts. arXiv:2302.11565v1, 2023. \url{https://arxiv.org/abs/2302.11565v1}
\bibitem{ref361} Disciplining deliberation a sociotechnical perspective on machine learning trade-offs. arXiv:2403.04226v1, 2024. \url{https://arxiv.org/abs/2403.04226v1}
\bibitem{ref362} Leveraging Multi-Method Evaluation for Multi-Stakeholder Settings. arXiv:2001.04348v1, 2019. \url{https://arxiv.org/abs/2001.04348v1}
\bibitem{ref363} Evaluatology The Science and Engineering of Evaluation. arXiv:2404.00021v1, 2024. \url{https://arxiv.org/abs/2404.00021v1}
\bibitem{ref364} Axes for Sociotechnical Inquiry in AI Research. arXiv:2105.06551v1, 2021. \url{https://arxiv.org/abs/2105.06551v1}
\bibitem{ref365} A Modest Proposal (Possible) Implications for Appropriation. arXiv:1306.5247v1, 2013. \url{https://arxiv.org/abs/1306.5247v1}
\bibitem{ref366} Buying time in software development how estimates become commitments. arXiv:2105.08018v1, 2021. \url{https://arxiv.org/abs/2105.08018v1}
\bibitem{ref367} Policy Scan and Technology Strategy Design methodology. arXiv:1806.03235v1, 2018. \url{https://arxiv.org/abs/1806.03235v1}
\bibitem{ref368} Beyond Interpretable Benchmarks Contextual Learning through Cognitive and Multimodal Perception. arXiv:2304.00002v2, 2022. \url{https://arxiv.org/abs/2304.00002v2}
\end{thebibliography}

\end{document}
